% This file was converted to LaTeX by Writer2LaTeX ver. 1.2
% see http://writer2latex.sourceforge.net for more info
\documentclass{paper}
\usepackage{ppc_utf8}

\usepackage[utf8]{inputenc}
\usepackage[T3,T1]{fontenc}
%\usepackage[portuges]{babel}
\usepackage[noenc]{tipa}
\usepackage{tipx}
\usepackage[geometry,weather,misc,clock]{ifsym}
\usepackage{pifont}
\usepackage{eurosym}
\usepackage{amsmath}
\usepackage{wasysym}
\usepackage{amssymb,amsfonts,textcomp}
\usepackage{array}
\usepackage{supertabular}
\usepackage{hhline}
%\usepackage[pdftex]{graphicx}
\usepackage{setspace}
\usepackage{xspace}
%\usepackage[pdftex, colorlinks]{hyperref}
\usepackage[space]{grffile} % permite espaços no nome de arquivo
\usepackage{todonotes}
\usepackage{pdfpages}% permite inserir páginas PDF
\usepackage{tocloft}
\usepackage{afterpage}
\usepackage[alf,abnt-repeated-author-omit=yes,abnt-etal-list=0]{abntex2cite}

\usepackage{caption}
\captionsetup{width=1.2\linewidth,textfont={bf}}

%%%%%%%%
\makeatletter
\newcommand\arraybslash{\let\\\@arraycr}
\makeatother
\setlength\tabcolsep{1mm}
\renewcommand\arraystretch{1.3}
\def\web#1{\href{#1}{#1}}

%%%%%%%%%%%%%%%%%%%%%%%%%%%%%%%%%%%%%%%%%%%%%%%%%%%%%%%%%%%%%%%%%%%%%%%%%%%%%%%%%%%%%
% Este arquivo é a base comum do PPC da ECA e da EME
% O resultado da compilação depende do arquivo PPC-flag.tex
% \def\COMPILAR{0}  :  valor=0: compila PPC ECA; valor<>0:compila PPC EME
% A compilação automática é feita com o script em Python: Compila_PPCv0.0.py
%%%%%%%%%%%%%%%%%%%%%%%%%%%%%%%%%%%%%%%%%%%%%%%%%%%%%%%%%%%%%%%%%%%%%%%%%%%%%%%%%%%%%
%
\input{PPC-flag}
\if\COMPILAR0
    \def\COORDENADOREIXO{FRANCISCO NÉLIO COSTA FREITAS\xspace}
\def\COORDENADORCURSO{Luiz Daniel Santos Bezerra\xspace}
\def\NOMEENG{Engenharia de Controle e Automação}
\def\NOMECURSO{Curso de Bacharelado em \NOMEENG\xspace}
\def\NOMEPROFISSIONAL{Bacharel em Engenharia de Controle e Automação\xspace}
\def\NUMVAGAS{30\xspace}
\def\CREDITO{20\xspace}

\def\CHBASICA{1280\xspace}
\def\CHOPT{920\xspace}
\def\CHOPTMIN{160\xspace}
\def\CHTOTALOB{4000\xspace}     % DISCIPLINAS OBRIGATÓRIAS + ESTÁGIO
\def\CHTOTALDOB{3440\xspace}    % DISCIPLINAS OBRIGATÓRIAS
\def\CHESTAGIO{400\xspace}
\def\CHPROF{2560\xspace}
\def\CHCOMP{400\xspace}
\def\CHTCC{40\xspace}
\def\CHTOTALGERAL{5160\xspace}    % DISCIPLINAS OBRIGATÓRIAS + ESTÁGIO + OPTATIVAS
\def\CHINT{0\xspace}            % DISCIPLINAS INTEGRADORAS

\def\CHBASICACR{64\xspace}
\def\CHOPTCR{46\xspace}
\def\CHOPTMINCR{8\xspace}
\def\CHTOTALOBCR{200\xspace}
\def\CHTOTALDOBCR{172\xspace}
\def\CHESTAGIOCR{20\xspace}
\def\CHPROFCR{128\xspace}
\def\CHCOMPCR{20\xspace}
\def\CHTCCCR{2\xspace}
\def\CHTOTALGERALCR{258\xspace}
\def\CHINTCR{0\xspace}

\def\CHBASICASE{6,4\xspace}
\def\CHOPTSE{4\xspace}
\def\CHOPTMINSE{4\xspace}
\def\CHTOTALOBSE{41,2\xspace}
\def\CHTOTALDOBSE{17,2\xspace}     % CHTOTALDOB/10/20
\def\CHESTAGIOSE{20\xspace}
\def\CHPROFSE{12,8\xspace}         % CHPROF/10/20
\def\CHCOMPSE{2\xspace}
\def\CHTCCSE{2\xspace}
\def\CHTOTALGERALSE{43,2\xspace}   % CHTOTALDOBSE + CHESTAGIOSE + CHOPTSE
\def\CHINTSE{0\xspace}

\def\CHBASICAPE{24,8\xspace}       % PERCENTUAIS EM RELAÇÃO À CHTOTALGERAL
\def\CHOPTPE{17,8\xspace}
\def\CHOPTMINPE{3,1\xspace}
\def\CHTOTALOBPE{77,5\xspace}
\def\CHTOTALDOBPE{66,7\xspace}
\def\CHESTAGIOPE{7,8\xspace}
\def\CHPROFPE{49,6\xspace}
\def\CHCOMPPE{7,8\xspace}
\def\CHTCCPE{0,8\xspace}
\def\CHTOTALGERALPE{100,0\xspace}
\def\CHINTPE{0\xspace}

\def\CHBASICAPEMIN{32\xspace}   % PERCENTUAIS EM RELAÇÃO À CHTOTALOB
\def\CHOPTPEMIN{23\xspace}
\def\CHOPTMINPEMIN{4,0\xspace}
\def\CHTOTALOBPEMIN{100,0\xspace}
\def\CHTOTALDOBPEMIN{86\xspace}
\def\CHESTAGIOPEMIN{10\xspace}
\def\CHPROFPEMIN{64\xspace}
\def\CHCOMPPEMIN{10\xspace}
\def\CHTCCPEMIN{1,0\xspace}
\def\CHTOTALGERALPEMIN{129\xspace}
\def\CHINTPEMIN{0\xspace}


\def\NOMEEGRESSO{Bacharel em Engenharia de Controle e Automação\xspace}
\def\COMISSAO{
Celso Rogério Schmidlin Júnior

Fabrício Bandeira da Silva

Geraldo Luis Bezerra Ramalho

José Daniel de Alencar Santos

Luiz Daniel Santos Bezerra

Samuel Vieira Dias
}
\def\NDE{
\begin{longtable}{|c|c|c|}
\hline 
\textbf{NOME} & \textbf{FUNÇÃO} & \textbf{SIAPE} \\ \hline

José Daniel de Alencar Santos & Presidente & 1442729 \\ \hline

Luiz Daniel Santos Bezerra & Membro Docente & 1842966 \\ \hline

Adriano Holanda Pereira & Membro Docente & 1556624 \\ \hline

Antônio Barbosa de Souza Júnior & Membro Docente & 2031223 \\ \hline

Daniel Barbosa de Brito & Membro Docente & 1960021 \\ \hline

Fábio Timbó Brito & Membro Docente & 1641803 \\ \hline
 
Jonatha Rodrigues da Costa & Membro Docente & 1075472 \\ \hline

Rafael Oliveira de Sousa & Membro Docente & 2188862 \\ \hline

Venício Soares de Oliveira & Membro Docente & 1659388 \\ \hline

\end{longtable}
}
\def\COLEGIADO{
\begin{longtable}{|c|c|c|}
\hline 
\textbf{NOME} & \textbf{FUNÇÃO} & \textbf{MATRÍCULA} \\ \hline

Luiz Daniel Santos Bezerra & Presidente & 1842966 \\ \hline

Leilane Lima Almeida Evangelista & Representante da CTP & 2165051  \\ \hline

Roseane Michelle de Lima Silveira & Suplente Representante da CTP & 1576780 \\ \hline

José Daniel de Alencar Santos & Docente da Área de Estudos Específicos & 1442729 \\ \hline

Rafael Oliveira de Sousa & Suplente de Docente da Área de Estudos Específicos & 2188862 \\ \hline

Adriano Holanda Pereira & Docente da Área de Estudos Específicos & 1556624 \\ \hline

Venceslau Xavier de Lima Filho & Suplente de Docente da Área de Estudos Específicos & 1544405 \\ \hline

Daniel Barbosa de Brito & Docente da Área de Estudos Específicos & 1960021 \\ \hline

Jonatha Rodrigues da Costa & Suplente de Docente da Área de Estudos Específicos & 1075472 \\ \hline

Francisca Ione Chaves & Docente da Área de Estudos Básicos & 1460033 \\ \hline

Francisco Ricardo Nogueira de Vasconcelos & Suplente de Docente da Área de Estudos Básicos & 1616139 \\ \hline

Ismael Fama Ribeiro & Discente & 20152045070200 \\ \hline

Maria Eliziane de Oliveira Tabosa & Discente Suplente & 20152045070014 \\ \hline

Laisla Firmino Teixeira Costa & Discente & 20152045070332 \\ \hline

Leonardo Adriano Vasconcelos de Oliveira & Discente Suplente & 20162045070208 \\ \hline
\end{longtable}
}

\def\COLABORADORES{
Adriano Holanda Pereira

Fábio Timbó Brito

Francisco Frederico dos Santos Matos

Francisco José dos Santos Oliveira

Francisco Nélio Costa Freitas

Pedro Pedrosa Rebouças Filho

Rodrigo Freitas Guimarães

Venceslau Xavier de Lima Filho
}
\def\PEDAGOGO{
Isabel Magda Said Pierre Carneiro
}

\def\PEDAGOGOS{
Samoel Rodrigues da Silva

Leilane Lima Almeida Evangelista
}

\def\RESOLUCAOAPROVACAODEZENOVE{n° 26/2019-CONSUP/IFCE de 24 de maio de 2019\xspace}
\def\RESOLUCAOAPROVACAONOVA{n° 117/2017-CONSUP/IFCE de 27 de novembro de 2017\xspace}
\def\RESOLUCAOAPROVACAO{n° 002/2014-CONSUP/IFCE de 31 de janeiro de 2014\xspace}
\def\PORTARIACOMISSAO{n° 023/GDG de 27 de março de 2013\xspace}
\def\PORTARIACOMISSAOALT{n° 042/GDG de 22 de abril de 2013\xspace}
\def\PORTARIARECONHECIMENTO{n° 430/Ministério da Educação de 03 de maio de 2021\xspace}
\def\PORTARIACOLEGIADO{n° 099/GDG de 23 de dezembro de 2010\xspace}
\def\PORTARIACOLEGIADOALT{a portaria n° 086/GDG de 26 outubro de 2011\xspace}
\def\PORTARIANDE{n° 098/GDG  de 23 de dezembro de 2010\xspace}
\def\PORTARIANDEALT{a portaria n° 083/GDG de 13 de outubro de 2011\xspace}

\def\PORTARIANDEATUAL{Nº 132/GAB-MAR/DG-MAR/MARACANAÚ, DE 23 DE NOVEMBRO DE 2020\xspace}
\def\PORTARIACOLEGIADOATUAL{Nº 69/GAB-MAR/DG-MAR/MARACANAÚ, DE 08 DE JULHO DE 2021\xspace}

\def\LMATdisc{Materiais}
\def\LIAFdisc{Materiais e Resistência dos materiais}
\def\LMETdisc{Metrologia}
\def\LTFdisc{Fisica II, Instrumentação}
\def\LIAdisc{Lógica de Programação, Linguagem de Programação, Circuitos Elétricos I e Manufatura Integrada por Computador, Métodos Numéricos, Processamento Digital de Imagens}
\def\LSHIPdisc{Acionamentos Hidráulicos e Pneumáticos, Controle I, Controle II, Identificação de Sistemas}
\def\LEEdisc{Circuitos Elétricos I, Eletrônica I (Dig.+Analógica Básica), Eletrônica II (Dig.+Analógica Avançada), Dispositivos Periféricos}
\def\LAMEPdisc{Circuitos Elétricos II, Eletrônica III (Industrial), Instalações Elétricas, Máquinas Elétricas e Acionamentos de Máquinas}
\def\LINCdisc{Instrumentação, Processamento Digital de Sinais (PDS), Controladores Lógicos Programáveis (CLP), Controle II (Discreto), Sistemas Embarcados, Sistemas Digitais de Controle Distribuído (SDCD)}
\def\LPROTdisc{Microcontroladores, Dispositivos Periféricos, Robótica I, Robótica II e Sistemas Embarcados}
\def\LPCdisc{Eletrônica I, Eletrônica II, Eletrônica III, Circuitos Elétricos I}
\def\LAMSCdisc{Identificação de Sistemas, Inteligência Computacional Aplicada, Microcontroladores, Robótica}

\def\SEMESTRE{2014.2}
\def\PAGINA{https://ifce.edu.br/maracanau/menu/cursos/superiores/bacharelados/controle\_automacao}
\else
    \input{PPC-EMEConfig}
\fi
\def\PUDcurso{\NOMEENG}


%*** GRAPHICS RELATED PACKAGES ***
\graphicspath{{figuras/}} 
\DeclareGraphicsExtensions{.eps,.png,.jpg,.bmp}

\pdfinfo{
   /Author (Eixo da Indústria, IFCE Campus Maracanaú)
   /Title (Projeto Pedagógico do Curso de Engenharia de Controle e Automação)
   /Creator (pdfLaTeX)
   /CreationDate (D:20130701100000)
   /Subject (versão jul/2013)
   /Keywords (PUD)
   /ModDate (D:\pdfdate)
}

%DEFINIR MARGENS DO DOCUMENTO
\geometry{a4paper,left=3cm,right=2cm,top=3cm,bottom=2cm} %INSERIDO POR CELSO SCHMIDLIN
 
\begin{document}
\thispagestyle{empty} 
\clearpage\setcounter{page}{1}
\onehalfspacing
\begin{center}
\begin{figure}[ht]
\centering

\includegraphics[width=2.845cm,height=2.845cm]{brasaoBR}
\end{figure}

\textbf{MINISTÉRIO DA EDUCAÇÃO}

\textbf{SECRETARIA DE EDUCAÇÃO PROFISSIONAL E TECNOLÓGICA}

\textbf{INSTITUTO FEDERAL DE EDUCAÇÃO, CIÊNCIA E TECNOLOGIA DO CEARÁ}

\textbf{\textit{CAMPUS} MARACANAÚ}

\vspace{1cm}

%\begin{figure}[ht]
%\centering
%
\includegraphics[width=4.916cm,height=1.928cm]{logo_ifce}
%\end{figure}

%PRÓ-REITORIA DE ENSINO

%DIRETORIA DE ENSINO

\vspace{4cm}

\textbf{\MakeUppercase{\large{PROJETO PEDAGÓGICO}}}

\textbf{\MakeUppercase{\large\NOMECURSO}}

\vspace{5cm}

Aprovado pela Resolução \RESOLUCAOAPROVACAO .

Atualizado pela Resolução \RESOLUCAOAPROVACAONOVA .

Atualizado pela Resolução \RESOLUCAOAPROVACAODEZENOVE .

\vspace{5cm}

\textbf{Maracanaú}

\textbf{2013}
\end{center}

% CONTRA-CAPA ------------------------------------------------------------

\input{PPC-ECAEME-ContraCapa}

% DADOS DA INSTITUIÇÃO ------------------------------------------------------

\clearpage
\thispagestyle{empty}
\begin{center}
\textbf{DADOS DA INSTITUIÇÃO}
\end{center}

\vspace{1cm}

\begin{itemize}
\item \textbf{Nome:} Instituto Federal de Educação, Ciência e Tecnologia do Ceará – IFCE / Campus Maracanaú

\item\textbf{CNPJ:} 10.744.098/0009-00

\item \textbf{UF:} {Ceará} ~~~~~~ \textbf{Município: }Maracanaú

\item \textbf{Endereço:} Av. Parque Central, S/N

\item\textbf{Bairro:} Distrito Industrial I

\item\textbf{CEP:} 61939-140

\item\textbf{Telefone(s):} (85) 3878 6300

\item\textbf{Fax:} (85) 3878 6311

\item\textbf{Comprovante do Imóvel: } \textbf{ANEXO \ref{anexoImovel}}.

%\item\textbf{Diretor Geral do Campus: }Júlio César da Costa Silva

\item\textbf{E-mail da Diretção Geral:} gabmaracanau@ifce.edu.br

\item\textbf{Página Institucional na Internet:} \web{http://ifce.edu.br/maracanau}

\end{itemize}

% DADOS DO CURSO ------------------------------------------------------

\clearpage
\thispagestyle{empty}
\begin{center}
\textbf{DADOS DO CURSO}
\end{center}

\vspace{1cm}

\begin{itemize}
\item \textbf{Denominação:} \MakeUppercase{Bacharelado em \NOMEENG}
\item \textbf{Titulação conferida:} \MakeUppercase{\NOMEEGRESSO}
\item \textbf{Nível:} \MakeUppercase{SUPERIOR}
\item \textbf{Modalidade:} \MakeUppercase{PRESENCIAL}
\item \textbf{Período de Integralização:} \MakeUppercase{10 semestres}
\item \textbf{Periodicidade: }{semestral (  )  anual ( X ) }
\item \textbf{Formas de ingresso: }{SISU ( X )  Vestibular (  ) Transferência ( X ) Diplomado ( X ) }
\item \textbf{Número de vagas ofertadas/ano: }{\NUMVAGAS}
\item \textbf{Turno de Funcionamento: }{integral ( X ) noturno ( )}
\item \textbf{Semestre de Início:} \MakeUppercase{\SEMESTRE}
\item\textbf{Página Institucional na Internet:} 

\web{\PAGINA}
\item\textbf{Telefone(s):} (85) 3878 6316 / 3878 6345
\item \textbf{Carga Horária dos Componentes Curriculares Obrigatórios: }{\CHTOTALDOB}
\item \textbf{Carga Horária Total de Disciplinas Optativas Previstas: }{\CHOPT}
\item \textbf{Carga Horária Obrigatória de Disciplinas Optativas: }{\CHOPTMIN}
\item \textbf{Carga Horária do Estágio Obrigatório: }{\CHESTAGIO}
\item \textbf{Carga Horária das Atividades Complementares Não Obrigatórias: }{\CHCOMP}
\item \textbf{Carga Horária Obrigatória para Integralização do Curso: }{\CHTOTALOB}
\item \textbf{Carga Horária do Trabalho de Conclusão do Curso: }{\CHTCC}
%\item \textbf{Carga Horária Total Geral: }{\CHTOTALGERAL}
\item \textbf{Sistema de Carga Horária: }{1 crédito = 20 horas}
%\item \textbf{Nome do Coordenador do Curso: }{\COORDENADORCURSO}
\end{itemize}

% REVISÕES ------------------------------------------------------

\clearpage
\thispagestyle{empty}
\begin{center}
\textbf{REVISÕES}
\end{center}
\bigskip
\begin{center}
\begin{tabular}{lp{2cm}p{11cm}}
Revisão & Data & Descrição\\
\hline
\if\COMPILAR0
    \input{PPC-ECARevisoes}
\else
    \input{PPC-EMERevisoes}
\fi
\hline
\end{tabular}

\clearpage
\thispagestyle{empty}
\begin{tabular}{lp{2cm}p{11cm}}
Revisão & Data & Descrição\\
\hline
\if\COMPILAR0
    \input{PPC-ECARevisoes2}
\else
    \input{PPC-EMERevisoes2}
\fi
\hline
\end{tabular}

\clearpage
\thispagestyle{empty}
\begin{tabular}{lp{2cm}p{11cm}}
Revisão & Data & Descrição\\
\hline
\if\COMPILAR0
    \input{PPC-ECARevisoes3}
\else
    \input{PPC-EMERevisoes3}
\fi
\hline
\end{tabular}

\footnotesize{Compilação mais recente: \today}
\end{center}
\bigskip
%%%%%%

% SUMÁRIO -----------------------------------------------------------------

\clearpage
\thispagestyle{empty}
\begin{center}
\textbf{SUMÁRIO}
\afterpage{\thispagestyle{empty}}
\end{center}

\setcounter{tocdepth}{2} %%%%
\setcounter{secnumdepth}{2}
\renewcommand\contentsname{}
\tableofcontents
\afterpage{\thispagestyle{empty}}

% APRESENTAÇÃO ------------------------------------------------------------

\clearpage
\thispagestyle{empty}
\begin{center}
\section{APRESENTAÇÃO}
%\textbf{APRESENTAÇÃO}
\end{center}

\bigskip

%/O presente documento visa apresentar e detalhar a proposta pedagógica do \NOMECURSO do Instituto Federal de Educação, Ciência e Tecnologia do Ceará – IFCE, Campus Maracanaú, segundo diretrizes curriculares em vigor da Lei de Diretrizes e Bases da Educação Nacional – LDB No 9.394/96, e de acordo com as diretrizes que regulamentam os cursos superiores de graduação tecnológica. Elaborado pela Comissão nomeada pela portaria \PORTARIACOMISSAO em conjunto com os membros do NDE e acompanhado pelos membros do Colegiado.
 
O presente documento visa apresentar e detalhar a proposta pedagógica do \NOMECURSO do Instituto Federal de Educação, Ciência e Tecnologia do Ceará - IFCE, Campus Maracanaú, ofertado na modalidade presencial. O projeto está fundamentado na Lei de Diretrizes e bases da Educação Nacional (LDB), Lei n° 9.394/96, bem como nas normativas legais em âmbito nacional e institucional que regulamentam os cursos superiores de graduação, especialmente nas normativas específicas para os cursos de bacharelado. 

O projeto foi elaborado pela comissão nomeada pela portaria \PORTARIACOMISSAO e aprovado pela Resolução \RESOLUCAOAPROVACAO. Foi atualizado por meio da Resolução \RESOLUCAOAPROVACAONOVA, fruto da atuação conjunta e contínua do Núcleo Docente Estruturante (NDE) e Colegiado do curso, cientes de que o Projeto Pedagógico do Curso (PPC) não é um documento imutável, mas sim discutido e reavaliado pela comunidade acadêmica, visando a contínua promoção da qualidade do curso. 

A formatação do referido projeto apresenta os objetivos, a organização curricular, os procedimentos metodológicos e de avaliação do processo de ensino e aprendizagem e do curso, entre outros aspectos também relevantes, visando à formação não somente de um \NOMEPROFISSIONAL, mas de um cidadão capaz de atuar no seu contexto social com competência técnica e humanamente comprometido com a construção de uma sociedade mais justa, solidária e ética, em consonância com a missão do IFCE e com os objetivos dos Institutos Federais, nos termos da Lei n° 11.892/2008.

% CONTEXTUALIZAÇÃO DA INSTITUIÇÃO --------------------------------------------

\clearpage\section{CONTEXTUALIZAÇÃO DA INSTITUIÇÃO}
Nos primeiros vinte anos após a Proclamação da República, as indústrias brasileiras já apresentavam algum crescimento, demandando a necessidade de mão de obra mais qualificada. As novas tarefas exigiam pessoas com conhecimentos especializados e apontavam para a necessidade de se estabelecer, de imediato, o ensino profissional.

Assim, em setembro de 1909, o então Presidente do Brasil, Nilo Peçanha, mediante Decreto-Lei n° 7.566, cria nas capitais dos estados da república, as Escolas de Aprendizes Artífices para o ensino profissional primário e gratuito.

A Instituição, batizada com a denominação de Escola de Aprendizes Artífices do Ceará, foi instalada no dia 24 de maio de 1910, na Av. Alberto Nepomuceno, onde funciona, atualmente, a Secretaria Estadual da Fazenda.

Em 1930 o governo provisório assume o poder e a educação passa a ser regulada pelo Ministério da Educação e Saúde Pública (MESP). As Escolas de Aprendizes Artífices, anteriormente ligadas ao Ministério da Agricultura, passaram, por consequência e de imediato, ao MESP e a receber subsídios do governo central.

Em 1937, na reforma do Ministério da Educação e Saúde Pública, o ministro Capanema, mediante a Lei n° 378 de 13 de janeiro, transforma as Escolas de Aprendizes Artífices em Liceus Profissionais, recebendo, no Ceará, a denominação de Liceu Industrial de Fortaleza.

Com a eclosão da Segunda Guerra Mundial, em primeiro de setembro de 1939, houve intensa redução na importação de produtos estrangeiros. Por esta razão, o Brasil passou a cuidar da implantação de indústrias básicas, incentivando a criação de estabelecimentos fabris e, consequentemente, adotou uma política paralela de incentivo à formação de mão-de-obra qualificada, para atender ao incipiente parque industrial.

Por despacho do Ministro da Educação, em 28 de agosto de 1941, houve uma outra modificação no nome dos Liceus. No Ceará, a denominação passou a ser Liceu Industrial do Ceará, nome que durou apenas um ano depois, em 1942, de acordo com o Decreto nº 4121, de 25 de fevereiro, recebeu o nome de Escola Industrial de Fortaleza.

A conjuntura nacional e internacional despertou o interesse do governo brasileiro em modernizar e melhorar o ensino profissional.

Em 1942, a Lei Orgânica do Ensino Industrial estabeleceu as bases da organização e do regime do ensino destinado à preparação profissional dos trabalhadores na indústria e definiu o ensino industrial como de 2º grau, em paralelo com o ensino secundário. Os cursos técnicos de três anos preparariam os alunos para uma nova modalidade de educação, que seria a formação técnica de segundo grau para a área industrial como atribuição das escolas técnicas industriais, que
naquele ano iniciaram suas atividades.

No estado do Ceará, a denominação Escola Técnica Federal do Ceará surge mediante a Lei nº 3.552 de 16 de fevereiro de 1953, alterada pelo Decreto-Lei nº 196, de 27 de agosto de 1969 vinculada ao MEC por intermédio da Secretaria de Educação Médio e Tecnológica - SEMTEC. É uma autarquia educacional, tendo se firmado no Estado como instituição de excelência no ensino técnico-profissional.

Cumpre salientar que tantas mudanças de nome foram decorrentes do sempre renovado papel da Instituição, para uma constante sintonia com os novos horizontes que eram delineados pela permanente dinâmica do progresso muito acelerada nas últimas décadas. A Escola Técnica Federal do Ceará teve inclusive seu campo de ação ampliado com a criação das UNED’s (Unidades Descentralizadas de Ensino) de Cedro e de Juazeiro do Norte (1994), viabilizando o ensino profissional em outras regiões do Estado.

A velocidade do desenvolvimento industrial do país e a inserção gradual de tecnologias avançadas demandam a formação de especialistas de diversos níveis, impondo um persistente reestudo na formação desses profissionais. Deste reestudo nascem os CEFET’s (Centros Federais de Educação Tecnológica) tendo por objetivo ministrar ensino em nível superior de graduação e pós-graduação, visando à formação de profissionais nas áreas de construção civil, industrial e tecnológica, a formação de professores e especialistas para o ensino médio e de formação profissional, formação de técnicos, promoção de cursos de extensão, aperfeiçoamento, atualização profissional e realização de pesquisas na área
técnico-industrial.

A denominação de Centro Federal de Educação Tecnológica do Ceará (CEFET-CE) foi oficializada pela Lei nº 8.948, de 8 de dezembro de 1994 e regulamentada pelo Decreto-Lei nº 2.406, de 27 de novembro de 1997 e pelo Decreto de 22/03/99 (DOU
de 22/03/99) que implantou a nova institucionalidade.

A necessidade de capacitação de novos profissionais levou o Governo Federal a sancionar a lei 11.892/08 que transformou os CEFET’s, Escolas Agrotécnicas e Técnicas em Institutos Federais de Educação, Ciência e Tecnologia (IF’s). Com o mesmo status das universidades federais, os IF’s serão obrigados a oferecer 20\% das vagas para a formação de professores, ou seja, os cursos de licenciaturas.

Os IF’s representam uma nova concepção da educação profissional e humana no Brasil e traduzem o compromisso do governo federal com os jovens e adultos. Esta nova rede de ensino tem um modelo institucional em que as unidades possuem autonomia administrativa e financeira. A nova instituição terá também forte inserção na área de Pesquisa e Extensão para estimular o desenvolvimento de soluções técnicas e tecnológicas.

O Instituto Federal de Educação, Ciência e Tecnologia do Ceará (IFCE) é uma Autarquia Educacional pertencente à Rede Federal de Ensino. Hoje, com 32 Campi, o Instituto Federal do Ceará se consolida como instituição de ensino inclusivo e de qualidade, cuja missão é produzir, disseminar e aplicar os conhecimentos científicos e tecnológicos na busca de participar integralmente da formação do cidadão, visando sua inserção social, política, cultural e ética.  O IFCE valoriza o compromisso ético com responsabilidade social, o respeito, a transparência e a excelência, em consonância com os preceitos básicos de cidadania e humanismo, com liberdade de expressão, cultura da inovação e idéias pautadas na sustentabilidade ambiental.

%\clearpage\section{BREVE HISTÓRICO DO CAMPUS MARACANAÚ}
\subsection{Breve Histórico do Campus Maracanaú}

Quando sancionada pelo Presidente Luis Inácio Lula da Silva, a Lei 11.892 criou trinta e oito Institutos Federais de Educação, Ciência e Tecnologia em 2008. O IFCE nasceu com nove Campi, dentre estes destaca-se o Campus Maracanaú, previamente estabelecido como CEFET - Uned Maracanaú em 2006.

Maracanaú é um município do Estado do Ceará que integra a região metropolitana de Fortaleza e constitui o maior Distrito Industrial do Ceará, caracterizada por um crescente contingente de empresas dos mais diversos setores, indústrias que vão desde o gênero alimentício e têxtil até a indústria metal-mecânica. O município de Maracanaú conta com aproximadamente 18.000 empresas instaladas (segundo dados do Governo do Estado do Ceará), ocupando a quarta posição no estado neste quesito.

O IFCE – Campus Maracanaú foi criado com o intuito de atender a demanda de mão de obra qualificada para as empresas do Estado do Ceará, favorecido por sua localização, estando mais próximo das indústrias em desenvolvimento e das já
existentes. Os cursos ofertados pelo campus Maracanaú sempre estiveram alinhados ao arranjo produtivo da região onde está inserido. Seu primeiro curso oferecido foi o curso técnico em Desenvolvimento de Software, ainda em 2006 como CEFET - Uned Maracanaú. Atualmente, o campus de Maracanaú do IFCE oferta ao todo nove cursos, distribuídos entre os níveis técnico e superior. Na pós-graduação, o campus possui o mestrado em Energias Renováveis e também atua no mestrado em Ciência da Computação, que funciona no campus de Fortaleza.

Neste contexto, o Eixo Tecnológico da Indústria, com apoio da Diretoria Geral do Campus Maracanaú, constituiu uma comissão regulamentada pelas portarias  \PORTARIACOMISSAO (ANEXO \ref{anexoPortariaComissao}) e  \PORTARIACOMISSAOALT (ANEXO \ref{anexoPortariaComissaoAlt}), para implantar 
%conforme resolução n° 020 de 13 de novembro de 2008, reproduzida no ANEXO \ref{anexoResolucao},
o \NOMECURSO, por considerar ser uma função estratégica dos setores industrial e de prestação de serviços.

% TÓPICOS ESPECÍFICOS DE CADA CURSO -------------------------------------------
% JUSTIFICATIVA
% FUNDAMENTAÇÃO LEGAL
% OBJETIVOS DO CURSO

\if\COMPILAR0
    \input{PPC-ECAPerfilCurso}
\else
    \input{PPC-EMEPerfilCurso}
\fi

% FORMAS DE INGRESSO -----------------------------------------------------------

\clearpage
\thispagestyle{empty}
\section{FORMAS DE INGRESSO}

%\subsection{Forma de Acesso ao Curso}

São ofertadas, anualmente, \NUMVAGAS vagas para ingresso no \NOMECURSO. As vagas são preenchidas por meio do Sistema de seleção Unificada (SISU), com base nas notas obtidas pelos estudantes no ENEM (Exame Nacional do Ensino Médio) ou através de edital para ingresso como diplomados ou transferidos, seja por transferência interna ou externa, conforme estabelecido nas seções I, II (Subseções I, II, III e IV), III, IV e V do Capitulo I, Título III, do Regulamento da Organização Didática (ROD) de junho de 2015, ANEXO \ref{anexoRODtransferencia}.

% ÁREAS DE ATUAÇÃO -----------------------------------------------------------

\clearpage\section{ÁREAS DE ATUAÇÃO}

\if\COMPILAR0
    \input{PPC-ECAAreasAtuacao}
\else
    \input{PPC-EMEAreasAtuacao}
\fi

% PERFIL PROFISSIONAL DO EGRESSO ---------------------------------------------

\clearpage\section{PERFIL PROFISSIONAL DO EGRESSO}

\if\COMPILAR0
    \input{PPC-ECAPerfilEgresso}
\else
    \input{PPC-EMEPerfilEgresso}
\fi

% METODOLOGIA ----------------------------------------------------------------

\clearpage\section{METODOLOGIA}

Os procedimentos metodológicos usados em cada disciplina estão especificados, em linhas gerais, nos respectivos Programas de Unidade Didática (PUDs), mas dependem, adicionalmente, das características de cada professor. A grande maioria dos professores opta por aulas expositivas, conforme as necessidades de cada disciplina, com auxílio de quadro branco e pincel, intercaladas com o uso de projeções, aulas de exercícios, práticas em laboratórios, salas de informática, ou ainda visitas a setores do próprio campus ou externas a este. Recursos adicionais também estarão presentes, como o uso de ferramentas de simulação numérica (em determinadas disciplinas específicas) ou tecnologias que garantam a acessibilidade de docentes e/ou discentes (quando necessário).

Obedecendo a Resolução CNE/CES 11, de 11 de março DE 2002, as disciplinas do \NOMECURSO estão divididas em núcleo de conteúdos básicos, um núcleo de conteúdos profissionalizantes e um núcleo de conteúdos específicos. A descrição detalhada sobre cada núcleo é apresentada no item 10.1 que trata da Organização curricular.

Algumas disciplinas de conteúdo básico, como Introdução à Engenharia, HST, Ética e Projetos Sociais, assim como algumas disciplinas de gestão, respectivamente listadas nas Tabs. \ref{tab5} e \ref{tab6}, usarão diferentes metodologias, tais como: trabalho em equipe, seminários de apresentação de projetos por parte dos alunos, pesquisas diversas, trabalho de campo, convivência industrial no caso de estágios, entre outras atividades. Além disso, como já acontece nas outras disciplinas, o professor disponibiliza um horário extra-classe para tira-dúvidas e assim efetuar um melhor acompanhamento da aprendizagem dos alunos.

Tendo consciência que o perfil do egresso deve estar antenado com o mundo do trabalho em constante mutação, busca-se formar um profissional com o conhecimento específico de sua profissão, mas também com uma visão do todo. Este profissional deve saber buscar conhecimento a todo momento, ficar atento a novas tecnologias e desenvolvimentos, e possuir habilidades de comunicação com outras áreas, liderança, administração, entre outras. Neste contexto, o aluno através da interdisciplinaridade existente no curso tem a possibilidade de cursar várias disciplinas do \NOMECURSO em outras engenharias e/ou outras áreas distintas, favorecendo uma troca de experiências e uma visão mais ampla durante sua formação. Além disso, a interdisciplinaridade ocorre entre as mais variadas disciplinas do curso em que um conteúdo é útil para uma outra disciplina subsequente e assim sucessivamente. Neste sentido, foi inserido um campo nos PUDs e na matriz do curso para que seja destacada a interdisciplinaridade de cada disciplina, como forma de deixar claro ao aluno, bem como de incentivar o professor a buscar que esta seja efetivamente alcançada.

Os conteúdos pertinentes às políticas de educação ambiental (lei nº 9.795, de 27 de abril de 1999 e Decreto Nº 4.281 de 25 de junho de 2002), de educação em direitos humanos (Resolução Nº 1, de 30 de maio de 2012) e de educação das relações étnico-raciais (Resolução CNE/CP N° 1, de 17 de junho de 2004), conforme as normatizações vigentes, são contemplados em sua magnitude, da seguinte forma:
%Os conteúdos obrigatórios, recomendados pelas Resoluções CNE Nº2, de 15 de Junho de 2012, CNE Nº1, de 30 de maio de 2012 e CNE Nº1, de 17 de junho de 2004, são contemplados em sua magnitude da seguinte forma:

\begin{itemize}
\item \textbf{Educação Ambiental:} Disciplina de Ética e Cidadania 
\item \textbf{Relações Étnicas-Raciais:} Disciplina de Projetos Sociais 
\item \textbf{Direitos Humanos:} Disciplina de Projetos Sociais
\end{itemize}

A forma de abordagem das temáticas mencionadas, além da abordagem enquanto conteúdos nos componentes curriculares citados, também poderá levar em consideração alguns aspectos, a saber: incentivo a pesquisas envolvendo as temáticas; desenvolvimento de projetos de extensão; organização de eventos, palestras e realização de visitas técnicas.

Nas disciplinas que envolvem práticas de laboratório, haverá um contato maior com os equipamentos didáticos, colocando o aluno em contato direto com os fenômenos físicos, envolvendo ainda recursos de informática para a aquisição e tratamento de dados, bem como para a confecção de relatórios. Em geral, o aluno de \NOMEENG dispõe de um grande arsenal de ferramentas de informática que vão auxiliar diretamente em seus estudos, juntamente com os recursos da Internet, da Biblioteca, bem como da Biblioteca Virtual (\href{http://bvu.ifce.edu.br/login.php}{http://bvu.ifce.edu.br/login.php}). 

Adicionalmente, um elevado número de alunos faz parte dos corpos técnicos dos laboratórios e grupos de pesquisa, dispondo de computadores, envolvimento com trabalhos correlatos, que permitem um melhor acompanhamento das disciplinas cursadas naquele momento. 

Finalmente, o uso das TIC's (Tecnologia da Informação e Comunicação) oferecem um conjunto de recursos tecnológicos no processo de ensino-aprendizagem, tais como e-mail, grupos de discussão on-line e redes sociais em salas de vídeo conferência que possibilitam o intercâmbio de informações e a geração de novos conhecimentos e competências entre todos os envolvidos. Além disso, estas também garantem a acessibilidade de docentes e/ou discintes com necessidades especiais.

Atividades de monitoria são realizadas sob orientação de um professor orientador para alunos que estejam com dificuldade de aprendizagem e, assim, contribuir para um maior envolvimento dos alunos com o IFCE, propiciando uma melhor formação acadêmica ao aluno, além de estimulá-los à participação no processo educacional e nas atividades relativas ao ensino.

Durante o processo de ensino-aprendizagem e seguindo as orientações legais dispostas nos Decretos N° 5.296/2004, N° 6.949/2009, N° 7.611/2011, na Portaria N° 3.284/2003, na Lei N° 12.764, de 27 de dezembro de 2012, são identificados possíveis estudantes com necessidades especiais, os quais, terão um tratamento diferenciado, com o devido apoio ao discente prestado por uma equipe multidisciplinar do campus.   

Outras atividades de apoio ao discente são comentadas no tópico \ref{APOIO}, APOIO AO DISCINTE.

% ESTRUTURA CURRICULAR ----------------------------------------------------------

\clearpage\section{ESTRUTURA CURRICULAR}
A estrutura curricular do \NOMECURSO é dividida nos seguintes itens: 

\begin{itemize}
\item Organização Curricular 
\item Matriz curricular 
\item Fluxograma
\item Estágio
\item Atividades complementares 
\item Trabalho de conclusão de curso. 
\end{itemize}

\subsection{Organização Curricular}
Na organização curricular do curso de \NOMECURSO, foram considerados os seguintes princípios, considerando os objetivos do curso e perfil profissional do egresso a ser formado, a saber: 
\begin{itemize}
    \item Ensino com uma formação básica bastante sólida, fornecida por um conjunto de disciplinas obrigatórias fundamentais para a área do \NOMECURSO ;
    \item Flexibilidade Curricular: permitir que o futuro profissional tenha uma formação básica forte e que complemente esta formação com disciplinas optativas e atividades diversas como estágios, iniciação científica, entre outras, na sua área de interesse específico.
    \item Atualidade: permitir que novas tecnologias e novos conceitos sejam facilmente agregados ao currículo através de disciplinas de caráter optativo;
    \item Qualidade da Formação: além das atividades didáticas em sala de aula, o currículo prevê uma série de outras atividades, como estágios, trabalho de conclusão do curso (TCC), disciplinas integradoras, atividades de iniciação científica, que buscam o aperfeiçoamento individual do aluno e o seu amadurecimento como um profissional especializado, mas com sólida formação básica.
    \item Interdiciplinaridade: várias disciplinas do núcleo básico do \NOMECURSO podem ser cursadas em outros cursos do IFCE, além dos conteúdos abordados em uma disciplina servirem de base para várias outras subsequentes.
\end{itemize}

De um modo geral, a proposta curricular é formada por um conjunto de disciplinas obrigatórias, bem como um conjunto de disciplinas optativas, com as cargas definidas na Tab. \ref{tab4}, fruto da atuação conjunta e contínua do Núcleo Docente Estruturante (NDE) e Colegiado do curso.

\clearpage
\thispagestyle{empty}
    \begin{longtable}{|p{0.3cm}|p{7.1cm}|p{2.7cm}|p{1.8cm}|p{1.0cm}|}
    \caption{Distribuição da Carga Horária do Curso}\label{tab4}
    
    \vspace{-0.25cm} \\
    %\multicolumn{5}{c}{\textbf{Tabela 1. Distribuição da Carga Horária do Curso}}\label{tab4} \\ 
    \hline 
    \cellcolor{\cinza} & \cellcolor{\cinza}Atividades & \cellcolor{\cinza}máximo de horas-aula semanais & \cellcolor{\cinza}horas-aula totais & \cellcolor{\cinza}[\%] \\ 
    \hline
     1 &  Disciplinas Obrigatórias  &  \CHTOTALDOBSE  &  \CHTOTALDOB & \CHTOTALDOBPE \\ \hline
     2 &  Disciplinas Optativas Obrigatórias  &  \CHOPTMINSE  &  \CHOPTMIN & \CHOPTMINPE \\ \hline
     3 &  Estágio  &  \CHESTAGIOSE  &  \CHESTAGIO & \CHESTAGIOPE \\ \hline
     4 &  Total Obrigatório  &  \CHTOTALOBSE  &  \CHTOTALOB & \CHTOTALOBPE \\ \hline
%     5 &  Disciplinas Optativas  &  \CHOPTSE  &  \CHOPT & \CHOPTPE \\ \hline
     5 &  Atividades Complementares Não Obrigatórias &  \CHCOMPSE  &  \CHCOMP & \CHCOMPPE \\ \hline
      %&  Total Geral  &  \CHTOTALGERALSE  &  \CHTOTALGERAL & \CHTOTALGERALPE \\ \hline
   \end{longtable}

%Como é visto na Tab.\ref{tab4}, a carga horária total geral é de \CHTOTALGERALSE horas semanais, mesmo se o discente cursar uma disciplina optativa, o estágio e uma atividade complementar no corrente semestre. O percentual de 100\% é obtido somando-se as atividades complementares não obrigatórias e o restante da carga horária de disciplinas optativas não apresentadas na Tab.\ref{tab4}.

A Tab.\ref{tab5} apresenta o agrupamento das disciplinas de conteúdo básico, que fazem parte da proposta curricular. As Tabelas \ref{tab6} e \ref{tab7} apresentam, respectivamente, as disciplinas de conteúdo profissionalizante e de conteúdo específico (optativas). As denominações de conteúdos básico, profissionalizante e específico, seguem as denominações usadas na \citeonline{cne11}. As disciplinas de conteúdo básico também podem ser consideradas integradoras do ponto de vista do campus, pois são comuns a todos os cursos de engenharia, fazendo parte de uma integração entre os alunos do campus já nos semestres iniciais. Ressalta-se mais uma vez a atuação conjunta e contínua do Núcleo Docente Estruturante (NDE) e Colegiado do curso, que definem tanto as disciplinas que fazem parte da matriz do curso, quanto a respectiva carga horária de cada uma destas.

    \begin{longtable}{|p{0.8cm}|p{7.5cm}|p{2.0cm}|}
    \caption{Disciplinas de conteúdo básico}\label{tab5} 
    \vspace{-0.25cm}
    %\multicolumn{3}{c}{\textbf{Tabela 2. Disciplinas de conteúdo básico}}\label{tab5} 
    \\ \hline 
    \cellcolor{\cinza} & \cellcolor{\cinza}Disciplinas & \cellcolor{\cinza}horas-aula semanais  \\ 
    \hline
     1 &  Introdução à Engenharia  &  2 \\ \hline
     2 &  Cálculo I  &  4 \\ \hline
     3 &  Cálculo II  &  4 \\ \hline
     4 &  Cálculo III  &  4 \\ \hline
     5 &  Física I  &  4 \\ \hline
     6 &  Física II  &  4 \\ \hline
     7 &  Física III  &  4 \\ \hline
     8 &  Álgebra Linear  &  4 \\ \hline
     9 &  Química Geral  &  4 \\ \hline
     10 &  Inglês Técnico  &  2 \\ \hline
     11 &  Probabilidade e Estatística  &  4 \\ \hline
     12 &  Desenho Técnico  &  4 \\ \hline
     13 &  Desenho Auxiliado por Computador  &  4 \\ \hline
     14 &  Lógica de Programação  &  4 \\ \hline
     15 &  Linguagem de Programação  &  4 \\ \hline
     16 &  HST  &  2 \\ \hline
     17 &  Metodologia Científica e Tecnológica  &  2 \\ \hline
     18 &  Ética e Cidadania  &  2 \\ \hline
     19 &  Projetos Sociais  &  2 \\ \hline
      &  Total / (\% em relação à carga total obrigatória)  &  \CHBASICACR(\CHBASICAPEMIN) \\ \hline
   \end{longtable}

\if\COMPILAR0
    \begin{longtable}{|p{0.8cm}|p{7.5cm}|p{2.0cm}|}
    \caption{Disciplinas de conteúdo profissionalizante}\label{tab6} 
    \vspace{-0.25cm}
    %\multicolumn{3}{c}{\textbf{Tabela 3. Disciplinas de conteúdo profissionalizante}}\label{tab6} 
    \\ \hline 
    \cellcolor{\cinza} & \cellcolor{\cinza}Disciplinas & \cellcolor{\cinza}horas-aula semanais  \\ 
    \hline
     1 &  Metrologia  &  4 \\ \hline
     2 &  Eletrônica I  &  4 \\ \hline
     3 &  Materiais I  &  4 \\ \hline
     4 &  Circuitos Elétricos I  &  4 \\ \hline
     5 &  Métodos Numéricos  &  4 \\ \hline
     6 &  Eletrônica II  &  4 \\ \hline
     7 &  Circuitos Elétricos II  &  4 \\ \hline
     8 &  Instrumentação  &  4 \\ \hline
     9 &  Sistemas Lineares  &  4 \\ \hline
     10 &  Resistência dos Materiais I  &  4 \\ \hline
     11 &  Microcontroladores  &  4 \\ \hline
     12 &  Eletrônica III  &  4 \\ \hline
     13 &  Instalações Elétricas  &  4 \\ \hline
     14 &  Controle I  &  4 \\ \hline
     15 &  Processamento Digital de Sinais  &  4 \\ \hline
     16 &  Controladores Lógicos Programáveis  &  4 \\ \hline
     17 &  Máquinas Elétricas  &  4 \\ \hline
     18 &  Acionamentos Hidráulicos e Pneumáticos  &  4 \\ \hline
     19 &  Empreendedorismo  &  2 \\ \hline
     20 &  Controle II  &  4 \\ \hline
     21 &  Acionamento de Máquinas  &  4 \\ \hline
     22 &  Dispositivos Periféricos  &  4 \\ \hline
     23 &  Robótica I  &  4 \\ \hline
     24 &  Sistemas Digitais de Controle Distribuído  &  4 \\ \hline
     25 &  Manufatura Auxiliada por Computador  &  4 \\ \hline
     26 &  Gestão e Controle da Qualidade  &  4 \\ \hline
     27 &  Gestão da Manutenção Industrial &  4 \\ \hline
%     28 &  Tópicos Especiais em Engenharia I  &  4 \\ \hline
%     29 &  Tópicos Especiais em Engenharia II  &  4 \\ \hline
     28 &  TCC  &  2 \\ \hline
     29 &  Estágio Currícular  &  20 \\ \hline
      &  Total / (\% em relação à carga total obrigatória)  &  \CHPROFCR(\CHPROFPEMIN) \\ \hline
   \end{longtable}

%    \vspace{1cm}

    \begin{longtable}{|p{0.8cm}|p{7.5cm}|p{2.0cm}|}
    \caption{Disciplinas de conteúdo específico (optativas)}\label{tab7} 
    \vspace{-0.25cm}
    %\multicolumn{3}{c}{\textbf{Tabela 4. Disciplinas de conteúdo específico (optativas)}}\label{tab7} 
    \\ \hline 
    \cellcolor{\cinza} & \cellcolor{\cinza}Disciplinas & \cellcolor{\cinza}horas-aula semanais  \\ 
    \hline
     1 &  Manufatura Integrada por Computador  &  4 \\ \hline
     2 &  Modelagem de Sistemas a Eventos Discretos  &  4 \\ \hline
     3 &  Processamento Digital de Imagens  &  4 \\ \hline
     4 &  Robótica II  &  4 \\ \hline
     5 &  Identificação de Sistemas  &  4 \\ \hline
     6 &  Sistemas Embarcados  &  4 \\ \hline
     7 &  Controle III  &  4 \\ \hline
     8 &  Equações Diferênciais Ordinárias  &  4 \\ \hline
     9 &  Inteligência Computacional Aplicada  &  4 \\ \hline
     10 &  Mecânica dos Fluidos  &  4 \\ \hline
     11 &  Libras  &  4 \\ \hline
     12 &  Educação Física  &  2 \\ \hline
      &  Total / (\% em relação à carga total obrigatória)  &  \CHOPTCR(\CHOPTPEMIN) \\ \hline
   \end{longtable}

%    \begin{longtable}{|p{0.8cm}|p{7.5cm}|p{2.0cm}|}
%    \multicolumn{3}{c}{\textbf{Tabela 5. Disciplinas integradoras}}
%    \label{tab8} \\ \hline 
%    \cellcolor{\cinza} & \cellcolor{\cinza}Disciplinas & \cellcolor{\cinza}horas-aula semanais  \\ 
%    \hline
%     1 &  Estágio Currícular  &  20 \\ \hline
%     2 &  Libras (optativa)  &  2 \\ \hline
%     3 &  Educação Física (optativa)  &  2 \\ \hline
%      &  Total / (\% em relação à carga total obrigatória)  &  \CHINTCR(\CHINTPEMIN) \\ \hline
%   \end{longtable}
\else
    \begin{longtable}{|p{0.8cm}|p{7.5cm}|p{2.0cm}|}
    \caption{Disciplinas de conteúdo profissionalizante}\label{tab6} 
    \vspace{-0.25cm}
    %\multicolumn{3}{c}{\textbf{Tabela 3. Disciplinas de conteúdo profissionalizante}}    \label{tab6} 
    \\ \hline 
    \cellcolor{\cinza} & \cellcolor{\cinza}Disciplinas & \cellcolor{\cinza}horas-aula semanais  \\ 
    \hline
     1 &  Metrologia  &  4 \\ \hline
     2 &  Eletrônica I  &  4 \\ \hline
     3 &  Materiais I  &  4 \\ \hline
     4 &  Circuitos Elétricos I  &  4 \\ \hline
     5 &  Métodos Numéricos  &  4 \\ \hline
     6 &  Materiais II  &  4 \\ \hline
     7 &  Processos de Fabricação I  &  4 \\ \hline
     8 &  Processos de Fabricação II  &  4 \\ \hline
     9 &  Mecanismos  &  4 \\ \hline
     10 &  Resistência dos Materiais I  &  4 \\ \hline
     11 &  Resistência dos Materiais II  &  4 \\ \hline
     12 &  Termodinâmica  &  4 \\ \hline
     13 &  Soldagem  &  4 \\ \hline
     14 &  Instalações Elétricas  &  4 \\ \hline
     15 &  Manutenção Industrial  &  4 \\ \hline
     16 &  Mecânica dos Fluidos  &  4 \\ \hline
     17 &  Sistemas Mecânicos  &  4 \\ \hline
     18 &  Máquinas Elétricas  &  4 \\ \hline
     19 &  Empreendedorismo  &  2 \\ \hline
     20 &  Transferência de Calor  &  4 \\ \hline
     21 &  Acionamentos Hidráulicos e Pneumáticos  &  4 \\ \hline
     22 &  Planejamento e Controle da Produção  &  4 \\ \hline
     23 &  Manufatura Auxiliada por Computador  &  4 \\ \hline
     24 &  Refrigeração Industrial  &  4 \\ \hline
     25 &  Análise de Falhas  &  4 \\ \hline
     26 &  Máquinas Térmicas e de Fluxo  &  4 \\ \hline
     27 &  Gestão e Controle da Qualidade  &  4 \\ \hline
     28 &  Gestão da Manutenção Industrial &  4 \\ \hline
     29 &  Inspeção Industrial  &  4 \\ \hline
%     29 &  Tópicos Especiais em Engenharia I  &  4 \\ \hline
%     30 &  Tópicos Especiais em Engenharia II  &  4 \\ \hline
     30 &  TCC  &  2 \\ \hline
     31 &  Estágio Currícular  &  20 \\ \hline
      &  Total / (\% em relação à carga total obrigatória)  &  \CHPROFCR(\CHPROFPEMIN) \\ \hline
   \end{longtable}

%    \vspace{1cm}

    \begin{longtable}{|p{0.8cm}|p{7.5cm}|p{2.0cm}|}
    \caption{Disciplinas de conteúdo específico (optativas)}\label{tab7} 
    \vspace{-0.25cm}
    %\multicolumn{3}{c}{\textbf{Tabela 4. Disciplinas de conteúdo específico (optativas)}}\label{tab7} 
    \\ \hline 
    \cellcolor{\cinza} & \cellcolor{\cinza}Disciplinas & \cellcolor{\cinza}horas-aula semanais  \\ 
    \hline
     1 &  Corrosão e Proteção Anti-Corrosiva  &  4 \\ \hline
     2 &  Tribologia e Lubrificação  &  4 \\ \hline
     3 &  Transportadores Industriais  &  4 \\ \hline
     4 &  Gestão de Projetos  &  4 \\ \hline
     5 &  Projeto de Engenharia  &  4 \\ \hline
     6 &  Manufatura Integrada por Computador  &  4 \\ \hline
     7 &  Processamento Digital de Imagens  &  4 \\ \hline
     8 &  Equações Diferênciais Ordinárias  &  4 \\ \hline
     9 &  Inteligência Computacional Aplicada  &  4 \\ \hline
     10 &  Libras  &  4 \\ \hline
     11 &  Educação Física  &  2 \\ \hline
      &  Total / (\% em relação à carga total obrigatória)  &  \CHOPTCR(\CHOPTPEMIN) \\ \hline
   \end{longtable}

%    \begin{longtable}{|p{0.8cm}|p{7.5cm}|p{2.0cm}|}
%    \multicolumn{3}{c}{\textbf{Tabela 5. Disciplinas integradoras}}
%    \label{tab8} \\ \hline 
%    \cellcolor{\cinza} & \cellcolor{\cinza}Disciplinas & \cellcolor{\cinza}horas-aula semanais  \\ 
%    \hline
%     1 &  Estágio Currícular  &  20 \\ \hline
%     2 &  Libras (optativa)  &  2 \\ \hline
%     3 &  Educação Física (optativa)  &  2 \\ \hline
%      &  Total / (\% em relação à carga total obrigatória)  &  \CHINTCR(\CHINTPEMIN) \\ \hline
%   \end{longtable}

\fi

Apesar do \NOMECURSO ser anual, ressalta-se que as disciplinas de conteúdo básico são ofertadas todo semestre pelo campus, sejam elas no próprio curso, como também por outros cursos de engenharia. Tal política de combate à retenção e, consequentemente, da evasão, dá oportunidade ao discente, caso haja reprovação, desistência ou trancamento, de fazer tais disciplinas no semestre subsequente à perda da mesma. 

\subsection{Matriz Curricular}

%\subsubsection{Matríz}
%A seguir estão listadas todas as disciplinas que compõem a

A Matriz Curricular do \NOMECURSO do IFCE Campus Maracanaú, apresentada na Tab. \ref{MATRIZ}, contempla todas as disciplinas do curso. De um modo geral, o curso está organizado em 10 (dez) semestres letivos, cada um com intervalo de tempo de 100 dias de atividades de ensino. Ao longo de cada semestre são trabalhadas Unidades Curriculares, cada uma com seus objetivos e conjunto de habilidades a serem desenvolvidas, de forma que, ao final do curso, o graduando possa desenvolver as competências e habilidades necessárias à formação do futuro profissional.

A Matriz Curricular do \NOMECURSO contempla as disciplinas básicas, específicas e profissionalizantes necessárias para uma boa formação teórica e prática, conforme Tab. \ref{MATRIZ}. Ao longo do curso, o discente deverá cursar pelo menos duas disciplinas de conteúdo específico (optativas), perfazendo um total de 160 horas. Vale ressaltar que, a partir da demanda de discentes e docentes, procura-se que em todos os semestres seja ofertada pelo menos uma das disciplinas optativas previstas, de forma que todo o conteúdo optativo seja disponibilizado aos discentes, sendo esta uma das ações para garantir a flexibilidade curricular. Ressalva-se ainda, que devem ser respeitados os pré-requisitos exigidos, conforme elencados na matriz curricular.

%A Matriz Curricular do \NOMECURSO, Tabela \ref{MATRIZ}, contempla as disciplinas básicas, específicas e profissionalizantes necessárias para uma boa formação teórica e prática. Pelo menos duas disciplinas de Tópicos Especiais (160 horas) garantem a flexibilidade da abordagem de assuntos específicos e de interesse para a formação dos alunos. As referidas disciplinas equivalem aos conteúdos optativos a serem ofertados
%conforme a demanda por parte dos alunos e professores. Na matriz que segue na Tabela \ref{MATRIZ}, essas disciplinas estão dispostas nos semestres 8 e 9 como sugestão de carga horária mínima de disciplinas optativas. No entanto, ao longo de cada semestre sempre são ofertadas de uma a duas disciplinas optativas, de forma que todo o conteúdo optativo seja disponibilizado aos discentes, garantindo assim uma maior flexibilização do curso. Ressalva-se que o conteúdo ofertado deve respeitar os pré-requisitos exigidos.

O aluno que obtiver aprovação em qualquer das disciplinas específicas ou optativas de cursos de nível equivalente, ofertadas no Eixo Tecnológico da Indústria, terá direito a aproveitar a referida disciplina no \NOMECURSO.

Em obediência ao Decreto Federal n° 5.626, de 22 de dezembro de 2005, Capítulo II, Artigo 3°, § 2°, a matriz curricular apresenta a disciplina de Linguagem Brasileira de Sinais – LIBRAS para \NOMECURSO. O ANEXO \ref{anexoDecreto}, apresenta o referido Decreto que em seu Capítulo II, Artigo 3°, § 2° diz: ``A Libras constituir-se-á em disciplina curricular optativa nos demais cursos de educação superior e na educação profissional, a partir do ano da publicação deste Decreto''.

As componentes curriculares do \NOMECURSO estão detalhadas no Programa de Unidade Didática (PUD) das disciplinas. O PUD é um documento que descreve os pré-requisitos exigidos, carga horária (total, teórica e prática), número de créditos, período, interdisciplinaridade, temas atuais, ementa, objetivos, programa, metodologia, avaliação, recursos e as bibliografias básica e complementar.

O PUD das disciplinas é atualizado sempre que for detectada a necessidade de melhorias, adequando a disciplina à realidade dos alunos do curso e também às exigências do mercado de trabalho.

A relação completa dos PUDs das disciplinas do \NOMECURSO é listada no ANEXO \ref{anexoPUD}.

%\subsubsection{Matrícula}
A matrícula é requerida pelo interessado e operacionalizada por Unidades Curriculares no prazo estabelecido em calendário escolar do Campus Maracanaú. O regime de matrícula consta nas Seções I e II, Capitulo II, Título III, do Regulamento da Organização Didática (ROD) de junho de 2015, ANEXO \ref{anexoRODmatricula}.

A escolha de disciplinas só pode ser realizada a partir do segundo semestre, onde, no primeiro, o aluno é obrigatoriamente matriculado em todas as disciplinas do referido semestre. O período para integralização do curso é de 10 semestres letivos.

%\newpage

% matriz do curso
\if\COMPILAR0
    
 \begin{longtable}{|c|c|l|c|c|p{1.5cm}|}
 \caption{MATRIZ CURRICULAR ECA (5719)} \label{MATRIZ}
 
 \vspace{-0.25cm}
 %\multicolumn{6}{c}{\textbf{Tabela 5. MATRIZ CURRICULAR ECA (5719)}} \label{MATRIZ}
 
 \\ \hline 

 \cellcolor{\cinza}Sem & \cellcolor{\cinza}Cod & \cellcolor{\cinza}Disciplina & \cellcolor{\cinza}CH & \cellcolor{\cinza}CR & \cellcolor{\cinza}PR \\ 

 \hline 
\hline

S1&  \hyperref[PUD:ACE002]{04507.1} & Álgebra Linear & 80 & 4 &  \\ 

S1&  \hyperref[PUD:ACE003]{04507.2} & Cálculo 1 & 80 & 4 &  \\ 

S1&  \hyperref[PUD:ACE004]{04507.3} & Inglês Técnico & 40 & 2 &  \\ 

S1&  \hyperref[PUD:ACE005]{04507.4} & Química Geral & 80 & 4 &  \\ 

S1&  \hyperref[PUD:ACE017]{04507.5} & Introdução à Engenharia & 40 & 2 & \\ 

S1&  \hyperref[PUD:EME105]{04507.6} & Desenho Técnico & 80 & 4 & \\ 

& & & \fbox{\textbf{400}} & & \\ 
\hline

%S2&  \hyperref[PUD:ACE006]{04507.7} & Cálculo 2 & 80 &  \hyperref[PUD:ACE002]{04507.1}  \hyperref[PUD:ACE003]{04507.2}  \\ 
S2&  \hyperref[PUD:ACE006]{04507.7} & Cálculo 2 & 80 & 4 &  \hyperref[PUD:ACE003]{04507.2} \\ 

S2&  \hyperref[PUD:ACE008]{04507.8} & Física 1 & 80 & 4 &  \hyperref[PUD:ACE003]{04507.2} \\ 

S2&  \hyperref[PUD:ACE013]{04507.9} & Probabilidade e Estatística & 80 & 4 &  \hyperref[PUD:ACE003]{04507.2} \\ 

S2&  \hyperref[PUD:ECA208]{04507.10} & Lógica de Programação & 80 & 4 &  \\ 

S2&  \hyperref[PUD:EME106]{04507.11} & Desenho Auxiliado por Computador & 80 & 4 &  \hyperref[PUD:EME105]{04507.6}  \\ 

& & & \fbox{\textbf{400}} & & \\ 
\hline

S3&  \hyperref[PUD:ACE010]{04507.12} & Cálculo 3 & 80 & 4 &  \hyperref[PUD:ACE006]{04507.7} \\ 

S3&  \hyperref[PUD:ACE011]{04507.13} & Física 2 & 80 & 4 &  \hyperref[PUD:ACE008]{04507.8} \\ 

S3&  \hyperref[PUD:ECA209]{04507.14} & Linguagem de Programação & 80 & 4 &  \hyperref[PUD:ECA208]{04507.10} \\ 

S3&  \hyperref[PUD:ECA211]{04507.15} & Eletrônica I & 80 & 4 & \\ 

S3&  \hyperref[PUD:EME107]{04507.16} & Metrologia & 80 & 4 & \\ 

& & & \fbox{\textbf{400}} & & \\ 
\hline

S4&  \hyperref[PUD:ACE009]{04507.17} & HST & 40 & 2 & \\ 

S4&  \hyperref[PUD:ACE012]{04507.18} & Física 3 & 80 & 4 &  \hyperref[PUD:ACE011]{04507.13} \\ 

S4&  \hyperref[PUD:ACE014]{04507.19} & Metodologia Científica e Tecnológica & 40 & 2 &  \\ 

S4&  \hyperref[PUD:ECA210]{04507.20} & Circuitos Elétricos I & 80 & 4 &  \hyperref[PUD:ACE003]{04507.2} \\ 

S4&  \hyperref[PUD:ECA212]{04507.21} & Métodos Numéricos & 80 & 4 &  \hyperref[PUD:ECA209]{04507.14}  \\ 

S4&  \hyperref[PUD:EME108]{04507.22} & Materiais & 80 & 4 &  \hyperref[PUD:ACE005]{04507.4}  \\ 

& & & \fbox{\textbf{400}} & & \\ 
\hline

S5&  \hyperref[PUD:ECA214]{04507.23} & Eletrônica II & 80 & 4 &  \hyperref[PUD:ECA210]{04507.20} \\ 

S5&  \hyperref[PUD:ECA215]{04507.24} & Circuitos Elétricos II & 80 & 4 &  \hyperref[PUD:ECA210]{04507.20} \\ 

S5&  \hyperref[PUD:ECA219]{04507.25} & Instrumentação & 80 & 4 & \\ 

S5&  \hyperref[PUD:ECA220]{04507.26} & Sistemas Lineares & 80 & 4 &  \hyperref[PUD:ACE002]{04507.1} \\ 

S5&  \hyperref[PUD:EME109]{04507.27} & Resistência dos Materiais & 80 & 4 &  \hyperref[PUD:EME108]{04507.22}  \\ 

& & & \fbox{\textbf{400}} & &  \\ 
\hline

S6&  \hyperref[PUD:ECA213]{04507.28} & Microcontroladores & 80 & 4 &  \hyperref[PUD:ECA209]{04507.14}  \\ 

S6&  \hyperref[PUD:ECA216]{04507.29} & Eletrônica III & 80 & 4 &  \hyperref[PUD:ECA214]{04507.23} \\ 

S6&  \hyperref[PUD:ECA222]{04507.30} & Instalações Elétricas & 80 & 4 &  \hyperref[PUD:ECA210]{04507.20} \\ 

S6&  \hyperref[PUD:ECA233]{04507.31} & Controle I & 80 & 4 & \hyperref[PUD:ECA220]{04507.24} \\ 

& & & \fbox{\textbf{320}} & & \\ 
\hline

S7&  \hyperref[PUD:ACE007]{04507.33} & Ética e Cidadania & 40 & 2 & \\ 

S7&  \hyperref[PUD:ECA205]{04507.34} & Processamento Digital de Sinais & 80 & 4 &  \hyperref[PUD:ECA220]{04507.26}\\ 

S7&  \hyperref[PUD:ECA221]{04507.35} & Controladores Lógicos Programáveis & 80 & 4 & \hyperref[PUD:ECA213]{04507.28}\\ 

S7&  \hyperref[PUD:ECA223]{04507.36} & Máquinas Elétricas & 80 & 4 & \hyperref[PUD:ACE012]{04507.18} \\ 

S7&  \hyperref[PUD:ECA228]{04507.37} & Acionamentos Hidráulicos e Pneumáticos & 80 & 4 &  \hyperref[PUD:ECA214]{04507.30}  \\ 

& & & \fbox{\textbf{360}} & & \\ 
\hline

S8&  \hyperref[PUD:ACE015]{04507.38} & Empreendedorismo & 40 & 2 & \\ 

%S8&  \hyperref[PUD:ACE018]{04507.39} & Tópicos Especiais em Engenharia I & 80 &  \\ 

S8&  \hyperref[PUD:ECA202]{04507.40} & Controle II & 80 & 4 & \hyperref[PUD:ECA233]{04507.31} \\ 

S8&  \hyperref[PUD:ECA224]{04507.41} & Acionamento de Máquinas & 80 & 4 & \hyperref[PUD:ECA233]{04507.31} \\ 

S8&  \hyperref[PUD:ECA225]{04507.42} & Dispositivos Periféricos & 80 & 4 &  \hyperref[PUD:ECA213]{04507.28}  \\

& & & \fbox{\textbf{280}} & & \\ 
\hline

%S9&  \hyperref[PUD:ACE019]{04507.43} & Tópicos Especiais em Engenharia II & 80 &  \\ 

S9&  \hyperref[PUD:ECA226]{04507.44} & Robótica I & 80 & 4 &  \hyperref[PUD:ECA225]{04507.42}  \\ 

S9&  \hyperref[PUD:ECA227]{04507.45} & Sistemas Digitais de Controle Distribuído & 80 & 4 &  \\ 

S9&  \hyperref[PUD:EME123]{04507.46} & Manufatura Auxiliada por Computador & 80 & 4 &  \hyperref[PUD:EME106]{04507.11} \\ 

& & & \fbox{\textbf{240}} & & \\ 
\hline

S10&  \hyperref[PUD:ACE016]{04507.47} & Projetos Sociais & 40 & 2 &  \\ 

S10&  \hyperref[PUD:ACE020]{04507.48} & TCC & 40 & 2 & \\ 

S10&  \hyperref[PUD:EME120]{04507.49} & Gestão e Controle da Qualidade & 80 & 4 &  \\ 

S10&  \hyperref[PUD:EME127]{04507.50} & Gestão da Manutenção Industrial & 80 & 4 & \\ 

& & & \fbox{\textbf{240}} & & \\ 
\hline

Opt&  \hyperref[PUD:ACE001]{04507.51} & Libras & 80 & 4 & \\ 

Opt&  \hyperref[PUD:ECA201]{04507.52} & Manufatura Integrada por Computador & 80 & 4 & \\ 

Opt&  \hyperref[PUD:ECA203]{04507.53} & Modelagem de Sistemas a Eventos Discretos & 80 & 4 & \\ 

Opt&  \hyperref[PUD:ECA204]{04507.54} & Processamento Digital de Imagens & 80 & 4 &  \hyperref[PUD:ECA209]{04507.14} \\ 

Opt&  \hyperref[PUD:ECA206]{04507.55} & Robótica II & 80 & 4 & \hyperref[PUD:ECA226]{04507.44} \\ 

Opt&  \hyperref[PUD:ECA218]{04507.56} & Identificação de Sistemas & 80 & 4 &  \hyperref[PUD:ECA233]{04507.31} \\ 

Opt&  \hyperref[PUD:ECA230]{04507.57} & Sistemas Embarcados & 80 & 4 &  \hyperref[PUD:ECA225]{04507.42} \\ 

Opt&  \hyperref[PUD:ECA235]{04507.58} & Controle III & 80 & 4 & \hyperref[PUD:ECA202]{04507.40} \\ 

%Opt&  \hyperref[PUD:ECA234]{04507.32} & \color{gray} Projeto Eletrônico Auxiliado por Computador & 80 &  \\ 
Opt&  \hyperref[PUD:ACE021]{04507.59} & Educação Física & 40 & 2 & \\ 

Opt&  \hyperref[PUD:EME131]{04507.60} & Equações Diferenciais Ordinárias & 80 & 4 &  \hyperref[PUD:ACE006]{04506.7} \\ 

Opt&  \hyperref[PUD:EME132]{04507.61} & Inteligência Computacional Aplicada & 80 & 4 &  \hyperref[PUD:ECA209]{04507.14} \\ 

Opt&  \hyperref[PUD:EME118]{04507.62} & Mecânica dos Fluidos & 80 & 4 &  \hyperref[PUD:ACE011]{04507.13} \\ 

& & & \fbox{\textbf{\CHOPT}} & & \\ 
 \hline 

 & &\color{gray} 47 Discisplinas Obrigatórias & \color{gray}\CHTOTALDOB & \color{gray}\CHTOTALDOBCR cr. & \\ 

 & &\color{gray} 12 Discisplinas Optativas & \color{gray}\CHOPT & \color{gray}\CHOPTCR cr. & \\ 

 & &\color{gray} 29 Discisplinas Profissionalizantes + Estágio & \color{gray}\CHPROF & \color{gray}\CHPROFCR cr. & \\ 

 & &\color{gray} 19 Discisplinas Comuns/Básicas & \color{gray}\CHBASICA & \color{gray}\CHBASICACR cr. & \\ 

 \hline 

 \end{longtable} 

%\footnotesize{\textbf{Fonte:} A Comissão de Elaboração.}
%\label{MATRIZ}
\else
    \input{MatrizEME}%\label{MATRIZ}
\fi

\subsection{Fluxograma}
Para um melhor entendimento da Matriz Curricular do \NOMECURSO, apresenta-se a disposição gráfica da estrutura curricular contendo a carga horária do componente curricular, a quantidade de créditos e o fluxo de pré-requisitos.

\if\COMPILAR0
	\includepdf[pages={1}]{Diagramas_Matriz_EME_ECA_AUT_MAN/CursoECA_Celso.pdf}
\else
	\includepdf[pages={1}]{Diagramas_Matriz_EME_ECA_AUT_MAN/CursoEME_Celso.pdf}
\fi

\subsection{Estágio Curricular}

Dentre as inovações propostas na organização curricular dos Cursos Superiores de Tecnologia, destaca-se o Estágio Supervisionado com carga horária mínima de 400 horas, a ser cursado em empresas relacionadas à área de formação do profissional. Essa atividade é obrigatória, ficando a diplomação do aluno(a) condicionada à realização da mesma.

O Estágio Supervisionado pode ser realizado após a conclusão do 5{\r{}} semestre. Neste momento o estudante faz seu primeiro contato com a realidade da empresa, saindo do ambiente acadêmico com seus princípios teóricos e vislumbrando a complexidade daquele novo mundo, suas tecnologias, procedimentos, cultura e ambiente. Neste contexto a teoria é colocada à prova e a capacidade de relacionamento do estudante é exigida.

O Estágio Supervisionado tem como finalidades principais:

\begin{itemize}
\item Esclarecer às diversas realidades no ambiente de trabalho;
\item Motivar o aluno ao permitir que ele possa avaliar o confronto “teoria x prática”;
\item Propiciar uma consciência das suas necessidades teóricas e comportamentais;
\item Criar uma visão geral do setor produtivo e da empresa em especial;
\item Identificar áreas de interesse para a sua própria especialização no decorrer e após o término do curso.
\end{itemize}

O aluno será acompanhado por um professor orientador de estágio que terá uma dedicação de 4 horas mensalmente, dentro do período letivo estabelecido pela instituição, destinadas ao acompanhamento do estágio. Essa carga horária é distribuída na forma de reuniões que podem ser realizadas na empresa ou no próprio Instituto Federal de Educação, Ciência e Tecnologia do Ceará – Campus Maracanaú. As reuniões devem sempre ocorrer com a apresentação de um relato das atividades que ele
está realizando e do desempenho apresentado na execução dessas atividades.

Ao término do estágio o aluno deverá apresentar um Relatório Final, até 7 (sete) dias antes do término do período letivo estabelecido pela instituição de ensino.

A avaliação final do estágio será feita pelo professor orientador de estágio através dos conceitos SATISFATÓRIO ou INSATISFATÓRIO, considerando a avaliação da empresa, a compatibilidade das atividades executadas com o currículo da habilitação e a coerência das atividades desenvolvidas na carga horária prevista.

Em caso de parecer INSATISFATÓRIO o professor orientador de estágio poderá pedir ao estagiário um novo relatório ou a realização de um novo estágio.

Em atendimento aos documentos demandados pela Portaria MEC No 544, de 16 de junho de 2020; Ofício Conjunto Circular nº 03-2020 PROEN-PROEXT-PRPI-REITORIA - 01 de junho 2020; Ofício Conjunto Circular no 5/2020 PROEN/PROEXT/PRPI/REITORIA - Plano de Trabalho Específico; Ofício Conjunto Circular nº 03-2020 PROEN-PROEXT-PRPI-REITORIA - 01 de junho 2020; Ofício Conjunto Circular nº 1/2021 PROEN/PRPI/PROEXT/REITORIA de 17 de março de 2021 e Ofício-Circular nº 3/2021/CEAE/DEA/PROEXT/REITORIA-IFCE de 30 de março de 2021, o Colegiado do curso autorizou a realização dos estágios obrigatórios e não obrigatórios de forma remota conforme ATA de reunião do Colegiado de 05/03/2021, conforme Plano de trabalho anexado a este PPC (ANEXO II: PLANO DE TRABALHO ESPECÍFICO - PRÁTICA DE ESTÁGIO OBRIGATÓRIO/NÃO OBRIGATÓRIO).

\subsection{Atividades Complementares}

%As atividades educacionais complementares devem privilegiar a construção de comportamentos sociais e profissionais que
%as atividades acadêmicas tradicionais, de sala de aula ou de laboratório não têm condições de propiciar.
%
%Nesta perspectiva, podem ser inseridas as atividades de cunho comunitário e de interesse coletivo e também serem
%privilegiadas atividades de monitoria acadêmica e de iniciação científica ou tecnológica que propiciem a participação
%do estudante na vida da instituição. Podem aqui também serem desenvolvidas atividades esportivas e culturais, além de
%intercâmbios com instituições estrangeiras.
%
%São atividades diversas, de cunho acadêmico, tecnológico e cultural, que fazem parte da vida escolar do aluno e que são
%relacionadas com o exercício profissional. Segundo a Resolução CES/CNE 11 (MEC, 2002), poderão também ser estimuladas
%atividades complementares tais como trabalhos de iniciação científica, projetos multidisciplinares, visitas técnicas,
%trabalhos em equipe, desenvolvimento de protótipos, monitorias, participação em empresas juniores e outras atividades
%empreendedoras.
%
%As Atividades Complementares ou Extra-Curriculares do \NOMECURSO constituem um conjunto
%de atividades didático-pedagógicas que permitem, no âmbito do currículo, a articulação entre teoria e prática e a
%complementação dos saberes e habilidades necessárias, a serem desenvolvidas durante o período de formação do
%profissional.
%
%É facultada ao aluno a participação em atividades complementares e a carga horária equivalente não será computada para
%efeitos de conclusão do curso.

As Atividades Complementares ou Extra-Curriculares do \NOMECURSO constituem um conjunto de atividades didático-pedagógicas que permitem, no âmbito do currículo, a articulação entre teoria e prática, bem como a complementação das habilidades e saberes necessários a serem desenvolvidos durante o período de formação do profissional.

Tratam-se de atividades diversas, de cunho acadêmico, tecnológico e cultural, que fazem parte da vida escolar do aluno e que são relacionadas com o exercício profissional. De acordo com a \citeonline{cne11}, poderão também ser estimuladas atividades complementares tais como trabalhos de iniciação científica, projetos multidisciplinares, visitas técnicas, trabalhos em equipe, desenvolvimento de protótipos, monitorias, participação em empresas juniores e outras atividades empreendedoras.

As atividades educacionais complementares devem privilegiar a construção de comportamentos sociais e profissionais que as atividades acadêmicas tradicionais, de sala de aula ou de laboratório, não têm condições de propiciar.

Nesta perspectiva, podem ser inseridas as atividades de cunho comunitário e de interesse coletivo, além de privilegiar atividades de monitoria acadêmica e de iniciação científica ou tecnológica que propiciem a participação do estudante na vida da instituição. Podem aqui também serem desenvolvidas atividades esportivas e culturais, além de intercâmbios com instituições estrangeiras.

%Embora haja todo o incentivo, existe a preocupação de que todo aluno siga essas atividades, sem exceção, como forma de enriquecer a sua formação. Portanto, é facultada ao aluno a participação em atividades complementares e a carga horária equivalente não será computada para efeitos de conclusão do curso.

A carga mínima de atividades complementares não obrigatórias recomendada para o curso abrange um total de \CHCOMP horas-aula. Na sequência são descritas as atividades complementares mais comuns, os critérios para o aproveitamento de carga horária e as exigências para aproveitamento destas atividades complementares. 

Vale destacar que é possível e até incentivada a inclusão e/ou retirada de atividades complementares, como forma do curso se adaptar às novas exigências do mundo do trabalho, sendo que isto será definido pelo NDE e Colegiado do curso.

\subsubsection{Visitas técnicas}

Acontecem a partir do primeiro semestre cursado, com o intuito de facilitar o processo ensino-aprendizagem das
disciplinas cursadas para garantir um bom aproveitamento da mesma.

As visitas técnicas a empresas do Distrito Industrial de Maracanaú e da região metropolitana de Fortaleza são realizadas semestralmente. Uma vez por ano é realizada uma visita técnica a uma empresa de grande porte localizada em regiões fora do estado do Ceará.

\subsubsection{Feiras, Seminários, Congressos e Semanas Tecnológicas}

Os alunos são estimulados a participarem de Seminários, Congressos, Palestras e a participação como Monitor (Auxiliar) em Eventos. Alunos de iniciação científica tem seus trabalhos publicados em Eventos de nível nacional e internacional, participando como apresentadores.

%\subsubsection{Semana de Tecnologias  Industriais (SET-IN)}

O Eixo Tecnológico da Indústria do Campus Maracanaú, com o apoio da direção geral desta unidade e da reitoria do IFCE, realiza a SETIN (Semana de Tecnologias Industriais), que é um evento anual de atividades de pesquisa e extensão. O evento também conta com o apoio de outras instituições e parceiros, tais como o IFCE – Campus Fortaleza, a AEDI (Associação das Empresas dos Distritos Industriais do Estado do Ceará) e o SEBRAE (Serviço Brasileiro de Apoio às Micro e Pequenas Empresas).


\subsubsection{Programa de Monitoria e Bolsas de Trabalho}

A monitoria é uma atividade desenvolvida por alunos de graduação, integrantes de projetos orientados para a diminuição dos índices de evasão e repetência, como também para a melhoria do padrão de qualidade dos cursos de graduação, coordenados por docentes. Além dos monitores bolsistas, remunerados com recursos orçamentários do IFCE, outros alunos podem se integrar aos projetos aprovados, na condição de monitores voluntários.

A Coordenação do \NOMECURSO, juntamente com o Diretoria de Ensino do Campus IFCE / Maracanaú, tem envidado esforços no sentido de fortalecer a componente prática da formação dos nos alunos, futuros engenheiros. Pela própria especificidade do Curso, uma integração eficiente entre a teoria e a prática no processo ensino-aprendizagem é indispensável à formação, com qualidade, dos profissionais exigidos pelo mercado de trabalho. Além disso, as atividades de caráter experimental se constituem, indubitavelmente, em fortes elementos de motivação para os estudantes em nível de Graduação.

O trabalho experimental possibilita o contato e a familiarização com equipamentos, montagens, circuitos, dispositivos e instrumentos de medição. Propicia a comprovação, no laboratório, dos conhecimentos teóricos adquiridos na sala de aula ou por outros meios. Permite ao estudante compreender as limitações e nuances dos modelos teóricos em face da prática de situações reais. Tais aspectos são fundamentais à formação do engenheiro, em particular do \NOMEEGRESSO. A atividade  experimental, instigando o interesse pela investigação científica, também contribui para despertar vocações para a pesquisa.

As disciplinas em que os monitores geralmente atuam constituem a base indispensável ao preparo dos alunos do Curso para o prosseguimento e aprofundamento dos seus estudos no campo da \NOMEENG. Evidencia-se a necessidade de que seja fortalecida a atividade de Monitoria no \NOMECURSO, ao lado de outras iniciativas objetivando incrementar a integração teoria-prática.

No \NOMECURSO, o Programa de Monitoria e de bolsas de trabalho tem os seguintes objetivos principais:

\begin{itemize}
\item Proporcionar um maior equilíbrio entre teoria e prática no \NOMECURSO, contribuindo para a formação de engenheiros capacitados a enfrentar e resolver problemas colocados pela realidade;
\item Fortalecer a componente experimental das disciplinas teóricas-práticas, em particular as de formação básica;
\item Motivar os monitores e demais alunos no estudo das disciplinas - não raro excessivamente teóricas - objetivando a redução dos níveis de evasão no Curso;
\item Permitir a redução do número de alunos em cada turma de laboratório - viabilizada pela presença de monitores - o que corresponderá a um melhor rendimento, com consequente melhoria da qualidade do ensino ministrado;
\item Propiciar o surgimento e florescimento de vocações para a docência e a pesquisa, além de promover a cooperação acadêmica entre discentes e docentes.
\end{itemize}

%Atualmente o \NOMECURSO conta com três (3) bolsistas de monitoria e com 20 (vinte) bolsistas de trabalho atuando no apoio às atividades laboratoriais do curso.

\subsubsection{Iniciação Científica com Bolsa ou de Forma Voluntária}

A iniciação científica é a atividade complementar mais importante desenvolvida no curso, onde
o aluno passa a fazer parte de uma equipe de pesquisa, tornando-se responsável pelo
desenvolvimento de um tema. Esse tema se encaixa em um trabalho maior,
envolvendo outros alunos de graduação e de mestrado. O aluno passa a
aprender técnicas não desenvolvidas em sala de aula e passa a se especializar em
determinadas áreas. Além do conhecimento adquirido, existe um grande progresso em
nível individual, quanto à capacidade de trabalho, independência e responsabilidade.

O IFCE oferece Bolsas de Iniciação Científica através dos Programas Institucionais de Bolsas de Iniciação Científica sendo elas do Conselho Nacional de Desenvolvimento Científico e Tecnológico – PIBIC/CNPq, ou da Fundação Cearense de Apoio ao Desenvolvimento Científico e Tecnológico, PIBIC/FUNCAP ambos destinado aos pesquisadores do IFCE com titulação de doutor, para as cotas PIBIC/FUNCAP. As bolsas oferecidas pelo Programa Institucional de Bolsas de Iniciação Científica do Instituto Federal de Educação, Ciência e Tecnologia do Ceará – PIBIC/IFCE são destinadas aos pesquisadores do IFCE com titulação de doutor, mestre ou especialista para as cotas PIBIC/IFCE.

Segundo a conceituação formal do CNPq, ``o PIBIC é um programa centrado na iniciação científica de novos talentos em todas as áreas do conhecimento, administrado diretamente pelas instituições. Voltado para o aluno de graduação e servindo de incentivo à formação, privilegia a participação ativa de bons alunos em projetos de pesquisa com qualidade acadêmica, mérito científico e orientação adequada, individual e continuada. Os projetos culminam com um trabalho final avaliado e valorizado, fornecendo retorno imediato ao bolsista, com vistas à continuidade de sua formação, de modo particular na pós-graduação''.

Além disso, o CNPq menciona que as bolsas de iniciação científica permitem que pesquisadores produtivos engajem estudantes de cursos superiores no processo acadêmico, otimizando a capacidade de orientação à pesquisa na instituição; promovem o aumento da produção científica, com o envolvimento de novos orientadores nas atividades de iniciação à pesquisa científica. Despertam vocação científica e incentivam talentos potenciais entre estudantes de cursos superiores, mediante suas participações em projetos de pesquisa, introduzindo o jovem graduando no domínio do método científico; proporcionam ao bolsista, orientado por pesquisador qualificado, a aprendizagem de técnicas e métodos científicos, bem como estimulam o desenvolvimento do pensar científico e da criatividade, decorrentes das condições criadas pelo confronto direto com os problemas de pesquisa; Despertam no bolsista uma nova mentalidade em relação à pesquisa além de preparar os estudantes para a pós-graduação.

Também são ofertadas bolsas de fomento na Iniciação em Desenvolvimento Tecnológico e Inovação do Conselho Nacional de Desenvolvimento Científico e Tecnológico – PIBITI/CNPq e do Programa Institucional de Bolsas de Iniciação em Desenvolvimento Tecnológico e Inovação do Instituto Federal de Educação, Ciência e Tecnologia do Ceará – PIBITI/IFCE destinados aos pesquisadores do IFCE com titulação de doutor, ou perfil equivalente, e, para as cotas PIBITI/CNPq, pesquisador com titulação de doutor, mestre ou especialista para as cotas PIBITI/IFCE.

As bolsas de Iniciação em Desenvolvimento Tecnológico e Inovação, PIBITI, propiciam à instituição um instrumento de formulação de sua política de inovação tecnológica, através da iniciação tecnológica na graduação, contribuem para a formação e a inserção de estudantes em atividades de pesquisa, desenvolvimento tecnológico e inovação, formação e o engajamento de recursos humanos para atividades de pesquisa, desenvolvimento tecnológico e inovação; contribuem para a formação de recursos humanos que se dedicarão ao fortalecimento da capacidade inovadora das empresas no País, formação do cidadão pleno, com condições de participar de forma criativa e empreendedora na sua comunidade; possibilitam maior interação entre atividades de desenvolvimento tecnológico e inovação, desenvolvidas na graduação e na pós-graduação além de envolver os pesquisadores nas atividades de formação de desenvolvimento tecnológico e inovação.

\subsubsection{Grupos de Pesquisa}

Hoje o \NOMECURSO conta com o Núcleo de Pesquisa Aplicada e Inovação Industrial (NAI)
que contempla os grupos de pesquisa que dão suporte aos alunos no desenvolvimento de pesquisa em nível de iniciação científica. Atualmente são 3 grupos de pesquisa cadastrados no CNPq e certificados pela instituição.

\vspace{0.5cm}

\noindent \textbf{1. Grupo de Inspeção e Análise de Falhas (GIAF)}

Atua no desenvolvimento de pesquisa em inspeção e análise de falhas aplicado ao setor industrial, visando a melhoria contínua dos processos, produtos manufaturados e plantas industriais, de forma a maximizar a vida útil dos equipamentos empregados na indústria. Contribui para a formação de profissionais especializados em inspeção e manutenção. O GIAF conta com a participação de 5 (cinco) pesquisadores e 4 (quatro) estudantes do \NOMECURSO como bolsistas de iniciação científica e 1 (um) estudante como voluntário.

As linhas de pesquisa do GIAF são:
\begin{itemize}
\item Análise Microestrutural de Metais
\item Corrosão
\item Inspeção e Ensaios em Materiais
\item Manutenção Preventiva e Preditiva
\end{itemize}

\noindent \textbf{2. Grupo de Pesquisa em Sistemas Inteligentes (GPSI)}

Atua no desenvolvimento de inovações científicas e tecnológicas aplicadas às áreas de Automação Industrial, Energia e Manutenção Industrial. Contribui para a melhoria contínua dos processos, plantas industriais e para a formação de profissionais especializados, atendendo às necessidades regionais e ao crescimento do parque industrial. O GPSI conta com a participação de 11 (onze) pesquisadores e 4 (quatro) estudantes do \NOMECURSO como bolsistas de iniciação científica e 1 (um) estudante como voluntário.

As linhas de pesquisa do GPSI são:
\begin{itemize}
\item Aproveitamento de Energias Alternativas e Eficiência Energética
\item Controle e Automação Aplicados a Processos de Fabricação
\item Inspeção e Manutenção Preditiva
\end{itemize}


\noindent \textbf{3. Grupo de Pesquisa Energias Renováveis (GPER)}

Atua no desenvolvimento de inovações científicas e tecnológicas aplicadas às áreas de Automação Industrial, Energia, Manutenção Industrial. Contribui para a melhoria contínua dos processos de geração de energia, atuando no desenvolvimento tecnológico e sustentável do parque industrial, visando diminuir os impactos da crise de energia e meio ambiente, além de contribuir para a formação de profissionais especializados, atendendo às necessidades regionais. O GPER conta com a participação de 11 (onze) pesquisadores e 4 (quatro) estudantes do \NOMECURSO como bolsistas de iniciação científica e 1 (um) estudante como voluntário.

As linhas de pesquisa do GPER são:
\begin{itemize}
\item Controle e Processamento de Energia
\item Mecânica Aplicada à conservação do meio ambiente
\item Energia Eólica
\item Energia fotovoltáica
\end{itemize}

\subsubsection{Bolsa de Monitoria e de Trabalho}

O \NOMECURSO dispõe através de programas de auxílio ao discente, bolsistas para atividades de monitoria e de trabalho atuando no apoio às atividades laboratoriais do curso. A atividade de monitoria é desenvolvida dentro de uma disciplina, por um aluno que já a tenha cursado e que tenha conseguido uma nota mínima de 7,0. Nesta atividade, há o contato com colegas mais novos, desenvolvendo no aluno monitor aspectos mais abrangentes de caráter didático-pedagógico, bem como a necessidade de aprofundamento na disciplina em questão.

\subsubsection{Estágios não Obrigatórios}
Estágios de curta duração também estão disponíveis para o aluno de graduação.
Nesses estágios diferentes empresas e diferentes processos produtivos podem ser
conhecidos, dando um maior embasamento e maior conhecimento no campo de
trabalho futuro do aluno.


%\subsubsection{Viagens de Estudo}
%Atividades como viagens de estudo já são usados no presente momento como elementos motivadores e %instrumentos pedagógicos complementares do curso de graduação. As visitas ocorrem aos parques industriais %do estado do Ceará. A programação é feita dentro do contexto de cada disciplina, havendo o acompanhamento %do professor responsável.

\subsubsection{AeroDesign}

O IFCE, Campus Maracanaú, conta com uma equipe de projeto de aeromodelos, a AirLamp Aerodesign - IFCE. Esta é formada, atualmente, por alunos dos cursos de Engenharia Mecânica e Engenharia de Controle e Automação do campus. A AirLamp visa tirar a Engenharia do papel projetando aeromodelos eficientes, com o menor custo possível, materiais inovadores a fim de competir na competição organizada pela SAE Brasil. A inexistência de um prérequisito tem como objetivo a participação de alunos de diversos níveis, incentivando o trabalho em equipes e uma maior adaptação do aluno ao ambiente do curso.

\subsubsection{Enactus}

A Enactus trata-se de uma organização internacional que mobiliza estudantes, acadêmicos e líderes de negócios que estão comprometidos a usar o poder da ação empreendedora para possibilitar o progresso no mundo. Guiados por professores conselheiros e especialistas em negócios, os estudantes
participantes formam times em seus campi para criar e implementar projetos comunitários que empoderam pessoas para melhorar sua qualidade e padrão de vida. Essa experiência não somente transforma vidas, mas também ajuda os estudantes a desenvolverem o tipo de talento e perspectiva necessários para tornarem-se líderes eficazes e direcionados ao valor.

Uma série anual de campeonatos regionais e nacionais fornece um fórum para que os times possam apresentar os resultados dos seus projetos, e para que sejam avaliados por líderes de negócios, que atuam como juízes. Os campeões nacionais de cada país avançam para a prestigiada Enactus World Cup. Além do aspecto comunitário do programa, oportunidades significativas são criadas para o aprendizado e trocas de experiências entre gerações, além de expor estudantes e ex-alunos às empresas que procuram um novo talento.

O Time Enactus IFCE Maracanaú tem como objetivo ajudar comunidades do entorno do campus em seu desenvolvimento. A nível de instituição, a participação no time é considerada uma atividade de extensão. A inexistência de um prérequisito tem como objetivo a participação de alunos de diversos níveis, incentivando o trabalho em equipes e uma maior adaptação do aluno ao ambiente do campus.


%\subsubsection{Empresa Júnior}
%A empresa júnior de engenharia mecânica – atualmente denominada I9 Consultoria –
%foi criada em 1995 e tem a finalidade de desenvolver as vocações empresariais dos
%estudantes do curso. A empresa atua diretamente no mercado como consultora e
%executora de projetos. O suporte é dado pelos professores e pelos laboratórios do
%EMC, sempre que houver necessidade. Trata-se também de uma atividade integradora,
%onde os estudantes são treinados para o dia a dia de trabalho em suas atividades
%futuras de engenheiros-empreendedores.

%\subsubsection{Programa Especial de Treinamento – PET}
%O Grupo PET Metrologia e Automação da UFSC, foi fundado em 1980 e conta com a
%participação de estudantes de Engª Mecânica. Inclui a organização e a participação
%seminários, palestras e cursos com o intuito de aproximar os estudantes ao mercado
%de trabalho. Tem também como meta a realização de projetos técnico-científicos de
%engenharia, em parceria com pequenas, médias e grandes empresas do setor metalmecânico
%brasileiro.


\subsubsection{Cooperação Internacional}

O Programa de Bolsas IFCE Internacional visa consolidar a internacionalização do IFCE, propiciando a interiorização destas ações, bem como possibilitar a participação de alunos de diferentes níveis de ensino, oportunizando a participação de discentes do ensino técnico cuja oferta para mobilidade internacional é quase inexistente.
A fim de intensificar as atividades já desenvolvidas com instituições de ensino estrangeiras parceiras do IFCE, os discentes selecionados pelo presente programa através de edital serão enviados para cursar um semestre acadêmico em instituições de ensino de excelência em diferentes países.
Além destas parcerias já consolidadas, outras instituições e indústrias têm sido utilizadas pelos alunos, colocando-se atualmente, como uma necessidade para a formação, tanto pelo aprendizado de novas línguas, quanto pelo contato com outras culturas.

\subsubsection{Outras Atividades Complementares}
Tratam-se de outras atividades complementares também aceitas no curso:
\begin{itemize}
    \item Assistir a palestras
    \item Participação como debatedor em eventos na área do curso
    \item Apresentação de trabalhos como expositor em eventos na área
    \item Participação em projetos e programas de extensão promovidos ou não pelo IFCE
    \item Participação em cursos de extensão na área do curso de graduação ministrados ou não pelo IFCE
    \item Participação em cursos de extensão em geral
    \item Participação em atividades ou eventos culturais organizados pelo IFCE ou por outras instituições de Ensino Superior
    \item Participação em atividades nos centros acadêmicos ou diretórios de estudantes
    \item Participação em órgãos de direção de entidade de natureza acadêmica
    \item Representação em colegiados acadêmicos ou administrativos do IFCE
    \item Aprovação em disciplinas conexas
    \item Cursos de ensino a distância em áreas afins ao curso
    \item Outras atividades relativas a quaisquer colaborações em situações acadêmicas
\end{itemize}

Desta-se novamente que é possível e até incentivada ainclusão e/ou retirada de atividades complementares, como forma do curso se adaptar às novasexigências do mundo do trabalho, sendo que isto será definido pelo NDE e Colegiado do curso.

\subsubsection{Carga Horária por Modalidade de Atividade Complementar}

O aproveitamento da carga horária seguirá os critérios estabelecidos na Tab. \ref{tab12} Deverá ser respeitado o limite de carga horária por cada Atividade Complementar descrita. A carga horária que exceder o cômputo geral, de acordo com as modalidades, não será aproveitada.

\begin{longtable}{|p{8.5cm}|p{3.5cm}|p{2cm}|}

\caption{Carga horária por modalidade de atividade complementar}\label{tab12} 

\vspace{-0.25cm}

%\multicolumn{3}{c}{\textbf{Tabela 6. Carga horária por modalidade de atividade complementar}} 

\\ \hline 

 \cellcolor{\cinza}Atividade & \cellcolor{\cinza}CH por Atividade & \cellcolor{\cinza}CH Total \\ 
 
 \hline

 Visita técnica & Até 4h por visita & Até 60h \\ 

 \hline 
 
 Feiras, Seminários, Congressos e Semanas Tecnológicas & Até 20h por evento & Até 60h \\
 
 \hline
 
 Programa de Monitoria e Bolsas de Trabalho & Até 30h por semestre & Até 60h \\
 
 \hline
 
 Iniciação Científica com Bolsa ou de Forma Voluntária & Até 20h por pesquisa & Até 80h \\
 
 \hline
 
 Grupos de Pesquisa & Até 20h por pesquisa & Até 80h \\
 
 \hline
 
 Bolsa de Monitoria e de Trabalho & Até 30h por semestre & Até 60h \\
 
 \hline
 
 Estágios não Obrigatórios & Até 70h & Até 70h \\
 
 \hline
 
 AeroDesign & Até 30h por semestre & Até 60h \\ 
 
 \hline
 
 Enactus & Até 30h por semestre & Até 60h \\
 
 \hline
 
 Cooperação Internacional & Até 40h por disciplina & Até 200h\\
 
 \hline
 
 Assistir a palestras & Até 4h por evento & Até 60h \\
 
 \hline
 
 Participação como debatedor em eventos na área do curso & Até 8h por evento & Até 60h \\
 
 \hline
 
 Apresentação de trabalhos como expositor em eventos na área & Até 20h por trabalho & Até 60h \\
 
 \hline
 
 Participação em projetos e programas de extensão promovidos ou não pelo IFCE & Até 20h por atividade & Até 80h \\
 
 \hline
 
 Participação em cursos de extensão na área do curso de graduação ministrados ou não pelo IFCE & Até 30h por curso & Até 60h \\
 
 \hline
 
 Participação em cursos de extensão em geral & Até 5h por curso & Até 20h \\
 
 \hline
 
 Participação em atividades ou eventos culturais organizados pelo IFCE ou por outras instituições de Ensino Superior & Até 10h por atividade & Até 40h \\
 
 \hline
 
 Participação em atividades nos centros acadêmicos ou diretórios de estudantes & Até 4h por atividade & Até 30h \\
 
 \hline
 
 Participação em órgãos de direção de entidade de natureza acadêmica & Até 10h por semestre & Até 40h \\
 
 \hline
 
 Representação em colegiados acadêmicos ou administrativos do IFCE & Até 10h por semestre & Até 40h \\
 
 \hline
 
 Aprovação em disciplinas conexas & Até 40h por disciplina & Até 80h \\
 
 \hline
 
 Cursos de ensino a distância em áreas afins ao curso & Até 60h & Até 60h \\
 
 \hline
 
 Outras atividades relativas a quaisquer colaborações em situações acadêmicas & Até 40h & Até 40h \\

 \hline
 
 \end{longtable} 

\subsubsection{Aproveitamento de Atividades Complementares}

Ficam estabelecidas as seguintes exigências para o aproveitamento das Atividades Complementares, conforme Tab. \ref{tab13}. O controle do cumprimento dos créditos referentes às Atividades Complementares é de responsabilidade do coordenador do curso, a quem cabe avaliar a documentação exigida para a validação da atividade.

\begin{longtable}{|p{8.5cm}|p{5.5cm}|}

\caption{Aproveitamento de atividades complementares}\label{tab13}

\vspace{-0.25cm}

%\multicolumn{2}{c}{\textbf{Tabela 7. Exigências para aproveitamento de atividades complementares}} 

\\ \hline 

 \cellcolor{\cinza}Atividade & \cellcolor{\cinza}Documentação Necessária \\ 
 
 \hline

 Visita técnica & Declaração emitida pelo professor responsável atestando a quantidade de horas e que a mesma não foi registrada em diário como atividade de alguma disciplina \\ 

 \hline 
 
 Feiras, Seminários, Congressos e Semanas Tecnológicas & Certificado de presença \\
 
 \hline
 
 Programa de Monitoria e Bolsas de Trabalho & Relatório do professor orientador \\
 
 \hline
 
 Iniciação Científica com Bolsa ou de Forma Voluntária & Relatório do professor orientador \\
 
 \hline
 
 Grupos de Pesquisa & Relatório do professor orientador \\
 
 \hline
 
 Bolsa de Monitoria e de Trabalho & Relatório do professor orientador \\
 
 \hline
 
 Estágios não Obrigatórios & Declaração emitida pelo setor de estágios \\
 
 \hline
 
 AeroDesign & Declaração de participação emitida pelos líderes da equipe \\ 
 
 \hline
 
 Enactus & Declaração de participação emitida pelos líderes do time \\
 
 \hline
 
 Cooperação Internacional & Declaração emitida pelo setor responsável \\
 
 \hline
 
 Assistir a palestras & Atestado de participação \\
 
 \hline
 
 Participação como debatedor em eventos na área do curso & Atestado de participação \\
 
 \hline
 
 Apresentação de trabalhos como expositor em eventos na área & Atestado de participação \\
 
 \hline
 
 Participação em projetos e programas de extensão promovidos ou não pelo IFCE & Declaração emitida pelo professor responsável atestando a quantidadede horas e que a mesma não foi registrada em diário como atividade de alguma disciplina \\
 
 \hline
 
 Participação em cursos de extensão na área do curso de graduação ministrados ou não pelo IFCE & Atestado de participação \\
 
 \hline
 
 Participação em cursos de extensão em geral & Atestado de participação \\
 
 \hline
 
 Participação em atividades ou eventos culturais organizados pelo IFCE ou por outras instituições de Ensino Superior & Atestado de participação \\
 
 \hline
 
 Participação em atividades nos centros acadêmicos ou diretórios de estudantes & Atas de reunião e relatório de atividades \\
 
 \hline
 
 Participação em órgãos de direção de entidade de natureza acadêmica & Atas de reunião e relatório de atividades \\
 
 \hline
 
 Representação em colegiados acadêmicos ou administrativos do IFCE & Atas de reunião e relatório de atividades \\
 
 \hline
 
 Aprovação em disciplinas conexas & Histórico constando a nota obtida e o programa da disciplina\\
 
 \hline
 
 Cursos de ensino a distância em áreas afins ao curso & Certificado de conclusão\\
 
 \hline
 
 Outras atividades relativas a quaisquer colaborações em situações acadêmicas & Certificado de realização\\

 \hline
 
 \end{longtable} 

Ao longo do semestre letivo, o aluno deverá apresentar os comprovantes cabíveis e suas respectivas cópias ao coordenador do curso, que os apreciará, podendo recusar a atividade se considerar insatisfatória e/ou o desempenho do aluno. Sendo aceita a atividade realizada pelo aluno, cabe ao coordenador do curso atribuir a carga horária correspondente.

%Quando da apresentação dos comprovantes, o coordenador deverá atestar as cópias, mediante o documento original e arquivá-las na pasta de Atividades Complementares do aluno.

É vedado o cômputo concomitante ou sucessivo, como Atividade Complementar, de cargas horárias ou conteúdos, trabalhos, atividades ou práticas próprias das disciplinas do currículo pleno, ou destinado à elaboração e defesa do TCC, ou desenvolvidos nos estágios curriculares.

De atos ou decisões do coordenador do curso caberá recurso à Direção de Ensino do campus.

\subsection{Trabalho de Conclusão de Curso (TCC)}

O TCC trata-se de um trabalho acadêmico de caráter obrigatório que tem o objetivo de promover a consolidação dos conhecimentos adquiridos no decorrer do curso. Além disso, visa a iniciação e envolvimento do aluno de graduação no campo da pesquisa científica.

Quando da elaboração do trabalho para apresentação, o aluno deve se matricular na disciplina de TCC, sendo que, para tal, é obrigatório que este tenha cumprido no minimo de 80\% dos crétidos obrigatórios constantes na matriz do curso.  
%O aluno deve se matricular na disciplina de TCC no 10° semestre%, mesmo que já venha realizando atividades do trabalho nos semestres anteriores. 
A entrega de um pré-projeto de TCC constitui pré-requisito obrigatório para que possa matricular-se na disciplina. 

Desta feita, o professor da disciplina de TCC, além de cumprir o programa constante no respectivo PUD, atuará como um coordenador do trabalho implementado pelo aluno e seu professor orientador.

O aluno pode optar por apresentar um artigo que tenha sido aceito para publicação em periódico ou artigo que tenha sido apresentado em um evento científico, desde que conste o nome do \textbf{professor orientador} na lista de autores. O artigo será submetido à análise do professor de TCC que emitirá um parecer quanto à relevância e atualidade do trabalho. Caso o artigo seja aceito \textbf{sem alterações} como substituto da monografia, o aluno poderá passar para a fase de apresentação pública do trabalho. Caso contrário, o aluno deverá iniciar o desenvolvimento de uma monografia (baseada no Manual de Normalização de Trabalhos Acadêmicos do IFCE, aprovado através da Resolução 034/Consup, de 27 de março de 2017 e disponível no endereço \web{https://ifce.edu.br/proen/bibliotecas/normalizacao-de-trabalhos-academicos}) ou artigo (científico ou técnico) visando a publicação em periódico de circulação nacional ou internacional com avaliação Qualis A ou B nas áreas de engenharia, computação ou interdisciplinar da CAPES.

O desenvolvimento do trabalho será acompanhado por um professor orientador que proporciona aos alunos subsídios para construir algo, tais como: definição do tema, acompanhamento das atividades práticas e/ou teóricas, revisão da parte escrita e conclusão do trabalho. A disciplina de TCC terá um professor responsável que atuará como coordenador das atividades implementadas pelo aluno e seu professor orientador, o qual acompanhará o cumprimento das seguintes etapas:

\begin{itemize}
\item Entrega do Pré-Projeto do TCC ou artigo publicado;
\item Entrega do termo de aceite de orientação assinado pelo professor orientador;
\item Entrega de relatórios da execução de atividades referentes à monografia ou artigo;
\item Entrega da versão impressa para apresentação pública;
\item Apresentação pública da monografia ou artigo;
\item Entrega da versão final da monografia ou artigo.
\end{itemize}

Os prazos de cada uma das etapas acima serão definidos pelo professor da disciplina de TCC juntamente com a coordenação
do \NOMECURSO. Somente após a conclusão da última etapa pelo aluno será lançada a nota da disciplina no sistema acadêmico.

\subsubsection{Entrega do Pré-Projeto do TCC}

O Pré-Projeto é um documento que apresenta a percepção do aluno quanto ao trabalho que será executado. Este documento é importante para orientar as atividades e permitir um bom planejamento do aluno. A viabilidade do projeto será analisada pelo orientador e pelo professor da disciplina de TCC, o qual poderá sugerir ajustes com o objetivo de melhorar as chances de conclusão no tempo previsto.

O Pré-Projeto deve ser apresentado à coordenação de curso como pré-requisito para matrícula na disciplina de TCC. A estrutura mínima do documento contemplar:

\begin{itemize}
\item Capa;
\item Tema;
\item Problema;
\item Objetivos;
\item Justificativa;
\item Metodologia;
\item Cronograma;
\item Referências Bibliográficas.
\end{itemize}

\subsubsection{Entrega do termo de aceite de orientação assinado pelo professor orientador}

O termo é um compromisso assumido pelo professor orientador e seu respectivo orientando com relação ao cumprimento das
etapas e do cronograma de elaboração e entrega do TCC. É importante ressaltar que o professor orientador é responsável
pela qualidade e garantia de aprovação do TCC. O termo deverá obedecer ao modelo descrito a seguir.

\begin{verbatim}
                         TERMO DE ACEITE DE TCC

             Eu, <nome do professor>, professor do Instituto Federal de Educação,
Ciência e Tecnologia de Ceará – Campus Maracanaú, Eixo Tecnológico da Indústria,
declaro para os devidos fins, que aceito orientar o(a) aluno(a) <nome completo
do(a) aluno(a)> do <nome do curso> na execução do Trabalho de Conclusão de Curso
cujo título é <título do trabalho>.

                                        Maracanaú – CE, dia / mês / ano

                                        ----------------------------------
                                        assinatura do professor orientador

\end{verbatim}

\subsubsection{Entrega de relatórios da execução de atividades referentes ao TCC}

Os relatórios têm o objetivo de auxiliar o acompanhamento das atividades previstas para a elaboração do TCC. O
professor da disciplina de TCC poderá solicitar ao aluno e/ou ao professor orientador que façam ajustes necessários à
boa qualidade do trabalho.

\subsubsection{Entrega da versão para apresentação pública}

 De forma a atender as demandas do professor da disciplina de TCC e da respectiva banca examinadora, devem ser entregues a versão digital (em pdf) do texto para apresentação, bem como cópias impressas com encadernação simples (caso necessárias) ao professor da disciplina de TCC, que terá a responsabilidade de repassar essas a cada membro da banca examinadora.

\subsubsection{Apresentação pública do TCC}

Trata-se de uma seção pública tendo uma banca examinadora composta pelo professor orientador, como presidente, e mais dois membros, sendo pelo menos um deles do Instituto Federal de Educação, Ciência e Tecnologia do Ceará. Um dos membros da banca, à exceção do presidente, pode ser de uma empresa industrial ou de outra instituição, pública ou privada, de ensino superior de graduação em áreas tecnológicas.  Todos os membros da banca devem possuir, pelo menos, pós-graduação em nível de especialização.  A apresentação do TCC é pública, portanto aberta para qualquer membro da sociedade que desejar assistir.

A apresentação é dividida em quatro etapas:

\begin{itemize}
\item Apresentação do trabalho pelo proponente: utiliza recursos multimídia para melhor visualização dos membros da banca e de todas as pessoas que estiverem presentes. O tempo de duração da apresentação deve ser de no máximo 30 (trinta) minutos.
\item Arguições e considerações por parte da banca: após a apresentação do trabalho, cada membro da banca inicia o processo de arguição e considerações, onde são apontadas sugestões para melhoria do trabalho e possíveis correções. Após todos os membros da banca concluírem suas arguições e considerações, o presidente pode determinar um tempo para questionamentos e considerações das pessoas que estão assistindo a defesa.
\item Reunião da banca com o professor da disciplina de TCC: é a última etapa da apresentação e ocorre para que os membros da banca e o professor da disciplina de TCC discutam, de maneira reservada, as características do trabalho apresentado e deliberem pela nota do trabalho.
\item Composição da nota da disciplina de TCC: cada membro da banca juntamente com o professor da disciplina de TCC atribuirão uma nota de 0 (zero) a 10 (dez), em formulário específico, e a média aritmética dessas notas será a nota final da disciplina de TCC. Essa etapa dará origem ao Parecer de Trabalho de Conclusão de Curso, que será o documento oficial a ser considerado para registro da nota final atribuída à disciplina de TCC.
\item Comprovação de participação dos membros da banca: cada membro da banca examinadora receberá uma declaração, pela Coordenação do \NOMECURSO, que comprovará a respectiva participação de cada um deles na defesa do TCC junto ao Instituto Federal de Educação, Ciência e Tecnologia do Ceará e a qualquer outra instituição de natureza pública ou privada.
\end{itemize}

\subsubsection{Entrega da versão final do TCC}

Após a apresentação do TCC, o(a) aluno(a) deve efetuar as correções e/ou melhorias propostas pela banca examinadora. A aceitação da versão final com suas respectivas correções e/ou melhorias será confirmada por meio de um consenso entre o professor orientador e o professor da disciplina de TCC.

A validade das notas atribuídas ao trabalho apresentado está condicionada à entrega %de 01 (uma) cópia impressa da versão final, encadernada em capa dura com letras douradas, assim como a entrega  
do arquivo eletrônico em mídia digital, o qual será disponibilizado para acesso público no Sistema de Bibliotecas do IFCE - SIBI (no endereço eletrônico \web{http://biblioteca.ifce.edu.br/}). É importante ressaltar que nesta versão final deve conter a ficha catalográfica, fornecida pela biblioteca do \textit{campus}, e o Parecer de Trabalho de Conclusão de Curso, fornecido pela coordenação do \NOMECURSO.

A nota da disciplina de TCC será distribuída igualmente na média dos 02 (dois) períodos, N1 e N2, do semestre letivo.

% APROVEITAMENTO E VALIDAÇÃO DO CONHECIMENTO ----------------------------------------------

\clearpage
\section{APROVEITAMENTO E VALIDAÇÃO DO CONHECIMENTO}

No que se refere ao aproveitamento de componentes curriculares cursados, o IFCE assegurará aos estudantes ingressantes e veteranos o direito de aproveitamento, mediante análise, desde que haja compatibilidade de programa e carga horária de, no mínimo, 75\% do total estipulado para o componente curricular a ser aproveitado. O componente curricular apresentado deve estar no mesmo nível de ensino ou em um nível de ensino superior ao do componente curricular a ser aproveitado, devendo ser solicitado no máximo uma vez. 

No aproveitamento, deverão ser considerados os conhecimentos adquiridos não só para as disciplinas do semestre em curso, como também para as de semestres posteriores, no caso de aluno recém-ingresso. Este, terá 10 (dez) dias após a sua matrícula, para requerer o aproveitamento de disciplina. Quanto ao aluno veterano, o aproveitamento será para o semestre/ano posterior, devendo a solicitação ser feita durante os 30 (trinta) primeiros dias do semestre em curso. E devem ser considerados, ainda, os demais critérios de aproveitamento determinados no Titulo III, Capitulo IV, Seção I, do ROD, que trata do aproveitamento de componentes curriculares.

Já no que se refere à validação de conhecimentos, o IFCE validará conhecimentos adquiridos em estudos regulares ou em experiência profissional de estudantes do IFCE com situação de matrícula em matriculado, mediante avaliação teórica ou prática. O requerente poderá estar matriculado ou não no componente curricular para o qual pretende validar conhecimentos adquiridos. 
Não poderá ser solicitada validação de conhecimento para estágio curricular, trabalho de conclusão de curso e atividades complementares, assim como para estudantes que tenham sido reprovados no IFCE no componente curricular cuja validação de conhecimentos foi solicitada.

A solicitação de validação de conhecimentos deverá ser feita mediante requerimento protocolado e enviado à coordenadoria do curso, juntamente com o envio dos seguintes documentos: declaração, certificado ou diploma - para fins de validação em conhecimentos adquiridos em estudos regulares, cópia da Carteira de Trabalho (páginas já preenchidas) ou declaração do empregador ou de próprio punho, quando autônomo - para fins de validação de conhecimentos adquiridos em experiências profissionais anteriores e documentação complementar, caso seja solicitado pela comissão avaliadora.

O calendário do processo de validação de conhecimentos deverá ser instituído pelo próprio campus. Porém, a validação deverá ser solicitada nos primeiros 30 (trinta) dias do período letivo em curso e todo o processo de validação deverá ser concluído em até 50 (cinquenta) dias letivos do semestre vigente, a contar da data inicial de abertura do calendário do processo de validação de conhecimentos, definida pelo campus. 

A validação de conhecimentos de um componente curricular só poderá ser solicitada uma única vez e devem ser considerados, ainda, os demais critérios de aproveitamento determinados no Titulo III, Capitulo IV, Seção II, do ROD, que trata da validação de conhecimentos.

% AVALIAÇÃO DO PROJETO DO CURSO ----------------------------------------------

\clearpage
\section{AVALIAÇÃO DO PROJETO DO CURSO}

A avaliação do curso de \NOMECURSO deve ocorrer de forma contínua e dialogada através da interação de professores, gestores e estudantes. Ao longo do desenvolvimento das atividades curriculares, todos os segmentos, docentes, discentes e técnicos-administrativos, juntamente com a Coordenação do Curso, NDE e Colegiado do curso, devem atuar na direção da consolidação de mecanismos que possibilitem a permanente avaliação dos objetivos do curso.

\subsection{Avaliação do Docente}

O corpo docente que atua no curso passa por avaliações semestrais através de questionários respondidos pelos estudantes, por meio do sistema eletrônico Q-acadêmico, para cada componente curricular nos quais estão regularmente matriculados no semestre letivo. 

No instrumento são observados pontos como Pontualidade, Assiduidade, Domínio de conteúdo, Metodologia de Ensino, Avaliação e Relação Professor-Aluno.

A partir dos resultados obtidos são gerados relatórios para cada docente, na sua área de acesso no Q-Acadêmico, referentes a cada componente curricular, que devem ser percebidos pelo professor como instrumentos de autoavaliação e reflexão sobre a sua prática docente, avaliando-a e, se for o caso, intervir de forma a enriquecê-la e viabilizar melhorias na aprendizagem dos estudantes.

%No ANEXO \ref{anexoAvaliacao}  é apresentado o modelo da avaliação do docente do \NOMECURSO.

\subsection{Núcleo Docente Estruturante}

%No ANEXO \ref{anexoPortariaNDE} é apresentada a portaria \PORTARIANDE referente à criação do Núcleo Docente Estruturante (NDE) do \NOMECURSO e \PORTARIANDEALT referente às alterações no mesmo.

A constituição, funcionamento e atuação do NDE do \NOMECURSO estão em consonância com o disposto na Resolução CONSUP Nº 004, de 28 de janeiro de 2015, que determina a organização do Núcleo Docente Estruturante no IFCE, como sendo:
\begin{itemize}
\item Construir e acompanhar a execução do PPC;
\item Promover a revisão e atualização do PPC, tendo como principal objetivo a adequação do perfil profissional do egresso, devendo as alterações serem aprovadas pela maioria do NDE, e submetidas à análise e aprovação do colegiado do curso;
\item Analisar os resultados obtidos nas avaliações internas e externas (ENADE, Relatório de Avaliação para Reconhecimento de curso) e propor estratégias para o desenvolvimento da qualidade acadêmica do curso;
\item Zelar pelo cumprimento das Diretrizes Curriculares Nacionais para os cursos de graduação.
\end{itemize}

\subsection{Colegiado}

%No ANEXO \ref{anexoPortariasColegiado} é apresentada a portaria  \PORTARIACOLEGIADO referente à criação do Colegiado do \NOMECURSO e \PORTARIACOLEGIADOALT referente às alterações no mesmo.

A constituição, funcionamento e atuação do colegiado do \NOMECURSO estão em consonância com o disposto na Resolução N° 75, de 13 de agosto de 2018, que define as normas de funcionamento dos colegiados de curso do IFCE.
 
Conforme estabelece o Artigo 4° da referida Resolução, compete ao Colegiado do IFCE:

\begin{itemize}
\item Supervisionar as atividades curriculares, propondo aos órgãos competentes as medidas necessárias à melhoria do ensino, pesquisa e extensão;
\item Aprovar as propostas de estruturação e reestruturação do Projeto Pedagógico do Curso;
\item Avaliar o desenvolvimento do Projeto Pedagógico do Curso no tocante a sua atualização, primando pela sintonia com as demandas da sociedade e do mundo do trabalho;
\item Deliberar sobre as recomendações propostas pelos docentes, discentes e egressos sobre assuntos de interesse do curso;
\item Propor soluções para as questões administrativas e pedagógicas do curso, tais como aquelas que tratam de evasão, reprovação, retenção, entre outras;
\item Propor, conforme o caso, a flexibilização curricular, bem como a extinção e a alteração de componentes curriculares;
\item Coletar e analisar informações sobre as diferentes áreas do saber que compõem o curso, incluindo questões de cunho acadêmico;
\item Orientar acerca de qual perfil docente deve ser solicitado, por ocasião de concurso público e/ou de remoção de professores, vislumbrando as necessidades do curso e as características de seu Projeto Pedagógico;
\item Organizar e construir a sequência de afastamento docente no âmbito do curso, bem como deliberar acerca da efetivação deste afastamento, com base na regulamentação vigente;
\item Colaborar, sempre que solicitado, no auxílio, indicação e escolha de membros de banca de concurso público, junto à Comissão Coordenadora de Concurso da Instituição;
\item Receber, analisar e encaminhar demandas do corpo docente e discente e tomar decisões de natureza didático-pedagógicas sobre elas, desde que atendam à legislação em vigor.
\end{itemize}

\subsection{Encontros Pedagógicos}

Os Encontros Pedagógicos são realizados semestralmente, como uma ação formativa, em ambientes produtivos onde são preconizadas as práticas docentes e sua melhoria, cujo objetivo é provocar no professor a avaliação de sua prática docente para que ao longo do processo melhorem sua atuação pedagógica.

\subsection{Ações decorrentes dos processos de autoavaliação e avaliação externa}

De acordo com a Lei nº 10.861 de 14 de abril de 2004, que institui o Sistema Nacional de Avaliação da Educação Superior - SINAES, com o objetivo de assegurar processo nacional de avaliação das instituições de educação superior, a avaliação visa melhorar a qualidade da educação superior. Para isso, a avaliação institucional dos cursos e do desempenho dos alunos são mecanismos básicos para ponderações acerca da execução dos processos educativos na Educação Superior.

A Comissão Própria de Avaliação - CPA, do IFCE, instituída com base no art. 11 da lei nº 10.861/2004, tem a finalidade de implementar o processo de autoavaliação do Instituto, bem como a sistematização e a prestação das informações solicitadas pela Comissão Nacional de Avaliação da Educação Superior (CONAES). Esta comissão é formada por quatro representantes do corpo docente; quatro representantes do corpo técnico-administrativo; quatro representantes dos alunos; e quatro representantes da sociedade civil organizada, de acordo com o Regimento Interno da CPA do IFCE.

A autoavaliação institucional, tarefa da CPA, compõe-se da autoavaliação dos campi, que se dá em torno das seguintes informações: identificação e histórico do campus, identificação da subcomissão de avaliação e seus trabalhos e os resultados da autoavaliação por campus, com suas respectivas dimensões. 

Os instrumentos utilizados são questionários compostos de perguntas em cada dimensão descrita, aplicados com professores, alunos e técnicos. Os resultados são categorizados pelas respostas às perguntas em fragilidade, avaliação mediana ou potencialidade de cada grupo aplicado. Para conclusão, a classificação final se dá pela fragilidade, potencialidade, controvérsia ou tendência a um dos dois primeiros conceitos. 

A autoavaliação institucional prediz tomadas de decisões acerca do funcionamento do campus em torno de suas atividades, caracterizadas pelas dimensões avaliadas. Também representa a comunidade acadêmica e em geral na participação da melhoria do campus nos seus processos de criação e manutenção de mecanismos e estratégias corroboradoras de uma educação adequada à realidade local e global. 

Já a avaliação do curso, de caráter externo, compreende a avaliação do desempenho do estudante por meio do ENADE e das condições de ensino oferecidas a ele, especificamente o perfil do corpo docente, as instalações físicas e a organização didático-pedagógica. 

A periodicidade máxima de aplicação do ENADE é trienal, e consiste em contextualizar o perfil dos estudantes avaliados, o que fornece maiores subsídios para políticas e tomadas de decisões sobre o percurso do curso e da instituição na oferta do mesmo. O ENADE é componente curricular obrigatório nos cursos de graduação, sendo o mesmo explicitado aos estudantes pelo NDE, Colegiado e CPA. 

A avaliação externa, do SINAES, compreende o referencial básico para os processos de autorização, reconhecimento e renovação de reconhecimento de cursos de graduação. Portanto, o presente projeto vislumbra não somente a regulação do \NOMECURSO, como projeta a sua qualidade e estima diante da comunidade externa na participação de processos avaliativos que ensejem tal qualidade.

As avaliações internas e externas são importantes, pois a gestão do curso é realizada considerando a autoavaliação institucional e o resultado das avaliações externas como insumo para aprimoramento contínuo do planejamento do curso, com evidência da apropriação dos resultados pela comunidade acadêmica e existência de processo de autoavaliação periódica do curso.

Importante destacar que os dados coletados nas avaliações externas e da Comissão Própria de avaliação – CPA serão utilizadas como subsídio à Coordenadoria de Curso, à Coordenadoria Técnico-Pedagógica, à Diretoria de Ensino e aos próprios docentes para intervirem de forma a viabilizar melhorias no processo de ensino-aprendizagem. 

A partir do exame ENADE, por exemplo, é possível avaliar a qualidade do \NOMECURSO ofertado pelo IFCE Campus Maracanaú e o rendimento de seus discentes em relação aos conteúdos programáticos, a suas habilidades e competências. 

A partir das avaliações externas é possível verificar todos os tópicos contidos no Projeto Pedagógico do Curso e confrontar o que está escrito com as práticas docentes, infraestrutura, etc. 

Questões e resultados obtidos por meio da CPA, desempenho dos estudantes no ENADE, reconhecimento do curso e renovação do reconhecimento, serão pautados e discutidos amplamente pelo Colegiado do curso, divulgados à comunidade acadêmica e, sempre que necessário, serão feitos os devidos encaminhamentos ao NDE. 

%Os dados coletados das avaliações docentes servem de subsídio as Coordenadorias de Curso, Coordenadoria Técnico-Pedagógica, Diretoria de Ensino e aos próprios professores para intervirem de forma a viabilizar melhorias no processo de ensino-aprendizagem. Os feedbacks são geralmente realizados pela Coordenadoria Técnico-Pedagógica, quando julgar necessário, por meio de conversas individuais ou em grupo conforme demanda por parte dos estudantes, das Coordenações de Curso, da Direção de Ensino e dos próprios professores.

%A partir do exame ENADE é possível avaliar a qualidade do \NOMECURSO ofertado pelo IFCE - Campus Maracanaú e o rendimento de seus alunos em relação aos conteúdos programáticos, suas habilidades e competências. O resultado desta avaliação serve de base para atualização do projeto pedagógico do curso, principalmente da matriz curricular.

%Finalmente, através do reconhecimento do curso pelo MEC é possível verificar todos os tópicos contidos no projeto pedagógico e confrontar o que está escrito com as práticas docentes, infraestrutura etc. Esta avaliação é importante, pois seu resultado servirá como um dos elementos norteadores para melhoria do curso ofertado.

% AVALIAÇÃO DA APRENDIZAGEM ----------------------------------------------

\clearpage
\section{AVALIAÇÃO DA APRENDIZAGEM}

A avaliação é um processo minucioso que permeia todo o processo ensino-aprendizagem. Avaliar não consiste somente em fazer provas e atribuir notas, mas é um processo pedagógico contínuo, que ocorre dia após dia, buscando corrigir erros e construir novos conhecimentos.

Avaliar consiste em analisar o desempenho do aluno quanto ao domínio das competências previstas face ao perfil necessário à sua formação, através da adoção de vários instrumentos e técnicas de avaliação, que deverão estar diretamente ligados ao contexto da área objeto e utilizados de acordo com a natureza do que está sendo avaliado.

A avaliação da aprendizagem acontece para que o professor tenha noção se os conhecimentos e competências necessárias à formação foram internalizados pelos discentes, bem como também serve para que o docente possa executar uma autoavaliação acerca de sua didática e metodologia de ensino, sendo possível, dessa forma, verificar se o caminho que está percorrendo deve ser revisto. Tratada dessa forma, a avaliação permite diagnosticar a situação do discente, em face da proposta pedagógica da instituição e orientar decisões quanto à condução da prática educativa. Portanto, o seu propósito é subsidiar a prática do professor, oferecendo pistas significativas para a definição e redefinição do trabalho pedagógico.

Como tal, a avaliação deverá ser contínua, processual e cumulativa, considerando a prevalência de aspectos qualitativos sobre os quantitativos, assim como estabelece a Lei de Diretrizes e Base da Educação (n° 9.394/96), isso para que seja efetivada a sua função formativa, servindo para o discente como parâmetro de referência de suas conquistas, dificuldades e possibilidades de crescimento e tendo em vista que o desenvolvimento de competências não envolve apenas conteúdos teóricos, mas, sobretudo práticas e atitudes
Nesse contexto, o processo de avaliação do \NOMECURSO do Campus de Maracanaú é orientado pelos objetivos do curso e perfil profissional do egresso a ser formado, sendo definido no Programa de Unidade Didática de cada disciplina (PUD).

Nesse sentido, vale destacar que os processos, instrumentos, critérios e valores de avaliação adotados pelo professor deverão ser explicitados aos estudantes no início do período letivo, quando da apresentação do PUD, observadas as normas dispostas no Regulamento da Organização Didática (ROD) do instituto, onde estão definidos os critérios para atribuição de notas, as formas de recuperação, promoção e frequência do aluno, assim como na Lei de Diretrizes e Base da Educação Nacional.
 
Considerando que o desenvolvimento de competências envolve conhecimentos, práticas e atitudes, o processo avaliativo exige diversidade de instrumentos e técnicas de avaliação, que deverão estar diretamente ligados ao contexto da área objeto da educação profissional e utilizados de acordo com a natureza do que está sendo avaliado.

Desta forma, são utilizados instrumentos diversificados que possibilitam ao professor observar e intervir no desempenho do aluno considerando os aspectos que necessitem ser melhorados, orientando a este, no percurso do curso, diante das dificuldades de aprendizagem apresentadas, reconhecendo as formas diferenciadas de aprendizagem, em seus diferentes processos, ritmos, lógicas, exercendo, assim, o seu papel de orientador e mediador que reflete na ação e que age sobre a realidade. 

Serão considerados instrumentos de avaliação, os trabalhos de natureza teórico-práticos; observação diária dos estudantes pelos professores, durante a aplicação de suas diversas atividades, exercícios, trabalhos individuais e/ou coletivos; fichas de observações; relatórios; autoavaliação; provas escritas com ou sem consulta; provas práticas e provas orais; seminários; projetos interdisciplinares, resolução de exercícios, planejamento e execução de experimentos ou projetos, relatórios referentes a trabalhos; experimentos ou visitas técnicas; realização de eventos ou atividades abertas à comunidade; autoavaliação descritiva e outros instrumentos de avaliação considerando o seu caráter progressivo e que enfatizem a resolução de situações problema específicas do processo de formação do técnico.

Dentre esses vários instrumentos podemos destacar:
\begin{itemize}
\item Trabalho de pesquisa/projetos para verificar as capacidades de construir objetivos e alcança-los; caracterizar o que vai ser trabalhado; antecipar resultados; escolher estratégias mais adequadas à resolução do problema; executar ações; avaliar essas ações e as condições de execução; seguir critérios preestabelecidos;
\item Observação da resolução de problemas relacionados ao trabalho em situações simuladas ou reais, com o fim de verificar que indicadores demonstram a aquisição de competências mediante os critérios de avaliação previamente estabelecidos;
\item Seminários de exposição de conteúdos ou experiência prática de campo são procedimentos metodológicos importantes porque pressupõem o uso de ferramentas e técnicas para pesquisa, estudo e trabalho em equipe;
\item Análise de casos – os casos são desencadeadores de um processo de pensar, fomentador da dúvida, do levantamento e da comprovação de hipóteses, do pensamento inferencial, do pensamento divergente, entre outros.
\item Prova – visa verificar a capacidade adquirida pelos alunos de aplicar os conteúdos aprendidos. Como, por exemplo: analisar, classificar, comparar, criticar, generalizar e levantar hipóteses, estabelecer relações com base em fatos, fenômenos, ideias e conceitos. Para fins de promoção são avaliados tanto o desempenho quanto a assiduidade do aluno. O aluno será orientado na medida em que os resultados das atividades não sejam apenas comunicados, mas discutidos, indicando erros, identificando dificuldades e limitações, sugerindo possíveis soluções e rumos, considerando o caráter progressivo da avaliação.
A sistemática de avaliação no IFCE se desenvolverá em duas etapas e em cada uma será computada a média obtida pelo discente. Independentemente do número de aulas semanais, o docente deverá aplicar, no mínimo, duas avaliações por etapa. 
\end{itemize}

O estudante que não atingir o mínimo necessário para aprovação, poderá realizar avaliação de recuperação, conforme estabelecido no Regulamento da Organização Didática do IFCE. A sistemática de avaliação no IFCE é apresentada na Subseção I, Seção I, Capítulo III, Título III, do Regulamento da Organização Didática (ROD) de junho de 2015, ANEXO \ref{anexoRODavaliacao}.

% EMISSÃO DE DIPLOMAS ----------------------------------------------

\clearpage
\section{EMISSÃO DE DIPLOMAS}

O diploma de \NOMEEGRESSO será conferido aos egressos do \NOMECURSO . A emissão do diploma  está condicionada à conclusão com aprovação de todos os componentes curriculares obrigatórios que compõem a matriz curricular, observadas a carga horária mínima de componentes curriculares optativos, incluindo o Trabalho de Conclusão de Curso (TCC) e o estágio curricular obrigatório. 

Deve ser observada a regulamentação da certificação estabelecida na legislação vigente e em âmbito Institucional, por intermédio da Pró-reitoria de Ensino, que estabelecerá normas complementares, regulamentando os processos em relação a prazos e procedimentos.

% POLÍTICAS INSTITUCIONAIS NO ÂMBITO DO CURSO ----------------------------

\clearpage
\section{POLÍTICAS INSTITUCIONAIS NO ÂMBITO DO CURSO}

O \NOMECURSO do IFCE é um curso que vem ganhando renome no âmbito do estado, reconhecido pelos órgãos governamentais de ensino e pesquisa, pelo setor empresarial e, enfatize-se, pela sociedade. O curso também é bastante conhecido pelas pessoas que não estão relacionadas diretamente com a educação, mas que possuem consciência da importância dos bons cursos para a sociedade. Esse reconhecimento da sociedade é resultado de diversas políticas adotadas pelos dirigentes desde sua criação, das quais pode-se citar:

\begin{itemize}
\item Contratação de professores em regime de 40 horas e dedicação exclusiva (DE) em meados da década de 2010, fundamentado na crença de que bons cursos só podem ser desenvolvidos com professores inteiramente dedicados a ele e em constante reciclagem;

\item Institucionalização dos CEFETs, o que alavancou a pesquisa e a extensão nos atuais IFs;

\item Política de incentivo a qualificação docente através de cursos a nível de pós-graduação (mestrado, doutorado e pós-doutorado) dos professores com afastamento temporário, podendo os mesmos realizarem seus estudos no Brasil ou no exterior, de forma a prover recursos humanos para a condução de pesquisas de ponta;

\item Incentivo constante para a interação universidade/empresa, com o objetivo de transferir para a sociedade os conhecimentos adquiridos nas atividades de pesquisa bem como direcionar estas mesmas pesquisas à resolução de demandas reais, além estreitar laços que melhor permitam a inserção de nossos alunos no mercado de trabalho;

\item Preocupação em sempre admitir profissionais de qualidade acadêmica em seu quadro.
\end{itemize}
 
O processo de pós-graduação dos professores está com cerca de 70\% concluído, e espera-se que até o final de 2019 esteja com 100\% de professores doutores (ver seção \ref{CORPODOCENTE}).

Os esforços também passam por melhorias e aumento da infraestrutura de ensino e pesquisa. Os laboratórios, (ver seção \ref{INFRA}, atualmente em número de 13), foram criados, ampliados e especializados e a interação com o setor empresarial vem se consolidando. A Tab. \ref{tab11} representa o número de projetos de pesquisa e de extensão aprovados pelos professores do eixo da indústria. 

\begin{longtable}{|p{1.5cm}|p{4.5cm}|p{4.5cm}|}
    \caption{Número de projetos de pesquisa e de extensão}    \label{tab11} 
    \vspace{-0.25cm}
    %\multicolumn{3}{c}{\textbf{Tabela 9. Número de projetos de pesquisa e de extensão}}    \label{tab11} 
    \\ \hline 
    \cellcolor{\cinza} Ano & \cellcolor{\cinza}Projetos de Pesquisa & \cellcolor{\cinza}Projetos de Extensão  \\ 
    \hline
     2013 &  9  &  15 \\ \hline
     2014 &  13  &  34 \\ \hline
     2015 &  07  &  18 \\ \hline
     2016 &  12  &  07 \\ \hline
     2017 &  19  &  03 \\ \hline
     Total &  59  & 77 \\ \hline
   \end{longtable}

O ensino e a pesquisa, como deve a ser a regra, andam lado a lado no \NOMECURSO, propiciando o nascimento de um forte curso, uma vez que seus docentes, em função de suas atividades em pesquisa, sempre procuraram ministrar aos alunos, sólidos fundamentos da área da disciplina. Os laboratórios foram aparelhados com equipamentos diversos. Projetos em sua maioria de própria autoria da Gestão do Campus, mais também provindos de financiamento de empresas, do CNPq, do Finep e outros, forneceram os recursos para a aquisição dos diversos equipamentos hoje disponíveis, para a manutenção desses equipamentos, para a montagem de estruturas computacionais, bem como para a aquisição de material de consumo. Estes equipamentos estão, em sua maioria, listados na seção \ref{INFRA}. Ao mesmo tempo em que o apoio governamental direto, através do Ministério da Educação se restringia aos salários dos professores e a infraestrutura básica, os recursos dos projetos foram permitindo a contratação de pessoal técnico e também o pagamento de bolsas de iniciação científica, em muitos casos. Criou-se, assim, uma grande estrutura de pesquisa e de extensão, com reflexos diretos no ensino de graduação, através da iniciação científica e também pela própria participação do professor em sala de aula. 

A argumentação inicial de que a localização geográfica do IFCE Campus Maracanaú, em uma cidade praticamente provida de indústrias, seria benéfico ao curso, vem sendo comprovada no dia a dia do Campus, pois a influência do centro de pesquisa e de formação aqui instalado, foi muito além da fronteira municipal, ultrapassando a fronteira estadual e até nacional. Tal nacionalização e posterior internacionalização também vem sendo complementada pelo recém-criado Mestrado em Energias Renováveis do IFCE Campus Maracanaú, proposto e aprovado pelos professores do eixo da indústria, aqui mencionados, e do eixo da química e meio ambiente também presente no Campus. Não se tem empreendido um esforço direto, até o presente momento, para quantificar o tamanho desta influência na economia estadual e também nacional; entretanto sabe-se da sua importância nos meios acadêmicos e industriais. 

Assim, a qualidade de nossos alunos vem melhorando ano a ano, principalmente, após a intensificação de participação destes alunos como bolsistas de iniciação científica nos laboratórios. Esta atividade, além de fornecer subsídios de aprendizagem de conteúdos técnicos, extrapolando a sala de aula, e mostrando aplicações para os fundamentos lá vistos, permite ao aluno um aprendizado gradual e não aparente da cultura da pesquisa, do comportamento do engenheiro perante os problemas e acelera seu processo de maturidade, pois em geral ele convive com mestrandos e professores. Vale ressaltar que a inserção dos alunos de graduação nos laboratórios de ensino e pesquisa, permite o seu contato com estas atividades e com as atividades de extensão. Esses alunos são orientados e estão em contato direto com os alunos de mestrado, auxiliando-os a desenvolver os temas de pesquisa destes últimos. Os trabalhos de extensão, também são resolvidos nesses ambientes, com a participação de todos. A discussões que surgem, os experimentos que são montados, os resultados que são obtidos, criam uma boa atmosfera para o crescimento dos estudantes. Além de seu grupo de colegas da mesma fase, o aluno inserido no laboratório passa a fazer parte de outro grupo, dando-lhe mais maturidade, tendo como retorno o melhoramento do seu desempenho nas diversas disciplinas.

Outras atividades incentivadas ao longo do curso e que também motivam os alunos, são: Projetos AeroDesign e Enactus; Empresa Júnior; Centro Acadêmico; organização de eventos técnicos, culturais e esportivos; Cooperação Internacional; viagens de estudos; estágios; etc.

A Tabe. \ref{tab10} apresenta os investimentos financeiros relacionados ao \NOMECURSO desde o ano de 2012. 

    \begin{longtable}{|p{1.0cm}|p{6.0cm}|p{3.0cm}|p{3.0cm}|}
    
    \caption{Investimentos financeiros relacionados ao curso}\label{tab10} 
    \vspace{-0.25cm}
    
    %\multicolumn{4}{c}{\textbf{Tabela 8. Investimentos financeiros relacionados ao curso}}\label{tab10} 
    \\ \hline 
    \cellcolor{\cinza}Ano & \cellcolor{\cinza}Equipamento & \cellcolor{\cinza}Empenho, Pregão, Contrato & \cellcolor{\cinza}Investimento [R\$]  \\ 
    \hline
     2012 &  Bloco da Oficina Mecânica & Nº40/2012  &  3.743.946,00 \\ \hline
     2013 &  Equipamentos para os Laboratórios  & PE 14/2013 e PE 06/2013 &  70.167,70 \\ \hline
     2014 &  Maquinário da Oficina Mecânica & PE 08/2014 &  938.260,88 \\ \hline
     2016 &  Maquinário da Oficina Mecânica  & PE 11/2016 & 135.004,76 \\ \hline
     2016 &  Livros & 2016NE800268, 2016NE800331, 2016NE800269 &  113.700,30 \\ \hline
      &  Total & & 5.001.079,64 \\ \hline
   \end{longtable}

% APOIO AO DISCENTE ---------------------------------------------

\clearpage
\section{APOIO AO DISCENTE}\label{APOIO}

Para o apoio aos discentes do \NOMECURSO, o Campus Maracanaú dispõe, além da coordenação do curso, de outros setores que promovem ações de orientação e acompanhamento pedagógico e psicológico, assim como serviços de assistência social, enfermagem e nutrição e núcleos de inclusão e acessibilidade.
 
O envolvimento, a participação e a colaboração dos setores como Coordenadoria de Assuntos Estudantis (serviço social, psicologia, enfermagem e nutrição), Coordenadoria de Controle Acadêmico, Coordenadoria Técnico-Pedagógica, Estágio, Biblioteca, entre outros, que também lidam com o corpo discente do campus, colaboram com a redução contínua da evasão e da retenção acadêmica, permitindo assim, que o estudante permaneça na instituição e conclua seu curso com êxito. As ações realizadas por cada setor ou serviço estão listadas a seguir:

%\clearpage
\subsection{Coordenação de Curso}
No que tange ao apoio discente, a Coordenação do \NOMECURSO visa ser facilitadora nas ações acadêmicas relacionadas ao curso e na resolução de possíveis demandas pelos alunos. Para tanto, lança mão de ações sistematizadas que vão desde o atendimento aos discentes, à promoção de estratégias para melhoria de processo de aprendizagem. 

Ademais, a coordenação do curso administra a potencialidade do corpo docente do seu curso, favorecendo a integração e a melhoria contínua. A atuação da coordenação do \NOMECURSO é detalhada em um plano de ação semestral ou anual elaborado com base nas diretrizes da Nota Técnica 4/2018/PROEN/REITORIA. O referido plano dispõe de indicadores de desempenho disponíveis e públicos e, de forma geral, esta é feita através de:

\begin{itemize}
\item Resolver assuntos ligados ao aproveitamento de disciplinas e à entrada de alunos no curso por meio do Edital de Transferidos / Graduados;
\item Tratar de assuntos pedagógicos do curso junto a Coordenação Pedagógica;
\item Tratar de questões ligadas à matricula e situação regular dos alunos;
\item Acompanhar a vida acadêmica do corpo discente;
\item Coordenar as atividades relacionadas ao reconhecimento do curso;
\item Solicitar dos professores os Programas de Unidades Didáticas (PUD) de todas as disciplinas do curso, bem como mantê-los atualizados;
\item Atuar juntamente com o NDE e o Colegiado nas discussões sobre alterações na matriz curricular, quando se fizer necessário;
\item Atuar juntamente com o NDE e o Colegiado na atualização do projeto pedagógico do curso, quando necessário;
%\item Coordenar as atividades desenvolvidas pelos monitores e pelos bolsistas de laboratórios;
\item Auxiliar ao setor administrativo em assuntos estratégicos, tais como: levantamento de demandas (infra-estrutura, equipamentos, etc.), elaboração de planos de trabalho, elaboração dos horários semestrais, levantamento de demandas de perfis de vagas para novos professores, planejamentos anuais, dentre outros.
\item Liderar as ações de divulgação do curso na sociedade;
\item Representar o curso em eventos e reuniões internas e externas, quando for o caso.
\item Acompanhar o registro de aulas no acadêmico;
\item Registrar aulas extras no acadêmico em comum acordo entre os professores e os discentes;
\item Atender às solicitações de reabertura de diários no acadêmico em comum acordo entre professores e discentes;
\item Dentre outras atividades detalhadas no plano de ação anteriormente mencionado.
\end{itemize}

\subsection{Coordenadoria de Controle Acadêmico}

A Coordenação de Controle Acadêmico (CCA), como órgão de execução, responde pelas questões operacionais junto ao Sistema Q-Acadêmico. Desse modo, define junto a Diretoria de Ensino/DIREN, a qual é subordinada, a execução dos processos de pré-matrícula, matrícula, criação de turmas e horários.
 
Pelo princípio da legalidade, executa procedimento em acordo com o Regulamento da Organização Didática/ROD, o que possibilita auxiliar coordenadores e estudantes quanto às diretrizes estabelecidos no regulamento, além de gerenciar procedimento de ingresso através do Sistema de Seleção Unificada – Sisu, através do acesso ao SISUGestão, o setor também controla e organiza arquivos de discentes. 

No atendimento ao público discente, emite documentação de situação acadêmica, como históricos, declarações e ementas das disciplinas aprovadas.

\subsection{Departamento de Extensão, Pesquisa, Pós-Graduação e Inovação}

%A Coordenadoria de Pesquisa, Pós-­Graduação e Inovação é um setor diretamente subordinado ao departamento de Extensão, Pesquisa, Pós-Graduação e Inovação – DEPPI, responsável por atividades de atendimento ao discente no que tange à concessão de auxílio acadêmico, auxílio didático-­pedagógico, bem como suporte aos discentes que participam de projetos de pesquisa como bolsistas com fomento ou bolsistas voluntários. Os auxílios são destinados aos alunos que participam de eventos científicos e/ou tecnológicos de âmbito nacional ou internacional.
%Em relação ao fomento da pesquisa, a coordenadoria divulga informações sobre editais internos e externos, além de dar orientação sobre os procedimentos de cadastro de bolsistas e acompanhar o andamento das atividades de pesquisa. Ademais, os discentes participantes do programa de Iniciação Científica recebem apoio para divulgar os resultados de suas pesquisas através do Seminário de Iniciação Científica e Tecnológica (SEMIC).

O Departamento de Extensão, Pesquisa, Pós-Graduação e Inovação (DEPPI) do  Campus Maracanaú é responsável por atividades de atendimento ao discente no que tange a concessão de auxílio acadêmico, auxílio didático-pedagógico, bem como suporte aos discentes que participam de projetos de extensão, pesquisa e inovação como bolsistas com fomento ou sem fomento por meio do auxílio acadêmico. O departamento apoia os discentes que aprovam trabalhos em eventos científicos e/ou tecnológicos de âmbito nacional ou internacional. 

Em relação ao fomento da extensão, pesquisa e inovação, o DEPPI divulga informações sobre editais internos e externos, além de dar orientação sobre os procedimentos de cadastro de bolsistas e acompanhar o andamento das atividades de pesquisa e dos projetos de extensão. 

Esse departamento, ainda, realiza captação de vagas de estágio por meio de seu programa de Relações Empresariais, no qual mantém estreita relações com as empresas locais e trabalha em conjunto com o Setor de Estágios e a Diretoria de Ensino (DIREN). 

Ademais, o DEPPI organiza eventos de extensão e pesquisa como a Semana de Integração Científica (SIC) e o Seminário de Iniciação Científica e Tecnológica (SEMIC). Nesses eventos, os discentes participantes dos programas de Iniciação Científica e de Extensão recebem apoio para divulgar os resultados de seus trabalhos para a comunidade.

\subsection{Coordenadoria Técnico Pedagógica}

A Coordenadoria Técnico-Pedagógica (CTP) do IFCE é o setor responsável pelo planejamento, acompanhamento, avaliação de ações pedagógicas desenvolvidas no campus com vistas à formulação e reformulação contínua de intervenções pedagógicas que favoreçam o alcance de resultados satisfatórios quanto ao processo ensino-aprendizagem.

A atuação da CTP é embasada nos fundamentos e pressupostos teóricos educacionais, nos princípios legais da Educação Brasileira. A atuação desse setor encontra-se em consonância com a Missão Institucional do IFCE. A seguir, apresentam-se as atribuições do referido setor:

\begin{itemize}
\item Realizar atendimento individual e/ou em grupo aos estudantes, professores, pais e ou responsáveis conforme necessidades observadas pelo setor e ou quando solicitado;
\item Analisar continuamente as causas da evasão e repetência para formular sistematicamente estratégias que visem à superação ou minimização do problema;
\item Acompanhar o desenvolvimento dos estudantes com baixo rendimento escolar (frequência e desempenho) propondo alternativas que favoreçam a superação e a minimização dessa problemática;
\item Mediar à relação professor e aluno e aluno - aluno voltada para o sucesso do desempenho discente solicitando apoio e parceria da Assistência Estudantil e Setor de Psicologia, quando necessário;
\item Realizar a cada período letivo, a Avaliação de Desempenho Docente, com vistas a promover a melhoria da prática docente por meio de análise dos resultados da avaliação e dos feedbacks que serão dados por meio de conversas individuais e capacitações;
\item Promover ações formativas (encontros pedagógicos, encontros de estudo, capacitações, orientações individuais, conselhos de classe, colegiados, entre outros) que provoquem no professor avaliação de sua prática docente para que ao longo do processo melhorem sua atuação pedagógica.
\end{itemize}

Convém destacar que as atribuições da CTP se articulam com as ações desenvolvidas por outros setores da instituição, como coordenações de cursos, coordenadoria de assistência estudantil (serviço social, psicologia, enfermagem e nutrição), coordenadoria de controle acadêmico, estágio, biblioteca, pesquisa, extensão, entre outros, que também lidam com o corpo discente do campus.

\subsection{Coordenadoria de Assuntos Estudantis}

A Assistência estudantil vem se consolidando no Instituto Federal de Educação, Ciência e Tecnologia do Ceará – IFCE como um conjunto de ações, configurando-se através de auxílios financeiros e serviços, visando ampliar as condições de permanência e apoio à formação acadêmica do corpo discente. Uma dessas ações diz respeito à disponibilização de serviços, caracterizados por ações continuadas, visando ao atendimento biopsicosocial do discente. Outra ação diz respeito aos auxílios sob a forma de pecúnia, sendo estes destinados, na sua maioria, ao discente, prioritariamente em condições de vulnerabilidade social, e operacionalizados por meio do regulamento dos auxílios. Tal regulamento é normatizado pelo programa de Auxílios, previsto na Política de Assistência Estudantil do IFCE ( aprovada pela resolução nº 024, de 22 de julho de 2015) e, institui ações de efetivação do Decreto nº 7.234, de 19 de junho de 2010, que dispõe sobre o Programa Nacional de Assistência Estudantil (PNAES).
 
As ações previstas na PNAES dizem respeito às seguintes áreas: moradia estudantil, alimentação, transporte, atenção à saúde, inclusão digital, cultura, esporte, creche, apoio pedagógico, acesso e participação e aprendizagem de estudantes com deficiência, transtornos globais do desenvolvimento e altas habilidades e superdotação (Decreto 7.234/2010, Art. 3º). Ressaltamos, ainda, que o referido decreto prevê que estas ações serão executadas por Instituições Federais de Ensino Superior, contemplando os IFs. 
Portanto, a assistência Estudantil no IFCE, vislumbrada mediante serviços ofertados (merenda escolar, atendimento psicológico, atendimento pedagógico, entre outras ações) e auxílios financeiros foram instituídos na perspectiva de “viabilizar a igualdade de oportunidades, contribuir para a melhoria do desempenho acadêmico e agir, preventivamente, nas situações de retenção e evasão decorrentes da insuficiência de condições financeiras.
  
O IFCE Campus Maracanaú dispões dos seguintes serviços, diretamente subordinados à Coordenadoria de Assuntos Estudantis, a saber: Serviço de Enfermagem, Serviço de Nutrição, Serviço de Psicologia e Serviço Social. As ações realizadas por cada serviço estão listas a seguir:

\subsubsection{Serviço de Enfermagem} 

No âmbito do IFCE, a Enfermagem destina-se a promoção da saúde com foco na educação em saúde, bem como a oferecer cuidados de primeiros socorros em situações de urgência e emergência, conforme ações elencadas a seguir: 

\begin{itemize}
\item contribuir para o desenvolvimento integral do(a) discente; 
\item colaborar no mapeamento da realidade socioeconômica, acadêmica e de saúde dos discentes; 
\item apoiar as estratégias de inclusão das pessoas com deficiência;
\item atuar na prevenção, promoção, tratamento e vigilância à saúde de forma individual e coletiva, colaborando com o processo de ensino-aprendizagem; 
\item realizar ações de prevenção e controle sistemático de situações de saúde e agravos em geral; 
\item desenvolver atividades de educação em saúde para a adoção de hábitos saudáveis, visando à melhoria da qualidade de vida e à promoção da saúde da comunidade acadêmica; 
\item participar de estratégias de combate à evasão escolar; 
\item participar do planejamento, execução e avaliação da programação das ações anuais de saúde; 
\item participar do processo de seleção de auxílios referente aos aspectos relativos às situações de saúde;
\item acompanhamento de discentes aos serviços de saúde, nas situações previstas nas diretrizes para atuação do enfermeiro no IFCE;
\item programa Saúde e Prevenção nas Escolas(SPE);
\item camisinha Card;
\item vacinação, sendo previstas a realização de 06 campanhas anuais;
\item campanhas semestrais de doação de sangue;
\item atendimento ambulatorial: aconselhamento em DST/HIV e AIDS; realização de curativos; aferição de pressão arterial; glicemia; auscultas cardíacas e sinais vitais; primeiros socorros; atendimento aos servidores e alunos com hipertensão e
diabetes.
\end{itemize}

\subsubsection{Serviço de Nutrição} 

O Serviço de Alimentação e Nutrição é responsável pela administração da Unidade de Alimentação no campus, a qual visa à oferta de uma alimentação adequada, compreendendo o uso de alimentos variados, seguros, que respeitem a cultura, as tradições e os hábitos alimentares saudáveis, contribuindo, assim, para melhoria do rendimento escolar, permanência do estudante no espaço educacional e promoção de hábitos alimentares saudáveis.

O Serviço de Nutrição ainda atua nos programas de educação e assistência nutricional, desenvolvendo ações com a equipe multiprofissional tendo em vista a promoção da saúde e segurança alimentar e nutricional, prestando, também, assessoria às atividades de ensino, pesquisa e extensão.


\subsubsection{Serviço de Psicologia} 

A psicologia escolar/educacional assume um papel de contribuir para a construção de uma educação de qualidade, baseada nos princípios do compromisso social, do respeito à diversidade e dos direitos humanos. Entende que a ação educativa é permeada por determinantes biopsicossociais que interferem, direta e indiretamente, no desenvolvimento do processo de aprendizagem de cada individuo, desse modo a ação educativa não se limitará a queixa, mas a busca constante de fomentar um ambiente escolar que promova saúde mental.

Neste sentido, o serviço de Psicologia do IFCE - Campus Maracanaú busca:

\begin{itemize}
\item Apoiar servidores no trabalho com a heterogeneidade de discentes;
\item Avaliar, acompanhar e orientar dentro do contexto institucional casos que requeiram encaminhamentos clínicos, estabelecendo um espaço de acolhimento, escuta e reflexão. No caso de demandas psicoterápicas, será realizado encaminhamento para outras instituições que ofereçam o tratamento adequado;
\item Fazer parte da equipe multiprofissional que envolve o processo de ensino e aprendizagem levando em conta o desenvolvimento global do discente;
\item Propiciar condições para que o discente expresse sua autonomia e consciência crítica, por meio da participação ativa na vida acadêmica, contribuindo para uma formação cidadã; 
\item Realizar acompanhamento dos discentes em situação de vulnerabilidade socioeconômica e dificuldade de aprendizagem para a realização das intervenções necessárias;
\item Identificar e analisar as causas e as motivações das reprovações, retenções e evasões dos discentes, a fim de subsidiar o direcionamento das intervenções, apreendendo quais os aspectos sociais, físicos, cognitivos e afetivos geram resistência no seu processo de aprendizagem elaborando condições para permanência da qualidade da aprendizagem;
\item Propiciar aos discentes espaços de reflexão e diálogo sobre as temáticas demandadas pelos diversos atores que compõem a comunidade acadêmica; 
\item Fomentar momentos de expressões artísticas, espirituais, culturais e esportivas do discente e comunidade acadêmica, propiciando as inter-relações e a circulação da palavra nas suas mais diferentes manifestações;
\item Estimular a criatividade e iniciativa dos discentes para criação de grupos autogeridos que trabalhem temáticas por eles definidas;
\item Favorecer a prevenção e promoção da saúde dos discentes e comunidade acadêmica, visando o alcance da discussão dos diversos aspectos que compõem o conceito ampliado de saúde, a partir de trabalhos preventivos que visem um processo de transformação pessoal e social;
\item Promover ações articuladas com a rede socioassistencial, educacional e de saúde do município, inserindo o campus Maracanaú como um dos pontos estratégicos de mobilização social do município.
\end{itemize}


\subsubsection{Serviço Social}

O Serviço Social no Campus de Maracanaú insere-se na promoção do Programa Nacional de Assistência Estudantil (PNAES – Decreto MEC Nº 7234), mediante elaboração e implementação de serviços, programas, projetos e auxílios (sob a forma de pecúnia), visando à ampliação das condições de acesso e de permanência, com enfoque numa formação crítica e autônoma.

A atuação do Serviço Social no Campus situa-se no âmbito da Assistência Estudantil, com destaque nas seguintes ações:
 
\begin{itemize}
\item \textbf{De caráter individual:} atendimento social, escuta qualificada, estudo social, análise socioeconômica, socialização de informações, orientações sociais, encaminhamento para outros serviços, seleção de estudantes para concessão de auxílios.
\item \textbf{De caráter coletivo:} atendimento coletivo, formação de grupos, reuniões, encontros, seminários, oficinas para alunos e técnicos, campanhas, realização de atividades de acolhimento e integração dos discentes à comunidade acadêmica, confecção de materiais educativos, mobilização e organização social e política, apoio à constituição das entidades estudantis, capacitação dos alunos e técnicos, participação nos espaços de controle social.
\end{itemize}

Destacamos que é de responsabilidade do Serviço Social, a concessão dos auxílios financeiros, a saber: 

\begin{itemize}
\item AUXÍLIO MORADIA - subsidia despesas com habitação para locação, sublocação de imóveis para discentes com referência familiar e residência domiciliar fora da Sede do município onde está instalado o campus;
\item AUXÍLIO ALIMENTAÇÃO - subsidia despesas de alimentação nos dias letivos;
\item AUXÍLIO TRANSPORTE – subsidia despesas no trajeto residência/campus/residência;
\item AUXÍLIO ÓCULOS – complementa despesas de aquisição de óculos ou lentes corretivas de deficiências oculares;
\item AUXÍLIO VISITAS/VIAGENS TÉCNICAS – subsidia despesas com alimentação e/ou hospedagem, em visitas e viagens técnicas;
\item AUXÍLIO ACADÊMICO – complementa despesas com alimentação, hospedagem, passagem e inscrição dos discentes para a participação em eventos acadêmicos;
\item AUXÍLIO DIDÁTICO-PEDAGÓGICO – subsidia a aquisição de material de uso individual e intransferível, indispensável à aprendizagem de determinada disciplina;
\item AUXÍLIO DISCENTES MÃES/PAIS – subsidia despesas de filho(s) de até 06 (seis) anos de idade ou com deficiência, sob sua guarda;
\item AUXÍLIO FORMAÇÃO – subsidia despesas relativas à ampliação da formação dos discentes em laboratórios/oficinas e em projetos caracterizados por ensino, pesquisa e extensão, vinculados ao seu curso.
\end{itemize}

Os auxílios têm por objetivos e finalidades ampliar as condições de permanência e apoio à formação acadêmica dos discentes, visando a reduzir os efeitos das desigualdades sociais; contribuir para reduzir a evasão; propiciar a melhoria do desenvolvimento acadêmico e biopsicossocial do discente.

\subsection{Biblioteca}

A Biblioteca Rachel de Queiroz, fundada em 12 de março de 2010, ocupa uma área de 390,11 m2, possui 62 assentos para estudo individual ou em grupo e oferece suporte para o ensino, pesquisa e extensão. Conta com um acervo de 1.850 títulos e 11.900 exemplares entre livros, cd-rom, disquete e folhetos, nas áreas técnicas, tecnólogos, licenciatura e engenharia e tecnologia, com ênfase em livros técnicos e didáticos. Conta ainda com 88 títulos e 407 exemplares de revistas doadas, além de acesso ao portal de Periódicos da CAPES\footnote{www.periodicos.capes.gov.br}, disponível para todos os computadores do IFCE Maracanaú.

Atualmente a Biblioteca Rachel de Queiroz disponibiliza a Biblioteca Virtual Universitária \footnote{https://ifce.bv3.digitalpages.com.br/login} contemplando todo o acervo bibliográfico da lista de títulos da Pearson e das editoras parceiras: Manole, Contexto, Ibpex, Papirus, Casa do Psicólogo, Ática, Scipione, MartinsFontes, Cia das Letras, Rideel, Educs e Jaypee. Todo acervo destas editoras pode ser acessado, de forma digital, por professores e alunos do IFCE. Dentre as muitas áreas beneficiadas destacam-se Ciências da Computação, Engenharias, Física, Matemática e Estatística. Além disso, a biblioteca presta serviços como o empréstimo domiciliar de todos os materiais que compõem o acervo; a consulta à base de dados tanto nos terminais de autoatendimento local quanto via internet; o acesso ao Portal de Periódicos Eletrônicos da Capes; a elaboração de catalogação na fonte; a orientação técnica para elaboração e apresentação de trabalhos acadêmicos, com base nas normas técnicas de documentação da ABNT; e levantamentos bibliográficos e referenciais para pesquisas.

A biblioteca conta com profissionais que registram e catalogam, classificam e indexam as novas aquisições e fazem a manutenção das informações bibliográficas no Sistema Sophia além de realizar preparação física (carimbos de identificação e registro, colocação de etiquetas) do material bibliográfico para empréstimo domiciliar.

A Biblioteca funciona de segunda a sexta-feira, nos seguintes horários:

\begin{itemize}
\item Empréstimo, estudo e leitura - 8h às 20h
\item Referência (Consulta Local) e Periódicos - 8h às 20h
\item Multimídia (computadores conectados à Internet) - 8h às 20h.
\end{itemize}

%
%A biblioteca possui um acervo com mais de 3.700 volumes, entre livros e CD’s, e mais 250 revistas periódicos nas áreas
%de ciências humanas, ciências puras, artes, literatura e tecnologia, com ênfase em livros técnicos e didáticos. São 904
%títulos nas mais diversas áreas do conhecimento.
%
%A biblioteca conta com profissionais que registram e catalogam, classificam e indexam as novas aquisições e fazem a
%manutenção das informações bibliográficas Sistema GNUTECA. Realizam, também, a preparação física (carimbos de
%identificação e registro, colocação de etiquetas, bolso e fichas de empréstimo) do material bibliográfico para
%empréstimo domiciliar.
%
%Principais serviços:
%
%\begin{itemize}
%\item Acesso à Base de Dados Gnuteca nos terminais locais e via Internet.
%\item Empréstimo domiciliar e renovação das obras e outros materiais.
%\item Consulta local ao acervo.
%\item Elaboração de catalogação na fonte.
%\item Orientação técnica para elaboração e apresentação de trabalhos acadêmicos, com base nas Normas Técnicas de
%Documentação da ABNT.
%\item Acesso ao Portal de Periódicos da Capes.
%\item Acesso à Internet.
%\item Levantamento bibliográfico.
%\end{itemize}
%O \textit{campus} conta ainda com um serviço adicional da Biblioteca Virtual, acessível via internet para todos os alunos
%matriculados. Este serviço disponibiliza um grande número de títulos para acesso por computador, dentro ou fora da
%insituição.

\subsubsection{Empréstimo}
O usuário poderá retirar, por empréstimo domiciliar, qualquer publicação constante do acervo bibliográfico, exceto as
obras de referência (enciclopédias, dicionários, atlas, periódicos, jornais, etc) ou outras publicações que, a critério
da Biblioteca, não podem sair.

\subsubsection{Setor Multimídia}
O Setor de Multimídia possui 5 computadores conectados à Internet para que o usuário possa fazer pesquisas online.
Funciona de segunda a sexta-feira, das 8h às 20h.

\subsection{Núcleo de Acessibilidade às Pessoas com Necessidades Educacionais Específicas} 

O Núcleo de Acessibilidade às Pessoas com Necessidades Educacionais Específicas (NAPNE) tem como objetivo disseminar uma cultura da “educação para convivência”, aceitação da diversidade e, principalmente, buscar a quebra das barreiras arquitetônicas, tecnológicas, educacionais e atitudinais.

Para tanto o NAPNE atua no sentido de:

\begin{itemize}
\item Buscar a quebra de barreiras arquitetônicas, comunicacionais, educacionais e atitudinais na Instituição de ensino, por meio de levantamentos e aplicação de questionários periodicamente;
\item Promover condições necessárias para o ingresso, a permanência e o êxito educacional de discentes com necessidades educacionais específicas no IFCE, realizando o acompanhamento dos estudantes;
\item Atuar junto às coordenações de cursos, à equipe pedagógica e aos colegiados dos cursos oferecendo suporte no processo de ensino-aprendizagem dos estudantes com necessidades educacionais específicas, colaborando com a adaptação dos referenciais teórico-metodológicos, colocando a equipe à disposição para prestar esclarecimentos e orientações;
\item Articular junto ao Campus e à PROEXT a disponibilização de recursos específicos para aquisições de materiais de consumo e permanente que possibilitem a promoção das atividades de ensino, pesquisa e extensão com qualidade;
\item Potencializar o processo ensino-aprendizagem por meio da utilização de novas tecnologias de informação e de comunicação (TICs) que facilitem esse processo, por meio da indicação dos recursos já existentes, assim como colaborando com projetos e pesquisas, e ainda promovendo campanha de conscientização e incentivo a ações inclusivas (Prêmio IFCE Inclusivo – premiação de honra ao mérito por ações, projetos e produtos desenvolvidos no IFCE Maracanaú);
\item Promover e participar de estudos, eventos e debates sobre Educação Inclusiva com o intuito de informar e sensibilizar a comunidade acadêmica no âmbito do IFCE e de outras instituições, realizando palestras e rodas de debates (Projeto Encontros Inclusivos), além do curso de Libras (Módulos I, II e III, totalizando 120hs);
\item Contribuir para a inserção da pessoa com necessidades educacionais específicas no IFCE e em espaços sociais, realizando a divulgação dos editais de seleção e dos cursos em instituições que atuem com pessoas com deficiência, além de fazer parceria com o Conselho Municipal dos Direitos da Pessoas com Deficiência de Maracanaú e Associações aproximando-os do campus.
\item Assessorar a Diretoria de Ingressos do IFCE especificamente nos casos de ingresso de estudantes e servidores com necessidades específicas, formando uma comissão para o acompanhamento da análise dos documentos dos cotistas no processo de matrícula.
\item Assessorar, quando necessário, no processo de alterações nas regulamentações que visem o ingresso e a permanência de pessoas com necessidades educacionais específicas no IFCE.
\end{itemize}

\subsection{Setor de Estágio} 

%O Setor de Estágio do campus de Maracanaú do Instituto Federal do Ceará é diretamente subordinado ao departamento de Extensão, Pesquisa, Pós-Graduação e Inovação – DEPPI e responsável pela administração do estágio discente, seja ele obrigatório ou não-obrigatório. Atua, em parceria, com a direção de ensino e coordenações de cursos, e conta com o apoio dos professores orientadores de estágio.
%Ainda, realiza o controle das documentações, acompanhamento dos relatórios e o cumprimento das regras de estágio conforme Lei n° 11.788, de 25 de Setembro de 2008.
%Também, faz a divulgação das ofertas de estágio pelas empresas para disseminar as oportunidades ao corpo discente.

O Setor de Estágio do IFCE-Campus de Maracanaú é diretamente subordinado à Diretoria de Ensino e é responsável pela administração do estágio discente, seja ele obrigatório ou não-obrigatório. Atua, em parceria, com o Departamento de Extensão, Pesquisa, Pós-Graduação e Inovação (DEPPI) e coordenações de cursos, e conta com o apoio dos docentes orientadores de estágio.

Ainda, realiza o controle das documentações, o acompanhamento dos relatórios e cumprimento das regras de estágio conforme Lei N° 11.788, de 25 de setembro de 2008, bem como a divulgação das ofertas de estágio pelas empresas para disseminar as oportunidades ao corpo discente.

\subsection{Setor de Educação Física e Esporte}

O Setor de Educação Física e Esporte (SEFE) oferece a toda a comunidade acadêmica do Campus Maracanaú além de uma avaliação física sistemática, diversas possibilidades para a prática de atividade física e esportes, entre elas: musculação, natação, hidroginástica, treinamento funcional, futebol de campo, futebol de salão, voleibol de quadra, voleibol de areia, futevôlei, basquetebol, handebol, tênis de mesa e jogos de tabuleiro.

O SEFE ainda possibilita ao público discente compor suas seleções esportivas e participar das competições a nível regional (jogos do IFCE sub-19 e aberto) e nacional (jogos dos IF sub-19). Além disso, possibilita também a socialização e integração entre discentes, docentes e comunidade por meio dos projetos de extensão desenvolvidos no setor.

% CORPO DOCENTE -----------------------------------------------

\clearpage
\section{CORPO DOCENTE}\label{CORPODOCENTE}

O detalhamento do corpo docente é dividido em três itens:
\begin{itemize}
\item Corpo docente necessário para desenvolvimento do curso;
\item Corpo docente existente;
\item Dados complementares dos docentes do eixo da indústria.
\end{itemize}

Os quais são apresentados a seguir.

\subsection{Corpo Docente Necessário ao Curso}

Na Tab.\ref{tab1} são apresentadas as áreas, subáreas e quantidade necessária de docentes para o atendimento a todas as disciplinas do curso conforme Tabela de Perfil Docente, disponível na \citeonline{perfildocente}.

\if\COMPILAR0
    \begin{longtable}{|p{3.0cm}|p{3.0cm}|p{1.0cm}|p{6.8cm}|}
    \caption{Quantitativo de Docentes Necessários ao Curso}\label{tab1} 
    \vspace{-0.25cm}
    %\multicolumn{4}{c}{\textbf{Tabela 8. Quantitativo de Docentes Necessários ao Curso}}    \label{tab1} 
    \\ \hline 
    \cellcolor{\cinza}ÀREA & \cellcolor{\cinza}SUBÁREA & \cellcolor{\cinza}QTE & \cellcolor{\cinza}DISCIPLINAS \\ 
    \hline
    Matemática   &  Matemática Básica  &  1  &  Cálculo I, Cálculo II, Cálculo III, Álgebra Linear, Probabilidade e Estatística, Métodos Numéricos, Equações Diferenciais Ordinárias (opt) \\ \hline
    Física   &  Física Geral e Experimental  &  1  &  Física I, Física II, Física III \\ \hline
    Letras   &  Lingua Inglesa  &  1  &  Inglês Técnico, Metodologia Científica \\ \hline
    Educação Física &  Ciências da Saúde  &  1  &  Educação Física \\ \hline
    Engenharia de Materiais e Metalurgia  &  Metarlugia de Transformação  &  1  &  Química Geral, Materiais I, Metrologia, Resistência dos Materiais I \\ \hline
    Engenharia Elétrica   &  Sistemas e Redes de Telecomunicações  &  3  &  Lógica de Programação, Linguagem de Programação, Processamento Digital de Sinais, Sistemas Digitais de Controle Distribuído, Modelagem de Sistemas a Eventos Discretos (opt), Identificação de Sistemas (opt), Sistemas Embarcados, Processamento Digital de Imagens (opt), Inteligência Computacional Aplicada (opt) \\ \hline
    Engenharia Elétrica   &  Circuitos Elétricos  &  2  &  Eletrônica I, Circuitos Elétricos I, Eletrônica II, Circuitos Elétricos II, Instrumentação, Sistemas Lineares, Microcontroladores, Eletrônica III, Controle I, Controladores Lógicos Programáveis, Máquinas Elétricas, Controle II, Dispositivos Periféricos, Robótica I, Robótica II (opt), Controle III (opt) \\ \hline
    Engenharia Elétrica   &  Eletromagnetismo  &  2  &  Desenho Técnico, Desenho Auxiliado por Computador, Instalações Elétricas, Máquinas Elétricas, Acionamentos Hidráulicos e Pneumáticos, Acionamento de Máquinas, Manufatura Auxiliada por Computador, Manufatura Integrada por Computador (opt) \\ \hline
    Engenharia Mecânica   &  Fenômenos de Transporte  &  1  &  Mecânica dos Fluidos (opt) \\ \hline
    Engenharia de Produção   &  Gerência de Produção  &  1  &  Introdução à Engenharia, HST, Ética e Cidadania, Empreendedorismo, Projetos Sociais, Gestão e Controle da Qualidade, Gestão da Manutenção Industrial \\ \hline
   \end{longtable}
\else
    \begin{longtable}{|p{3.0cm}|p{3.0cm}|p{1.0cm}|p{6.8cm}|}
    \multicolumn{4}{c}{\textbf{Tabela 8. Quantitativo de Docentes Necessários ao Curso}}
    \label{tab1} \\ \hline 
    \cellcolor{\cinza}ÀREA & \cellcolor{\cinza}SUBÁREA & \cellcolor{\cinza}QTE & \cellcolor{\cinza}DISCIPLINAS \\ 
    \hline
    Matemática   &  Matemática Básica  &  1  &  Cálculo I, Cálculo II, Cálculo III, Álgebra Linear, Probabilidade e Estatística, Métodos Numéricos, Equações Diferenciais Ordinárias (opt) \\ \hline
    Física   &  Física Geral e Experimental  &  1  &  Física I, Física II, Física III \\ \hline
    Letras   &  Lingua Inglesa  &  1  &  Inglês Técnico, Metodologia Científica \\ \hline
    Educação Física &  Ciências da Saúde  &  1  &  Educação Física \\ \hline
    Engenharia de Materiais e Metalurgia  &  Metarlugia de Transformação  &  1  &  Química Geral, Materiais I, Materiais II, Análise de Falhas, Corrosão e Proteção Anti-Corrosiva (opt), Tribologia e Lubrificação (opt) \\ \hline
    Engenharia Mecânica   &  Projeto de Máquinas  &  2  &  Desenho Técnico, Desenho Auxiliado por Computador, Metrologia, Resistência dos Materiais I, Resistência dos Materiais II, Mecânismos, Sistemas Mecânicos, Transportadores Industriais (opt), Projeto de Engenharia (opt) \\ \hline
    Engenharia Mecânica   &  Fenômenos de Transporte  &  1  &  Termodinâmica, Mecânica dos Fluidos, Transferência de Calor, Máquinas Térmicas e de Fluxo, Refrigeração Industrial \\ \hline
    Engenharia Mecânica   &  Processos de Fabricação  &  1  &  Processos de Fabricação I, Processos de Fabricação II, Soldagem, Manutenção Industrial, Manufatura Auxiliada por Computador, Manufatura Integrada por Computador (opt) \\ \hline
    Engenharia Elétrica   &  Circuitos Elétricos  &  1  &  Lógica de Programação, Linguagem de Programação, Eletrônica I, Circuitos Elétricos I, Instalações Elétricas, Máquinas Elétricas, Acionamentos Hidráulicos e Pneumáticos, Processamento Digital de Imagens (opt), Inteligência Computacional Aplicada (opt) \\ \hline
    Engenharia de Produção   &  Gerência de Produção  &  1  &  Introdução à Engenharia, HST, Ética e Cidadania, Empreendedorismo, Planejamento e Controle da Produção, Projetos Sociais, Gestão e Controle da Qualidade, Gestão da Manutenção Industrial, Inspeção Industrial, Gestão de Projetos (opt) \\ \hline
   \end{longtable}
\fi

\subsection{Corpo Docente Existente}

Na Tab.\ref{tab2} são apresentadas as principais informações do corpo docente existente no campus, conforme às disciplinas ministradas.

\if\COMPILAR0
    \begin{longtable}{|p{3.0cm}|p{2.5cm}|p{2.0cm}|p{6.3cm}|}
    \caption{Corpo Docente Existente no Campus}\label{tab2} 
    \vspace{-0.25cm}
    %\multicolumn{4}{c}{\textbf{Tabela 9. Corpo Docente Existente no Campus}}    \label{tab2} 
    \\ \hline 
    \cellcolor{\cinza}NOME & \cellcolor{\cinza}QUAL.PROF. / REG.TRAB. & \cellcolor{\cinza}TÍTULO & \cellcolor{\cinza}DISCIPLINAS MINISTRADAS \\ 
    \hline
    Adriano Holanda Pereira & Professor / DE  &  Mestrado Eng. Elétrica  & Sistemas Lineares, Circuitos Elétricos II, Instalações Elétricas, Máquinas Elétricas \\ \hline    
    Adriano Barros Carneiro &  Professor / DE  &  Mestrado Desporto  &  Educação Física \\ \hline
    Antônio Barbosa de Souza Júnior & Professor / DE  &  Doutorado Eng. Elétrica. &  Linguagem de Programação, Controladores Lógicos Programáveis (CLP), Robótica I, Trabalho de Conclusão de Curso (TCC) \\ \hline
    Antônio Carlos de Souza & Professor / DE  &  Especialização em Ensino de Física &  Física III \\ \hline
    Antônio Olívio Silveira Britto Júnior & Professor / DE  &  Mestrado Eng. Civil  & Ética e Cidadania, Projetos Sociais  \\ \hline
    Carlos Henrique Sales Martins  &  Professor / DE  &  Mestrado Matemática &  Cálculo II \\ \hline
    %Celso Rogério Schmidlin Júnior & Professor / DE  &  Mestre Eng. Elétrica  & Introdução à Engenharia, Circuitos Elétricos I, Circuitos Elétricos II \\ \hline
    Daniel Barbosa de Brito  &  Professor / DE  &  Doutorado Eng. Teleinf. &  Sistemas Lineares, Processamento Digital de Sinais \\ \hline
    David Carneiro de Souza   &  Professor / DE  &  Doutorado Matemática &  Cálculo I, Cálculo III, Álgebra Linear \\ \hline
    Eugênio Barreto Sousa e Silva & Professor / DE  &  Mestrado Administração &  Empreendedorismo \\ \hline
    Eurípedes Carvalho da Silva & Professor / DE  &  Doutorado Matemática &  Cálculo I, Cálculo II, Cálculo III, Equações Diferenciais Ordinárias (opt.) \\ \hline
    Fábio Timbó Brito & Professor / DE  &  Mestrado Eng. Elétrica  & Microcontroladores, Instrumentação, Dispositivos Periféricos, SDCD \\ \hline
    Fabrício Bandeira da Silva  &  Professor / DE & Doutorado Educação  &  Eletrônica I, Eletrônica II, CLP, Circuitos Elétricos I \\ \hline
    Francisca Ione Chaves   &  Professora / DE & Mestrado em Administração   &  Empreendedorismo \\ \hline
    Francisco Frederico dos Santos Matos   &  Professor / DE  &  Doutorado Eng. Mecânica  &  Mecânica dos Fluidos (opt) \\ \hline
    Francisco José dos Santos Oliveira   &  Professor / DE  &  Mestrado Eng. Materiais  &  Metrologia, Materiais I  \\ \hline
    Francisco Nélio Costa Freitas  &  Professor / DE  & Doutorado Eng. Materiais  &  Química Geral, Metrologia, Resistência dos Materiais I \\ \hline
    Jonatha Rodrigues da Costa  &  Professor / DE  & Mestrado Eng. Teleinf.  &  Lógica de Programação, Métodos Numéricos, Acionamentos Hidráulicos e Pneumáticos  \\ \hline
    José Ciro dos Santos & Professor / DE  &  Mestrado Energias Renováveis &  Desenho Auxiliado por Computador \\ \hline
    José Daniel de Alencar Santos & Professor / DE  &  Doutorado Eng. Teleinf. &  Circuitos Elétricos I, Circuitos Elétricos II, Controle I, Identificação de Sistemas (opt), Inteligência Computacional Aplicada (opt) \\ \hline
    Luis José Silveira de Sousa & Professor / DE  &  Doutorado Física  &  Física I, Física II \\ \hline
    Luiz Daniel Santos Bezerra & Professor / DE  &  Doutorado Eng. Elétrica  &  Eletrônica I, Eletrônica II, Eletrônica III, Controle II, Métodos Numéricos, Sistemas Embarcados (opt)  \\ \hline
    %Pedro Pedrosa Rebouças Filho & Professor / DE  &  Doutor Eng. Teleinf.  & Lógica de Programação, Linguagem de Programação, Robótica I, Sistemas Embarcados, Robótica II (opt), Processamento Digital de Imagens (opt) \\ \hline
    %Rodrigo Freitas Guimarães   &  Professor / DE  &  Doutor Eng. Materiais  &  Desenho Técnico, Desenho Auxiliado por Computador, Manufatura Auxiliada por Computador, Manufatura Integrada por Computador (opt) \\ \hline
    %Samuel Vieira Dias & Professor / DE  &  Doutor Eng. Elétrica  &  Métodos Numéricos, Modelagem de Sistemas a Eventos Discretos (opt), Controle III (opt) \\ \hline 
    Márcio Monteiro Cunha & Professor / DE  &  Mestrado Planejamento e Políticas Públicas  &  Metodologia Científica  \\ \hline
    Marcos Cirineu Aguiar Siqueira & Professor / DE  &  Especialização em Pesquisa Científica  &  Probabilidade e Estatística  \\ \hline
    Narcélio de Araújo Pereira & Professor / DE  &  Mestrado Eng. Produção  &  Higiene e Segurança do Trabalho (HST)  \\ \hline
    Rafael Oliveira de Sousa  &  Professor / DE &  Mestrado Eng. Elétrica  &  Introdução à Engenharia, Circuitos Elétricos I, Circuitos Elétricos II, Acionamentos de Máquinas, Instalações Elétricas   \\ \hline
    Thomas de Oliveira Praxedes  &  Professor / DE &  Mestrado Eng. Materiais  &  Resistência dos Materiais I   \\ \hline
    Venceslau Xavier de Lima Filho  &  Professor / DE &  Doutorado Eng. Materiais  &  Materiais I   \\ \hline
    Venício Soares de Oliveira  &  Professor / DE &  Mestrado Eng. Elétrica  &  Acionamentos Hidráulicos e Pneumáticos, Controladores Lógicos Programáveis (CLP)   \\ \hline
    Weber Chaves Fontoura  &  Professor / DE &  Mestrado Eng. de Produção  & Gestão e Controle da Qualidade, Gestão da Manutenção Industrial \\ \hline
    Teófilo Roberto da Silva & Professor / DE  &  Doutorado Linguística &  Inglês Técnico, Metodologia Científica \\ \hline
   \end{longtable}
\else
    \begin{longtable}{|p{3.0cm}|p{2.5cm}|p{2.0cm}|p{6.3cm}|}
    \caption{Corpo Docente Existente no Campus}\label{tab2} 
    \vspace{-0.25cm}
    %\multicolumn{4}{c}{\textbf{Tabela 9. Corpo Docente Existente no Campus}}    \label{tab2} 
    \\ \hline 
    \cellcolor{\cinza}NOME & \cellcolor{\cinza}QUAL.PROF. / REG.TRAB. & \cellcolor{\cinza}TÍTULO & \cellcolor{\cinza}DISCIPLINAS MINISTRADAS \\ 
    \hline
    David Carneiro de Souza   &  Professor / DE  &  Doutor Matemática &  Cálculo I, Cálculo II, Cálculo III, Álgebra Linear, Probabilidade e Estatística, Equações Diferenciais Ordinárias (opt) \\ \hline
    Luis José Silveira de Sousa & Professor / DE  &  Doutor Física  &  Física I, Física II, Física III \\ \hline
    Francisco Nélio Costa Freitas  &  Professor / DE  & Doutor Eng. Materiais  &  Introdução à Engenharia, Química Geral, Materiais I, Mecanismos, Corrosão e Proteção Anti-Corrosiva (opt) \\ \hline
    Francisco Frederico dos Santos Matos   &  Professor / DE  &  Doutor Eng. Mecânica  &  Termodinâmica, Mecânica dos Fluidos, Transferência de Calor, Máquinas Térmicas e de Fluxo, Refrigeração Industrial \\ \hline
    Rodrigo Freitas Guimarães   &  Professor / DE  &  Doutor Eng. Materiais  &  Materiais II, Soldagem, Transportadores Industriais (opt) \\ \hline
    Pedro Pedrosa Rebouças Filho & Professor / DE  &  Doutor Eng. Teleinformática  & Lógica de Programação, Linguagem de Programação, Processamento Digital de Imagens (opt), Inteligência Computacional Aplicada (opt) \\ \hline
    Samuel Vieira Dias & Professor / DE  &  Doutor Eng. Elétrica  &  Métodos Numéricos, Eletrônica I, Circuitos Elétricos I \\ \hline
   Fabrício Bandeira da Silva  &  Professor / DE & Doutor Educação  &  Planejamento e Controle da Produção, Gestão e Controle da Qualidade, Gestão da Manutenção Industrial, Inspeção Industrial, Gestão de Projetos (opt)    \\ \hline
    José Ciro dos Santos & Professor / DE  &  Mestre Ener. Renováveis  &  Desenho Técnico, Desenho Auxiliado por Computador, Manufatura Auxiliada por Computador, Manufatura Integrada por Computador (opt) \\ \hline
    Venceslau Xavier de Lima e Filho   &  Professor / DE  &  Doutor Eng. Materiais  &  Resistência dos Materiais I, Resistência dos Materiais II, Sistemas Mecânicos, Análise de Falhas, Tribologia e Lubrificação (opt) \\ \hline
    Francisco José dos Santos Oliveira   &  Professor / DE  &  Mestre Eng. Materiais &  Metrologia, Processos de Fabricação I, Processos de Fabricação II, Manutenção Industrial \\ \hline
    Adriano Holanda Pereira & Professor / DE  &  Mestre Eng. Elétrica  &  Instalações Elétricas, Máquinas Elétricas  \\ \hline
    Venício Soares de Oliveira  &  Professor / DE &  Mestre Eng. Elétrica  &  Acionamentos Hidráulicos e Pneumáticos, TCC, Projeto de Engenharia (opt)    \\ \hline
    Weber Chaves Fontoura  &  Professor / DE &  Mestre Eng. de Produção  &  Introdução à Engenharia, Planejamento e Controle da Produção, Gestão e Controle da Qualidade e Gestão da Manutenção Industrial \\ \hline
    Antônio Olívio Silveira Britto Júnior & Professor / DE  &  Mestre Eng. Civil  &  HST, Ética e Cidadania, Projetos Sociais  \\ \hline
    Eugênio Barreto Sousa e Silva & Professor / DE  &  Mestre Administração &  Empreendedorismo \\ \hline
    Teófilo Roberto da Silva & Professor / DE  &  Doutor Linguística &  Inglês Técnico, Metodologia Científica \\ \hline
    Adriano Barros Carneiro &  Professor / DE  &  Mestre Desporto  &  Educação Física \\ \hline
   \end{longtable}
\fi



\subsection{Dados Complementares dos Docentes do Eixo da Indústria}

Nesta seção são apresentadas informações adicionais sobre todos os docentes acima relacionados na Tab.\ref{tab2} e, além disso, também são listados outros docentes que também fazem parte do eixo da indústria do Campus Maracanaú, onde os mesmos colaboram com o curso de forma indireta.

\begin{itemize}

\item Nome Completo: Adriano Holanda Pereira

 CPF: 797.256.783-20

 Titulação máxima: (  ) Graduado (  ) Especialista  (X) Mestrado ( ) Doutorado

 Regime de trabalho: (X) DE ( ) 40horas ( ) 20horas

 Vínculo empregatício: (  ) Substituto (X) Efetivo  (  ) Voluntário

 Formação acadêmica: Graduação em Engenharia Elétrica e Mestrado em Engenharia Elétrica

 Endereço eletrônico do Lattes: \web{http://lattes.cnpq.br/9075415197983567}


\item Nome Completo: Adriano Barros Carneiro

 CPF: 032.097.813-38

 Titulação máxima: (  ) Graduado (  ) Especialista  (X) Mestrado ( ) Doutorado

 Regime de trabalho: (X) DE ( ) 40horas ( ) 20horas

 Vínculo empregatício: (  ) Substituto (X) Efetivo  (  ) Voluntário

 Formação acadêmica: Graduação em Educação Física, Mestrado em Ciência do Desporto e Doutorando em Ciência da Educação (previsão de conclusão do doutorado: 2018)

 Endereço eletrônico do Lattes: \web{http://lattes.cnpq.br/1063180174442960}

 
\item Nome Completo: Antonio Barbosa de Sousa Júnior

 CPF: 001.383.073-26

 Titulação máxima: (  ) Graduado  (  ) Especialista  (  ) Mestrado ( X ) Doutorado

 Regime de trabalho: ( X ) DE  (  ) 40horas (  ) 20horas

 Vínculo empregatício: (  ) Substituto  ( X ) Efetivo  (  ) Voluntário

 Formação acadêmica: Graduação em Engenharia de Controle e Automação, Mestrado e Doutodado em Engenharia Elétrica.

 Endereço eletrônico do Lattes: \web{http://lattes.cnpq.br/4260373627927424}

\item Nome Completo: Antônio Olívio Silveira Britto Júnior

 CPF: 234.457.933-87

 Titulação máxima: (  ) Graduado (  ) Especialista  (X) Mestrado ( ) Doutorado

 Regime de trabalho: (X) DE ( ) 40horas ( ) 20horas

 Vínculo empregatício: (  ) Substituto (X) Efetivo  (  ) Voluntário

 Formação acadêmica: Graduação em Agronomia, Especialização em Segurança do Trabalho e Mestrado em Engenharia Civil

 Endereço eletrônico do Lattes: \web{http://lattes.cnpq.br/3534055362942935}


\item Nome Completo: Celso Rogério Schmidlin Júnior

 CPF: 839.085.793-68

 Titulação máxima: (  ) Graduado  (  ) Especialista  (  ) Mestrado ( X ) Doutorado

 Regime de trabalho: (X) DE  (  ) 40horas (  ) 20horas

 Vínculo empregatício: (  ) Substituto  (X) Efetivo  (  ) Voluntário

 Formação acadêmica: Graduação em Engenharia Elétrica, Mestrado em Engenharia Elétrica e Doutorado em Engenharia
Elétrica

 Endereço eletrônico do Lattes: \web{http://lattes.cnpq.br/9279561215102176}


\item Nome Completo: David Carneiro de Souza

 CPF: 658.489.183-68

 Titulação máxima: (  ) Graduado  (  ) Especialista  (  ) Mestrado ( X ) Doutorado

 Regime de trabalho: ( X ) DE  ( ) 40horas (  ) 20horas

 Vínculo empregatício: (  ) Substituto  ( X ) Efetivo  (  ) Voluntário

 Formação acadêmica: Bacharelado em Matemática, Mestrado em Matemática e Doutorado em Matemática

 Endereço eletrônico do Lattes: \web{http://lattes.cnpq.br/9511010025717986}

\item Nome Completo: Eugênio Barreto Sousa e Silva

 CPF: 462.339.203-10

 Titulação máxima: (  ) Graduado (  ) Especialista  (X) Mestrado ( ) Doutorado

 Regime de trabalho: (X) DE ( ) 40horas ( ) 20horas

 Vínculo empregatício: (  ) Substituto (X) Efetivo  (  ) Voluntário

 Formação acadêmica: Graduação em Administração

 Endereço eletrônico do Lattes: \web{http://lattes.cnpq.br/5492252517721394}


\item Nome Completo: Fábio Timbó Brito

 CPF: 804.756.393-91

 Titulação máxima: (  ) Graduado  (  ) Especialista  ( X ) Mestrado (  ) Doutorado

 Regime de trabalho: ( X ) DE  (  ) 40horas (  ) 20horas

 Vínculo empregatício: (  ) Substituto  ( X ) Efetivo  (  ) Voluntário

 Formação acadêmica: Graduação Tecnológica em Mecatrônica, Especialização em Automação Industrial e Mestrado em
Engenharia Elétrica

 Endereço eletrônico do Lattes: \web{http://lattes.cnpq.br/1030417130389338}


\item Nome Completo: Fabrício Bandeira da Silva

 CPF: 619.669.593-20

 Titulação máxima: (  ) Graduado (  ) Especialista  (  ) Mestrado ( X ) Doutorado

 Regime de trabalho: ( X ) DE  (  ) 40horas (  ) 20horas

 Vínculo empregatício: (  ) Substituto  ( X ) Efetivo  (  ) Voluntário

 Formação acadêmica: Graduação em Engenharia Elétrica, Mestrado em Engenharia Elétrica e Doutorado em Educação

 Endereço eletrônico do Lattes: \web{http://lattes.cnpq.br/6536228936628118}


\item Nome Completo: Francisco Frederico dos Santos Matos

 CPF: 616.526.873-00

 Titulação máxima: (  ) Graduado  (  ) Especialista  (  ) Mestrado ( X ) Doutorado

 Regime de trabalho: ( X ) DE  (  ) 40horas (  ) 20horas

 Vínculo empregatício: (  ) substituto  ( X ) efetivo  (  ) Voluntário

 Formação acadêmica: Graduação em Engenharia Mecânica e Doutorado em Engenharia Mecânica

 Endereço eletrônico do Lattes: \web{http://lattes.cnpq.br/8204044573317813}

\item Nome Completo: Francisco José dos Santos Oliveira

 CPF: 491.827.143-04

 Titulação máxima: (  ) Graduado  (  ) Especialista  ( X ) Mestrado (  ) Doutorado

 Regime de trabalho: ( X ) DE  ( ) 40horas (  ) 20horas

 Vínculo empregatício: (  ) Substituto  ( X ) Efetivo  (  ) Voluntário

 Formação acadêmica: Graduação em Engenharia Mecânica, Especialização em Engenharia de Segurança do Trabalho, Mestrado em Engenharia e Ciência de Materiais e Doutorando em Engenharia e Ciência de Materiais (previsão para conclusão do doutorado: 2019)

 Endereço eletrônico do Lattes: \web{http://lattes.cnpq.br/3399911160895219}


\item Nome Completo: Francisco Nélio Costa Freitas

 CPF: 708.485.563-72

 Titulação máxima: (  ) Graduado  (  ) Especialista  (  ) Mestrado  ( X ) Doutorado

 Regime de trabalho: ( X ) DE  (  ) 40 horas  (  ) 20 horas

 Vínculo empregatício: (  ) Substituto  (  X ) Efetivo  (  ) Voluntário

 Formação acadêmica: Graduação em Engenharia Mecânica, Mestrado em Engenharia e Ciência de Materiais e Doutorado em Engenharia e Ciência de Materiais

 Endereço eletrônico do Lattes: \web{http://lattes.cnpq.br/1834964619080647}


\item Nome Completo: José Daniel de Alencar Santos

 CPF: 559.959.263-49

 Titulação máxima: (  ) Graduado (  ) Especialista  (  ) Mestrado ( X ) Doutorado

 Regime de trabalho: (X) DE ( ) 40horas ( ) 20horas

 Vínculo empregatício: (  ) Substituto (X) Efetivo  (  ) Voluntário

 Formação acadêmica: Graduação em Engenharia Elétrica, Especialização em Engenharia Elétrica, Mestrado em Engenharia de Teleinformática e Doutorado em Engenharia de Teleinformática

 Endereço eletrônico do Lattes: \web{http://lattes.cnpq.br/1123307163858590}


\item Nome Completo: Luiz Daniel Santos Bezerra

 CPF: 977.064.183-91

 Titulação máxima: (  ) Graduado  (  ) Especialista  (  ) Mestrado ( X ) Doutorado

 Regime de trabalho: ( X ) DE  (  ) 40horas (  ) 20horas  (  ) Voluntário.

 Vínculo empregatício: (  ) Substituto  ( X ) Efetivo  (  ) Voluntário

 Formação acadêmica: Graduação em Engenharia Elétrica, Mestrado em Engenharia Elétrica e Doutorado em Engenharia Elétrica

 Endereço eletrônico do Lattes: \web{http://lattes.cnpq.br/3261052136129353}


\item Nome Completo: Luis José Silveira de Sousa

 CPF: 820.207.793-15

 Titulação máxima: (  ) Graduado  (  ) Especialista  (  ) Mestrado ( X ) Doutorado

 Regime de trabalho: (X) DE  (  ) 40horas (  ) 20horas

 Vínculo empregatício: (  ) Substituto  (X) Efetivo  (  ) Voluntário

 Formação acadêmica: Licenciatura Plena em Física, Mestrado em Ciências Físicas Aplicadas e Doutorado em Física

 Endereço eletrônico do Lattes: \web{http://lattes.cnpq.br/6055826430173488}

%\item Nome Completo: Marcelo Monteiro Valente Parente
%
% CPF: 060.532.623-15
%
% Titulação máxima: (  ) Graduado (  ) Especialista  (  ) Mestrado ( X ) Doutorado
%
% Regime de trabalho: ( X ) DE  (  ) 40horas (  ) 20horas
%
% Vínculo empregatício: (  ) Substituto  ( X ) Efetivo  (  ) Voluntário
%
% Formação acadêmica: Graduação em Engenharia Química, Mestrado em Engenharia Metalúrgica e de Materiais e Doutorado em Química Inorgânica
%
% Endereço eletrônico do Lattes: \web{http://lattes.cnpq.br/6890118483209098}

\item Nome Completo: Pedro Pedrosa Rebouças Filho

 CPF: 007.626.423-86

 Titulação máxima: (  ) Graduado (  ) Especialista  (  ) Mestrado ( X ) Doutorado

 Regime de trabalho: ( X ) DE  (  ) 40horas (  ) 20horas

 Vínculo empregatício: (  ) Substituto  ( X ) Efetivo  (  ) Voluntário

 Formação acadêmica: Graduação Tecnológica em Mecatrônica, Mestrado em Engenharia de Teleinformática e Doutorado em
Engenharia de Teleinformática

 Endereço eletrônico do Lattes: \web{http://lattes.cnpq.br/4347965302097614}


\item Nome Completo: Rodrigo Freitas Guimarães

 CPF: 620.155.403-34

 Titulação máxima: (  ) Graduado  (  ) Especialista  (  ) Mestrado ( X ) Doutorado

 Regime de trabalho: ( X ) DE  (  ) 40horas (  ) 20horas

 Vínculo empregatício: (  ) Substituto  ( X ) Efetivo  (  ) Voluntário

 Formação acadêmica: Graduação em Engenharia Mecânica, Mestrado em Engenharia e Ciência de Materiais e Doutorado em
Engenharia e Ciência de Materiais

 Endereço eletrônico do Lattes: \web{http://lattes.cnpq.br/1434906331576002}


\item Nome Completo: Samuel Vieira Dias

 CPF: 856.463.343-49

 Titulação máxima: (  ) Graduado (  ) Especialista  (  ) Mestrado ( X ) Doutorado

 Regime de trabalho: ( X ) DE  (  ) 40horas (  ) 20horas

 Vínculo empregatício: (  ) Substituto  ( X ) Efetivo  (  ) Voluntário

 Formação acadêmica: Técnico em Mecânica Industrial, Graduação Tecnológica em Mecatrônica, Especialização em Educação, Mestrado em Engenharia Elétrica e Doutorado em Engenharia Elétrica

 Endereço eletrônico do Lattes: \web{http://lattes.cnpq.br/2651422114397357}


\item Nome Completo: Venceslau Xavier de Lima Filho

 CPF: 619.539.703-20

 Titulação máxima: (  ) Graduado  (  ) Especialista  (  ) Mestrado ( X ) Doutorado

 Regime de trabalho: ( X ) DE  (  ) 40horas (  ) 20horas

 Vínculo empregatício: (  ) Substituto  ( X ) Efetivo  (  ) Voluntário

 Formação acadêmica: Graduação em Engenharia Mecânica, Mestrado em Ciências e Engenharia dos Materiais e Doutorado em Engenharia e Ciência de Materiais

 Endereço eletrônico do Lattes: \web{http://lattes.cnpq.br/0209565051697344}


\item Nome Completo: Venício Soares de Oliveira

 CPF: 619.539.703-20

 Titulação máxima: (  ) Graduado  (  ) Especialista  ( X ) Mestrado (  ) Doutorado

 Regime de trabalho: ( X ) DE  (  ) 40horas (  ) 20horas

 Vínculo empregatício: (  ) Substituto  ( X ) Efetivo  (  ) Voluntário

 Formação acadêmica: Graduação em Tecnologia Eletromecânica e Mestrado em Engenharia Elétrica.

 Endereço eletrônico do Lattes: \web{http://lattes.cnpq.br/3263289316174386}
 
 \item Nome Completo: Weber Chaves Fontoura

 CPF: 827.281.273-87

 Titulação máxima: (  ) Graduado  (  ) Especialista  ( X ) Mestrado (  ) Doutorado

 Regime de trabalho: ( X ) DE  (  ) 40horas (  ) 20horas

 Vínculo empregatício: (  ) Substituto  ( X ) Efetivo  (  ) Voluntário

 Formação acadêmica: Graduação em Geografia pela Universidade Estadual do Ceará (2002) e em Engenharia Mecânica pela Universidade Federal do Ceará (2012) e Mestrado em Engenharia de Produção (Logística e Pesquisa Operacional) pela Universidade Federal do Ceará (2016).

 Endereço eletrônico do Lattes: \web{http://lattes.cnpq.br/9713957219078309}

\end{itemize}

% CORPO TÉCNICO ADMINISTRATIVO ------------------------------------------------

\clearpage
\section{CORPO TÉCNICO-ADMINISTRATIVO}

O corpo técnico-administrativo diretamente relacionado ao curso é apresentado na Tab. \ref{tab3}. Nela constam os nomes dos servidores, os setores nos quais estão lotados, seus cargos e titulação máxima.

%A Tab. \ref{tab3} apresenta os servidores técnico-administrativos do campus Maracanaú.

\begin{longtable}{|p{6cm}|p{5.5cm}|p{2.3cm}|}
\caption{Servidores Técnico-Administrativos Diretamente Relacionados ao Curso}\label{tab3} 
\vspace{-0.25cm}
%\multicolumn{4}{c}{\textbf{Servidores Técnico-Administrativos do Campus de Maracanaú}}\label{tab3} 
\\ \hline 
\cellcolor{\cinza}NOME & \cellcolor{\cinza}CARGO & \cellcolor{\cinza}TITULAÇÃO \\ 
\hline

%\begin{table}[H]
%\caption{Servidores Técnico-Administrativos do Campus de Maracanaú}
%\label{tab3}
%\begin{tabularx}{1.0\textwidth}{X|X|X|c}
%\hline
%\textbf{NOME} & \textbf{SETOR} & \textbf{CARGO} & \textbf{TITULAÇÃO} \\ \hline

\multicolumn{3}{|c|}{\textbf{Direção de Ensino}} \\
\hline
Anna Hilda Silva Melo 

(Assistente da Direção de Ensino) & Assitente em Administração & Especialista \\ \hline

\multicolumn{3}{|c|}{\textbf{Coordenadoria de Controle Acadêmico}} \\
\hline
Cristiano do Nascimento Lira & Auxiliar em Administração &	Graduado \\ \hline
Elder Kened Cardoso 

(Coordenador) &	Auxiliar em Administração & Graduado \\ \hline 
Débora Natazia Moreira Barbosa & Auxiliar em Administração & Graduada \\ \hline

\multicolumn{3}{|c|}{\textbf{Estágio}} \\
\hline
Andreia Cavalcante Rodrigues & Assistente em Administração & Especialista  \\ \hline
Sergina Mendes da Silva Flor & Assistente em Administração & Especialista \\ \hline

\multicolumn{3}{|c|}{\textbf{Coordenadoria Técnico-Pedagógica}} \\
\hline
Isabel Magda Said Pierre Carneiro 

(Coordenadora) & Pedagoga & Doutora \\ \hline
Leilane Lima Almeida Evangelista & Técnica em Assuntos Educacionais & Especialista \\ \hline
Roseane Michelle de Lima Silveira & Pedagoga & Especialista \\ \hline
Samoel Rodrigues da Silva & Técnico em Assuntos Educacionais & Especialista \\ \hline

\multicolumn{3}{|c|}{\textbf{Biblioteca}} \\
\hline
Luiz Carlos Silveira de Sousa 

(Coordenador) & Bibliotecário & Especialista \\ \hline 
Gláucio Barreto Lima & Bibliotecário & Mestre \\ \hline
Francisca Marta Mendes Oliveira & Auxiliar de Biblioteca & Especialista \\ \hline 
Antonia Ney da Silva Pereira & Auxiliar de Biblioteca & Nivel Médio \\ \hline

\multicolumn{3}{|c|}{\textbf{Coordenadoria de Assuntos Estudantis}} \\
\hline
Márcia Lorena Bezerra Peixoto 

(Coordenadora) & Assistente Social & Especialista \\ \hline
Keyla de Souza Lima Cruz & Assistente Social & Doutora \\ \hline
Diego Bastos do Nascimento Martins & Nutricionista & Mestre \\ \hline
Agnes Caroline Souza Pinto & Enfermeira & Doutora \\ \hline
Renata Alves Albuquerque & Psicóloga & Doutora \\ \hline
Lucélia Fernandes de Almeida Lima & Técnica de Enfermagem & Graduada \\ \hline

\multicolumn{3}{|c|}{\textbf{Setor de Educação Física e Esportes}} \\
\hline
Francisco Hermison Monteiro do Vale & Assistente em Administração & Graduado \\ \hline 
Sena Moreira do Nascimento & Assistente em Administração & Ensino Médio  \\ \hline

\multicolumn{3}{|c|}{\textbf{Núcleo de Atendimento às Pessoas com Necessidades Específicas}} \\
\hline
Emanuel Bruno Carioca Silva & Técnico em Tradução e Interpretação de LIBRAS & Nível Médio/ Nível Técnico  \\ \hline

\multicolumn{3}{|c|}{\textbf{Departamento de Extensão, Pesquisa, Pós-Graduação e Inovação}} \\
\hline
Iassodara Farias Leitão Pessoa & Assistente em Administração & Especialista \\ \hline
Débora Viana de Araújo & Assistente em Administração & Graduada \\ \hline
%\end{tabularx}
%\footnotesize{\textbf{Fonte:} A Comissão de Elaboração.}
\end{longtable}

% INFRAESTRUTURA ----------------------------------------------

\clearpage
\section{INFRAESTRUTURA}
\label{INFRA}
Para execução das atividades do \NOMECURSO está contemplada a seguinte infraestrutura:

\subsection{Laboratórios}

\begin{itemize}%[LAMEP]
\item [LMAT] LABORATÓRIO DE MATERIAIS: São realizadas análises microestruturais, tratamentos térmicos, ensaios mecânicos e ensaios de
corrosão. Disciplinas contempladas: \LMATdisc.

\item [LIAF] LABORATÓRIO DE INSPEÇÃO E ANÁLISE DE FALHAS: São realizados ensaios não-destrutivos com equipamentos tais como
ultra-som, kit de inspeção por partículas magnéticas, kit de inspeção por líquido penetrante, endoscópio industrial,
viscosímetro e espectrômetro de fluorescência de raios-X. Disciplinas contempladas: \LIAFdisc.

\item [LMET] LABORATÓRIO DE METROLOGIA DIMENSIONAL: São realizadas medições com uso de paquímetros digitais e analógicos, medições
com uso de micrômetros digitais e analógicos, medições com uso de goniômetros, medições com uso de relógios
comparadores digitais e analógicos, medições com uso de rugosímetros, medições com projetor de perfil, medições com uso
de traçadores de altura, medições com uso de réguas graduadas, medições e engenharia reversa com tecnologia ótica 2D e
medições por coordenadas e engenharia reversa com tecnologia de scanner a laser.  Disciplinas contempladas: \LMETdisc.

\item [LTF] LABORATÓRIO DE MÁQUINAS TÉRMICAS E DE FLUXO: Laboratório cuja finalidade é propiciar aos alunos uma visualização prática
do princípio de funcionamento das máquinas de fluxo e térmicas. Também ensaiam-se experimentos práticos e
computacionais que servirão de base para produção científica e desenvolvimento de mão de obra especializada. Disciplinas contempladas:\LTFdisc.

\item [LIA] LABORATÓRIO DE INFORMÁTICA APLICADA: Tem por finalidade desenvolver atividades voltadas à simulação de circuitos
elétricos e eletrônicos, desenvolvimento de algoritmos computacionais, bem como a criação e desenvolvimento de modelos
em 2D e 3D. Disciplinas contempladas: \LIAdisc.

\item [LSHIP] LABORATÓRIO DE SISTEMAS HIDRÁULICOS E PNEUMÁTICOS: Tem por finalidade possibilitar que os alunos conheçam os principais
atuadores pneumáticos e hidráulicos, dispositivos de geração e tratamento de ar comprimido, simular e implementar de
forma experimental circuitos pneumáticos e hidráulicos. Disciplinas contempladas: \LSHIPdisc.

\item [LAMEP] LABORATÓRIO DE ACIONAMENTOS DE MÁQUINAS E ELETRÔNICA DE POTÊNCIA: Este laboratório dispõe de equipamentos e dispositivos
que permitem a realização de aulas práticas em circuitos elétricos em corrente alternada (cargas R, RL, RC, RLC),
comandos e instalações elétricas em âmbito industrial, acionamentos e controladores industriais aplicados em motores de
indução monofásicos, trifásicos e servomotores CA, com e sem carga. Disciplinas contempladas:\LAMEPdisc.

\item [LEE] LABORATÓRIO DE ELETROELETRÔNICA: Este laboratório dispõe de equipamentos e dispositivos que permitem a realização de aulas práticas em circuitos elétricos em corrente contínua, bem como medições elétricas das principais grandezas (resistência, tensão, corrente, etc), módulos para práticas em circuitos de eletrônica digital e analógica, envolvendo amplificadores operacionais, diodos, transistores, etc. Disciplinas contempladas:\LEEdisc.

\item [LINC] LABORATÓRIO DE INSTRUMENTAÇÃO E CONTROLE: Tem por finalidade a realização de atividades e projetos da área de instrumentação e controle, através da utilização de equipamentos que simulam situações reais do ambiente industrial. Através~ de modernos sistemas eletrônicos de calibração e medição de variáveis de processo como nível, pressão, vazão e temperatura. Possibilita ao aluno associar a teoria aos conhecimentos práticos e reais da instrumentação industrial.
Disciplinas contempladas:\LINCdisc.

%\item [LPROT] LABORATÓRIO DE PROTÓTIPOS: Tem por finalidade permitir ao aluno a realização de projetos de protótipos de processos elétricos e eletrônicos, através da simulação~ em~ escala real ou reduzida, possibilitando uma posterior utilização em ambientes industriais ou sistemas embarcados.
%    Disciplinas contempladas:\LPROTdisc.

\item [LPC] LABORATÓRIO DE POTÊNCIA E CONTROLE: Tem por finalidade permitir ao aluno a realização de projetos de protótipos de processos elétricos e eletrônicos, através da simulação e concepção em escala real ou reduzida, possibilitando uma posterior utilização em ambientes industriais ou sistemas embarcados.
    Disciplinas contempladas:\LPCdisc.

\item [LAMSC] LABORATÓRIO DE APRENDIZAGEM DE MÁQUINAS E SIMULAÇÃO COMPUTACIONAL: Tem o objetivo de permitir ao aluno implementar e testar modelos computacionais para aplicações de aprendizagem de máquinas em problemas de controle e identificação de sistemas. Além disso, também possibilita ao aluno a montagem de projetos de protótipos em microcontroladores e robótica.
    Disciplinas contempladas:\LAMSCdisc.


%\item [LEE] LABORATÓRIO DE ELETROELETRÔNICA: Este laboratório dispõe de equipamentos e dispositivos que permitem a realização de
%aulas práticas em circuitos elétricos em corrente contínua, bem como medições elétricas das principais grandezas
%(resistência, tensão, corrente, etc), módulos para práticas em circuitos de eletrônica digital e analógica, envolvendo
%amplificadores operacionais, diodos, transistores, etc. Disciplinas contempladas: Eletricidade I, Eletrônica Analógica
%e Eletrônica Digital.
%
%\item [LINC] LABORATÓRIO DE INSTRUMENTAÇÃO E CONTROLE: Tem por finalidade a realização de atividades e projetos da área de
%instrumentação e controle, através da utilização de equipamentos que simulam situações reais do ambiente industrial.
%Através~ de modernos sistemas eletrônicos de calibração e medição de variáveis de processo como nível, pressão, vazão e
%temperatura, torna o aluno capaz de associar a teoria~ aos conhecimentos práticos e reais da instrumentação industrial.
%Disciplinas contempladas: Instrumentação e Controle, Redes Industriais e CLP.
%
%\item [LPROT] LABORATÓRIO DE PROTÓTIPOS: Tem por finalidade permitir ao aluno a realização de projetos de protótipos de processos
%elétricos e eletrônicos, através da simulação~ em~ escala real ou reduzida, possibilitando uma posterior utilização em
%ambientes industriais ou sistemas embarcados. Contempla todas as disciplinas do curso. (J.DANIEL)

\end{itemize}

\subsection{Ambientes Administrativos}

%A infraestrutura de salas de aula e laboratórios, %, além de complementar a teoria aprendida em sala de aula e laboratórios, também irá auxiliar
Os alunos do \NOMECURSO contam com uma infraestrutura administrativa composta de: % no desenvolvimento de atividades de pesquisa e extensão.

\begin{itemize}
\item SALA DE PROFESSORES
\item SALA DE ATENDIMENTO AO ALUNO
\item SALA DA COORDENADORIA DOS CURSOS DO EIXO TECNOLÓGICO DA INDÚSTRIA
\item 5 SALAS DE AULA
\item GABINETE DE PROFESSORES
\item AUDITÓRIO
\item SALA DE VIDEOCONFERÊNCIA
\item BIBLIOTECA
\item CENTRO ACADÊMICO DO CURSO
\end{itemize}

Essa infraestrutura também auxilia os alunos do \NOMECURSO no desenvolvimento de atividades de pesquisa e extensão.

\subsection{Bloco II do Eixo Tecnológico da Indústria}
O Bloco II do Eixo Tecnológico da Indústria, denominado Centro de Pesquisa e Tecnologia (CPT), conta com os seguintes ambientes:
\begin{itemize}
\item \textbf{Finalizados:} ajustagem, soldagem, manutenção eletromecânica, usinagem convencional, usinagem CNC, Laboratório de CAD e CAM;
\item \textbf{Em construção:} 08 (oito) laboratórios de pesquisa e desenvolvimento, 15 (quinze) gabinetes para professores, biblioteca setorial, auditório e salas de aula para futuros
programas de pós-graduação.
\end{itemize}

 Essa nova infraestrutura contribuirá de maneira extremamente significativa para a constante
atualização tecnológica do \NOMECURSO. Maiores detalhes podem ser apreciados no projeto do CPT apresentado no ANEXO \ref{anexoProjeto}.

\clearpage
\section{REFERÊNCIAS}
\bibliography{referencias}

%%%%%%%%%%%
\clearpage\appendix
\pagestyle{empty}

%%%%%%%%%%%
\clearpage\section{COMPROVANTE DO IMÓVEL}
\label{anexoImovel}
\includepdf[pages={3}]{anexos/Comprovante Imovel.pdf}

%%%%%%%%%%%
%\clearpage\section{RESOLUÇÕES E PORTARIAS}
\clearpage\section{PORTARIAS}
\label{anexoPortaria}

%\clearpage
\subsection{Portaria \PORTARIACOMISSAO}
\label{anexoPortariaComissao}
\if\COMPILAR0
	\includepdf[pages={1}]{anexos/Portaria23-27032013.pdf}
\else
	\includepdf[pages={1}]{anexos/Portaria24-27032013.pdf}
\fi

\subsection{Portaria \PORTARIACOMISSAOALT}
\label{anexoPortariaComissaoAlt}
\if\COMPILAR0
	\includepdf[pages={1}]{anexos/Portaria42-22042013.pdf}
\else
	\includepdf[pages={1}]{anexos/Portaria43-22042013.pdf}
\fi


\subsection{Portaria \PORTARIARECONHECIMENTO}
\label{anexoPortariaReconhecimento}
\if\COMPILAR0
\includepdf[pages={1-3}]{anexos/Portaria-430-03052021.pdf}
\else
\includepdf[pages={1-3}]{anexos/Portaria-430-03052021.pdf}
\fi

%\subsection{Portaria \PORTARIANDEATUAL}
%\label{anexoPortariaNDEATUAL}
%\if\COMPILAR0
%	\includepdf[pages={1}]{anexos/PORT093_NDE_ECA.pdf}
%\else
%	\includepdf[pages={1}]{anexos/PORT094_NDE_EME.pdf}
%\fi

%\subsection{Portaria \PORTARIACOLEGIADOATUAL}
%\label{anexoPortariaCOLEGIADOATUAL}
%\if\COMPILAR0
%	\includepdf[pages={1}]{anexos/PORT092_COLEGIADO_ECA.pdf}
%\else
%	\includepdf[pages={1}]{anexos/PORT091_COLEGIADO_EME.pdf}
%\fi


%\clearpage\subsection{Portarias n° 099/GDG de 23 de dezembro de 2010 e n° 086/GDG de 26 outubro de 2011}
%\label{anexoPortariaColegiado}
%
%\includepdf[pages={1}]{myfile.pdf}
%
%\clearpage\subsection{Portarias n° 098/GDG  de 23 de dezembro de 2010 e n° 083/GDG de 13 de outubro de 2011}
%\label{anexoPortariaNDE}
%
%\includepdf[pages={1}]{myfile.pdf}

%%%%%%%%%%%
\clearpage\section{REGULAMENTAÇÃO DA ORGANIZAÇÃO DIDÁTICA}
\label{anexoROD}

Reprodução de trechos do ROD que versam sobre matrícula, transferência e avaliação. O documento eletrônico original está disponível para acesso em \web{http://ifce.edu.br/espaco-estudante/regulamento-de-ordem-didatica}.

\subsection{Da Matrícula: Seções I e II, Capítulo II, Título III}
\label{anexoRODmatricula}

\begin{center}
Seção I - Da matrícula inicial
\end{center}

Art. 75. Matrícula é o ato formal pelo qual se dá a vinculação acadêmica do estudante ao IFCE após classificação em processo seletivo e convocação conforme número de vagas disponíveis, mediante apresentação dos documentos exigidos no edital.

Art. 76. Considera-se como matrícula inicial aquela realizada no período letivo de ingresso do estudante no IFCE para os cursos técnicos (integrados, concomitantes ou subsequentes) e de graduação (bacharelado, licenciatura ou tecnologia).

Art. 77. A matrícula inicial deverá ser efetivada de forma presencial pelo candidato classificado, quando maior de 18 (dezoito) anos, ou por seu representante legal, quando menor de 18 (dezoito) anos.

§ 1º Na ausência do estudante maior de 18 anos, a solicitação poderá ser realizada por seu representante legal, desde que apresente procuração com firma reconhecida.

§ 2º Na ausência do responsável legal pelo estudante menor que 18 anos, solicitação poderá ser realizada pelo representante do responsável legal, desde que apresente procuração com firma reconhecida.

§ 3º Uma vez realizada a matrícula pelo estudante, o horário da oferta dos componentes curriculares não poderá ser alterado. Em casos excepcionais, a alteração acontecerá somente mediante autorização do gestor máximo do ensino no campus.

Art. 78. Nos cursos de graduação do IFCE, é obrigatório ao estudante se matricular em todos os componentes curriculares do primeiro semestre.

Parágrafo único: Nos demais semestres o estudante deverá cumprir, no mínimo 12 (doze) créditos, salvo a condição de concludente ou em casos especiais autorizados pela coordenadoria de curso ou, na ausência desta, do gestor máximo do ensino no campus.

\begin{center}
Seção II - Da renovação periódica da matrícula
\end{center}

Art. 79. A renovação de matrícula é um procedimento obrigatório pelo qual o estudante confirma seu interesse em manter o vínculo acadêmico com um curso do IFCE no período letivo seguinte.

§ 1º O período letivo pode se referir a um semestre letivo ou a um ano letivo, a depender da periodicidade de oferta de disciplinas do curso.

§ 2º A renovação da matrícula de um curso com periodicidade semestral deverá ser realizada a cada semestre, enquanto que para os cursos com periodicidade anual a renovação só precisará ser realizada uma vez a cada ano letivo.

Art. 80. A renovação de matrícula para os cursos técnicos e de graduação do IFCE deve ser solicitada pelo estudante de forma on-line no sistema acadêmico da instituição, de acordo com as datas previamente definidas em calendário acadêmico.

§ 1º O processo de renovação da matrícula deverá prever uma fase para solicitar a renovação e outra para ajustar a matrícula realizada pela CCA.

§ 2º O processo de renovação da matrícula deverá ser concluído até o final do período letivo que antecede o período letivo para o qual a renovação da matrícula está sendo pleiteada.

Art. 81. O estudante, que não solicitar a renovação on-line da matrícula no prazo estabelecido, deverá comparecer à CCA no prazo de 5 (cinco) dias letivos, a contar do último dia do prazo para a renovação de matrícula, a fim de regularizar sua situação acadêmica.

Parágrafo único: O estudante que não solicitar a renovação on-line da matrícula, nem comparecer fisicamente à CCA para regularizar sua situação acadêmica, deverá ser considerado desistente do curso, tendo sua situação de matrícula alterada para ABANDONO no sistema acadêmico.

\begin{center}
Subseção II - Da renovação nos cursos de regime de créditos por disciplina
\end{center}

Art. 85. O estudante de um curso com regime de crédito por disciplina, no momento que solicitar a renovação de matrícula, deverá indicar quais componentes curriculares deseja cursar.

Parágrafo único: Os componentes curriculares a serem cursados podem ser selecionados entre aqueles:
I. obrigatórios da matriz curricular do curso;
II. optativos da matriz curricular do curso;
III. que constam em matrizes curriculares de outros cursos técnicos subsequentes ou concomitantes, desde que haja equivalência entre os componentes e que não haja choque de horário entre eles.

Art. 86. O estudante, durante a fase de ajuste de matrícula, poderá incluir ou excluir componentes curriculares para o período letivo a ser cursado.

Art. 87. Após o período de ajuste de matrículas, não deverá ser mais permitido:
I. que o estudante inclua algum componente curricular;
II. que haja alteração de horário de disciplina.

Parágrafo único: Em casos excepcionais, a alteração acontecerá somente mediante autorização do gestor máximo do ensino no campus.

Art. 88. O processo de renovação de matrícula deverá ser por componente curricular, priorizando a seguinte ordem de ocupação de vagas:
I. componentes pendentes dos estudantes finalistas;
II. componentes curriculares do semestre regular;
III. desempenho acadêmico do estudante, expresso pelo Índice de Rendimento Acadêmico (IRA).

§ 1º Entende-se por estudantes finalistas aqueles que para concluir o curso, dependem somente das disciplinas pleiteadas na renovação da matrícula.

§ 2º O Índice de Rendimento Acadêmico (IRA) é um valor quantitativo utilizado para medir o desempenho acumulado pelo estudante, nos componentes curriculares, ao longo do desenvolvimento de um curso. O cálculo do IRA é feito através de uma média ponderada das notas de cada componente curricular, levando-se em consideração a quantidade de créditos destes na matriz curricular. Este cálculo é realizado a cada fechamento de período letivo e atualizado pelo sistema acadêmico do IFCE. Para fins de cálculo, utiliza-se a seguinte fórmula:

\begin{equation*}
IRA=\frac{(Mf_1 x Cr_1)+(Mf_2 x Cr_2)+...+(Mf_n x Cr_n)}{(Cr_1+Cr_2+...+Cr_n)}
\end{equation*}
\vspace{0,3cm}

Onde:

\begin{itemize}
\item MF = Média final do componente curricular
\item Cr = Créditos do componente curricular
\end{itemize}

§ 3º As notas de componentes curriculares associadas a um período letivo em curso não serão consideradas no cálculo do IRA.

§ 4º Para efeito de cálculo do IRA estarão incluídos todos os componentes curriculares cursados pelo estudante, com exceção de disciplinas com situação de trancamento, aproveitamento ou dispensa.

\subsection{Do Ingresso: Seção I, Seção II (Subseções I, II, III e IV), Seção III, IV e V}
\label{anexoRODtransferencia}

\begin{center}
Seção II - Do ingresso de diplomados e transferidos
\end{center}

Art. 49. O IFCE poderá receber, em todos os seus cursos, estudantes oriundos de instituições devidamente credenciadas pelos órgãos normativos dos sistemas de ensino municipal, estadual e federal.

§ 1º O IFCE não receberá estudantes oriundos de cursos sequenciais.

Art. 50. O edital para ingresso de diplomados e transferidos deverá prever a seguinte ordem de prioridade de atendimento:
I. ingressantes por transferência interna;
II. ingressantes por transferência externa;
III. ingressantes diplomados.

Art. 51. Para os que pleiteiam ingresso por transferência, deverá ser considerada a seguinte ordem de prioridade no preenchimento das vagas existentes:
I. o maior número de créditos obtidos nos componentes curriculares a serem aproveitados;
II. o maior índice de rendimento acadêmico (IRA) ou índice equivalente; e
III. a maior idade.

Art. 52. No âmbito do IFCE, o ingresso de estudantes dos cursos técnicos ou de graduação, por meio de transferência, pode ser dos seguintes tipos:
I. transferência Interna
II. transferência Externa

\begin{center}
Subseção I - Do ingresso por transferência interna
\end{center}

Art. 53. O ingresso por transferência interna é o processo de entrada de estudante em um curso de um campus do IFCE, quando este é oriundo de outro curso do mesmo campus.

Art. 54. A transferência interna só deverá ser admitida quando:
I. houver, preferencialmente, similaridade entre o curso de origem e o pleiteado no que concerne à área de conhecimento ou eixo tecnológico;
II. atender aos pré-requisitos de escolaridade e as especificidades do curso definidos em edital, mediante comprovação;
III. o curso de origem e o curso pleiteado forem do mesmo nível de ensino.

Parágrafo único – A transferência interna só poderá ser pleiteada uma vez.

\begin{center}
Subseção II - Do ingresso por transferência externa
\end{center}

Art. 55. O ingresso por transferência externa é o processo de entrada de estudante em um curso de um campus do IFCE, quando este é oriundo de outro campus do instituto ou de outra instituição de ensino.

Art. 56. Para ter direito à matrícula, o estudante que pleiteia o ingresso por transferência deverá:
I. comprovar que foi submetido a um processo seletivo similar ao do IFCE;
II. apresentar guia de transferência ou histórico escolar com status transferido;
III. obter aprovação em teste de aptidão específica, quando o curso pretendido o exigir.

\begin{center}
Subseção III - Do ingresso por transferência EX OFFICIO
\end{center}

Art. 57. A transferência ex officio é a forma de atendimento ao estudante egresso de outra instituição de ensino congênere, independentemente da existência de vaga, do período e de processo seletivo, por tratar-se de servidor público federal, civil ou militar, inclusive seus dependentes, e quando requerida em razão de comprovada remoção ou transferência de ofício, acarretando mudança de domicílio para o município onde se situe a instituição recebedora, ou para a localidade mais próxima desta.

§ 1º São benficiários dessa forma de ingresso o cônjuge e os dependentes do servidor até a idade de 24 anos, como caracterizado no caput deste artigo, desde que comprovado o amparo da Lei No. 9.536, de 11 de dezembro de 1997.

§ 2º Conforme estabelecido no parágrafo único da Lei No. 9.536/97, essa regra não se aplica quando o interessado na transferência se deslocar para assumir cargo efetivo em razão de concurso público, cargo comissionado ou função de confiança.

Art. 58. A solicitação de transferência ex officio deverá ser feita mediante requerimento protocolado no campus de destino e encaminhado ao gestor máximo do ensino no campus do IFCE, sendo necessários os seguintes documentos:
I. cópia do ato de transferência ex officio ou remoção, publicado no Diário Oficial da União (DOU), ou órgão oficial de divulgação ou publicação da própria corporação;
II. declaração original da autoridade maior do órgão competente, comprovando a remoção ou transferência ex officio.

\begin{center}
Subseção IV - Do ingresso de diplomados
\end{center}

Art. 59. Entende-se por diplomados aqueles que possuem diploma de cursos de educação profissional técnica de nível médio ou diploma de cursos de graduação.

Art. 60. O requerente deverá ser diplomado no nível respectivo ou superior ao pretendido.

Art. 61. O ingresso de diplomados deverá ser concedido mediante o atendimento em pelo menos um dos seguintes critérios abaixo relacionados, desde que estes estejam definidos em edital estabelecido pelo campus:
I. maior número de créditos a serem aproveitados no curso solicitado;
II. classificação em entrevista ou prova;
III. classificação em teste de habilidades específicas, quando o curso o exigir.
    
Art. 62. O requerimento para ingresso de diplomado deverá ser acompanhado dos seguintes documentos, em cópia autenticada ou com a apresentação original para conferência:
I. documento oficial de identidade com foto;
II. cadastro de pessoa física (CPF);
III. cópia autenticada de diploma ou certidão de conclusão;
IV. histórico escolar;
V. programa dos componentes curriculares cursados, autenticados pela instituição de origem;
VI. outros documentos especificados em edital.

\begin{center}
Seção III - Do ingresso por matrícula especial
\end{center}

Art. 63. Deverá ser admitida matrícula especial, ao estudante que deseje cursar componentes curriculares nos cursos técnicos e de graduação, desde que haja vaga nos componentes curriculares constantes na solicitação e que o requerente seja diplomado no nível respectivo ou superior ao pretendido.

Art. 64. O estudante com matrícula especial poderá cursar no máximo 3 (três) componentes curriculares, podendo posteriormente aproveitá-los, caso efetive uma matrícula no IFCE.

Parágrafo único: Candidatos que possuam diploma estrangeiro de curso técnico ou de graduação e se submeteram a processo de revalidação de diplomas no IFCE, poderão cursar mais de três disciplinas, na qualidade de estudante especial, desde que seja uma recomendação da comissão avaliadora da revalidação, registrada em parecer técnico.

Art. 65. A solicitação de matrícula especial deverá ser feita mediante requerimento protocolado e encaminhado à coordenadoria do curso, nos primeiros 50 (cinquenta) dias letivos do período letivo imediatamente anterior ao que deverá ser cursado, devendo ser acompanhada dos seguintes documentos:
I. cópia do diploma para quem deseja matrícula na graduação, devidamente autenticada ou acompanhada do original;
II. cópia do diploma de conclusão do curso técnico de nível médio para quem deseja matrícula em curso técnico, devidamente autenticada ou acompanhada do original;
III. cópia do histórico escolar autenticada ou acompanhada do original.

§ 1º A coordenadoria do curso pleiteado pelo interessado deverá emitir o parecer no prazo de 30 (trinta) dias.

§ 2º Caberá à Proen encaminhar o parecer técnico ao gestor máximo do ensino no campus que, por conseguinte, deverá tomar as providências de efetivação de matrícula especial desses candidatos junto à sua CCA.

Art. 66. A matrícula especial não assegura, em qualquer hipótese, vínculo como estudante regular do IFCE.

Art. 67. O estudante com matrícula especial ficará sujeito às normas disciplinares e didático- pedagógicas, inclusive submetendo-se ao sistema de avaliação do componente curricular.

Art. 68. O estudante aprovado terá direito à declaração emitida pela CCA, constando: o componente curricular cursado, a carga horária, o período, a nota, a frequência e a ementa.

Art. 69. Em nenhuma hipótese, deverá ser permitido o ingresso informal de estudante ouvinte nos cursos do IFCE, sendo, portanto, o ingresso concedido somente ao aluno com matrícula especial, mediante documentação apresentada e parecer autorizativo.

\begin{center}
Seção IV - Do Reingresso
\end{center}

Art. 70. O IFCE concederá, em oportunidade única, o direito de reingresso a estudantes que abandonaram o curso, nas seguintes condições:
I. terem decorridos, no máximo, 5 (cinco) anos, a contar da data em que o estudante deixou de frequentar o curso;
II. existir vaga no curso;
III. apresentar em requerimento a quitação com a biblioteca (nada consta).

Art. 71. A solicitação de reingresso deverá ser feita mediante requerimento protocolado e enviado à coordenação de curso para análise e emissão de parecer.

§ 1º Em caso de deferimento da solicitação, o coordenador do curso deverá comunicar à CCA para que o estudante seja matriculado no sistema acadêmico.

§ 2º O estudante deverá receber um novo código de matrícula e ser vinculado à matriz curricular vigente do curso no qual está reingressando.

§ 3º A forma de ingresso do estudante a ser registrada no sistema acadêmico deverá ser REINGRESSO;

§ 4º Para aproveitar os componentes curriculares cursados com a matrícula anterior, o estudante deverá solicitar o aproveitamento de componentes curriculares, de acordo com os procedimentos estabelecidos no Capítulo IV - Seção I.

Art. 72. Não deverá ser permitido o reingresso de estudantes que deixaram de frequentar o curso:
I. no primeiro semestre – para cursos com periodicidade de oferta semestral de vagas;
II. no primeiro ano – para cursos com periodicidade de oferta anual de vagas.

\begin{center}
Seção V - Da ocupação de duas vagas em cursos do mesmo nível
\end{center}   

Art. 73. No âmbito do IFCE, em nenhuma hipótese deverá ser permitida aos estudantes de cursos de graduação, a ocupação de vagas em mais de um curso do mesmo nível de ensino.

Art. 74. Ao constatar que há estudante ocupando mais de uma vaga em cursos de graduação no IFCE, ou no IFCE e em outra instituição pública, a CCA deverá comunicar ao estudante a possibilidade de optar por uma das vagas no prazo de 5 (cinco) dias úteis, contado do primeiro dia útil posterior à comunicação.

§ 1º Caso o estudante não compareça no prazo assinalado neste artigo ou não opte por uma das vagas, a instituição providenciará o cancelamento:
I. da matrícula mais antiga, na hipótese da duplicidade ocorrer em instituições diferentes;
II. da matrícula mais recente, na hipótese da duplicidade ocorrer na mesma instituição.

§ 2º Concomitantemente ao cancelamento compulsório da matrícula na forma do disposto no § 1º deste artigo, deverá ser decretada a nulidade dos créditos adquiridos no curso cuja matrícula foi cancelada.

\subsection{Da Sistemática de Avaliação: Subseção I da Seção I, Capítulo III, Título III}
\label{anexoRODavaliacao}

\begin{center}
Seção I - Da sistemática de avaliação
\end{center}

Art. 94. Os processos, instrumentos, critérios e valores de avaliação adotados pelo professor deverão ser explicitados aos estudantes no início do período letivo, quando da apresentação do PUD, observadas as normas dispostas neste documento.

§ 1º As avaliações devem ter caráter diagnóstico, formativo, contínuo e processual, podendo constar de:
I. observação diária dos estudantes pelos professores, durante a aplicação de suas diversas atividades;
II. exercícios;
III. trabalhos individuais e/ou coletivos;
IV. fichas de observações;
V. relatórios;
VI. autoavaliação;
VII. provas escritas com ou sem consulta;
VIII. provas práticas e provas orais;
IX. seminários;
X. projetos interdisciplinares;
XI. resolução de exercícios;
XII. planejamento e execução de experimentos ou projetos;
XIII. relatórios referentes a trabalhos, experimentos ou visitas técnicas,
XIV. realização de eventos ou atividades abertas à comunidade;
XV. autoavaliação descritiva e outros instrumentos de avaliação considerando o seu caráter progressivo.

Art. 95. Ao estudante deverá ser assegurado o direito de conhecer os resultados das avaliações mediante vistas dos referidos instrumentos, apresentados pelos professores como parte do processo de ensino e aprendizagem.

§ 1º As avaliações escritas deverão ser devolvidas; e as demais, informadas ao estudante e registradas no sistema acadêmico, logo após a devida correção em um prazo máximo de até 10 (dez) dias letivos.

§ 2º A divulgação de resultados tem caráter individual, sendo vedada a sua exposição pública, salvo em casos de haver consentimento prévio do estudante.

Art. 96. O estudante que discordar do resultado obtido em qualquer avaliação da aprendizagem poderá requerer, à coordenadoria de curso, revisão no prazo de 2 (dois) dias letivos após a comunicação do resultado.

§ 1º A revisão da avaliação deverá ser feita pelo docente do componente curricular, juntamente com o coordenador do curso.

§ 2º Caso a revisão não possa ser feita pelo professor do componente curricular, o coordenador deverá designar outro docente para tal ação.

\begin{center}
Subseção I - Avaliação nos cursos com regime de créditos por disciplina
\end{center}

Art. 97. A sistemática de avaliação dos conhecimentos construídos, nos cursos com regime de crédito por disciplina, com periodicidade semestral, se desenvolverá em duas etapas.

§ 1º Deverá ser registrada no sistema acadêmico apenas uma nota para a primeira etapa (N1) e uma nota para a segunda etapa (N2), com pesos 2 e 3, respectivamente.

§ 2º O docente deverá aplicar, no mínimo, duas avaliações em cada uma das etapas.

§ 3º O critério para composição da nota de cada etapa, a partir das notas obtidas em cada uma das avaliações, ficará a cargo do docente da disciplina, em consonância com o estabelecido no PUD.

Art. 98. O cálculo da média parcial (MP) de cada disciplina deve ser feito de acordo com a seguinte equação:

\begin{equation*}
X_S=\frac{2X_1+3X_2}{5}\geq 7,0
\end{equation*}
\vspace{0,3cm}

Art. 99. Deverá ser considerado aprovado no componente curricular o estudante que, ao final do período letivo, tenha frequência igual ou superior a 75\% (setenta e cinco por cento) do total de horas letivas e tenha obtido média parcial (MP) igual ou superior a:
I. 6,0 (seis), para disciplinas de cursos técnicos concomitantes e subsequentes.
II. 7,0 (sete), para disciplinas de cursos de graduação.

Parágrafo único: Os estudantes aprovados com a nota da MP não precisarão realizar a avaliação final (AF) e sua média final (MF) deverá ser igual a sua média parcial (MP).

Art. 100. Deverão fazer avaliação final (AF) o estudante de curso técnico que obtiver MP inferior a 6,0 (seis) e maior ou igual a 3,0 (três), e o estudante de graduação que obtiver MP inferior a 7,0 (sete) e maior ou igual a 3,0 (três).

§ 1º A avaliação final deverá ser aplicada no mínimo 3 (três) dias letivos após o registro do resultado da MP no sistema acadêmico.

§ 2º A avaliação final poderá contemplar todo o conteúdo trabalhado no período letivo.

§ 3º A nota da avaliação final (AF) deverá ser registrada no sistema acadêmico.

§ 4º O cálculo da média final (MF) o estudante referido no caput deverá ser efetuado de acordo com a seguinte equação:

\begin{equation*}
MF=\frac{X_S+AF}{2}\geq 5,0
\end{equation*}
\vspace{0,3cm}

§ 5º Deverá ser considerado aprovado na disciplina o estudante que, após a realização da avaliação final, obtiver média final (MF) igual ou maior que 5,0 (cinco).

%\clearpage\section{MODELO DE FORMULÁRIO DA AVALIAÇÃO DOCENTE}
%\label{anexoAvaliacao}
%\includepdf[pages={1}]{anexos/FormularioAvaliacaoProfessor.pdf}

\clearpage\section{PROJETO DO BLOCO II DO EIXO TECNOLÓGICO DA INDÚSTRIA}
\label{anexoProjeto}
\includepdf[pages={1-3}]{anexos/ProjetoOficina.pdf}

%%%%%%%%%%%
\clearpage\section{LEGISLAÇÃO}
\label{anexoLegislacao}

\subsection{Decreto Federal n° 5.626, de 22 de dezembro de 2005}
\label{anexoDecreto}
\includepdf[pages={1-8}]{anexos/Decreto5626-22122005.pdf}


\clearpage\section{PROGRAMAS DE UNIDADE DIDÁTICA}
\label{anexoPUD}

%\pagestyle{empty}

\newpage
\if\COMPILAR0
    \input{PPC-ECAPUDs}
\else
    \input{PPC-EMEPUDs}
\fi
\newpage

\end{document}
